% Documento de Ingeniería del Proyecto – Fase 1 · Equipo B
% Compilar desde docs/documento/:
%   pdflatex -interaction=nonstopmode IS1-Fase1-EquipoB.tex   (dos veces, por el índice)
% Cada sección vive en secciones/ para que cada rol la edite en su rama sin conflictos.
% Entregables (formato del profesor): docs/entregas/IS1-Fase1-EquipoB-v<versión>.pdf y .docx
\newcommand{\versiondoc}{0.2}
\newcommand{\fechadoc}{11/11/2026}
% Preámbulo del Documento de Ingeniería del Proyecto – Fase 1 (Equipo B).
% Formato del machote 0.3: carta, márgenes de 2.5 cm, 11 pt, páginas numeradas.
\PassOptionsToPackage{unicode}{hyperref}
\PassOptionsToPackage{hyphens}{url}
\documentclass[11pt,letterpaper,spanish]{article}
\usepackage[margin=2.5cm]{geometry}
\usepackage[T1]{fontenc}
\usepackage[utf8]{inputenc}
\usepackage{textcomp}
\usepackage{lmodern}
\usepackage{xcolor}
\usepackage{amsmath,amssymb}
\usepackage{microtype}
\usepackage{parskip}
\usepackage{longtable,booktabs,array}
\usepackage{calc}
\newcounter{none} % tablas sin número (formato heredado de la conversión)
\newcolumntype{L}[1]{>{\raggedright\arraybackslash}p{#1}} % columna de ancho fijo alineada a la izquierda
\usepackage{etoolbox}
\makeatletter
\patchcmd\longtable{\par}{\if@noskipsec\mbox{}\fi\par}{}{}
\makeatother
\usepackage{footnotehyper}
\makesavenoteenv{longtable}
\usepackage{graphicx}
\usepackage{float}
\usepackage[labelformat=empty,font=small]{caption}
\usepackage{soul}
\usepackage[bidi=default,shorthands=off]{babel}
\usepackage{newunicodechar}
\newunicodechar{☐}{$\square$}
\newunicodechar{✓}{\checkmark}
\newunicodechar{≤}{$\leq$}
\newunicodechar{≥}{$\geq$}
\newunicodechar{→}{$\rightarrow$}
\usepackage{fancyhdr}
\pagestyle{fancy}
\fancyhf{}
\fancyhead[L]{\small IS1 · Equipo B · Sistema de Órdenes de Servicio y Citas}
\fancyhead[R]{\small v\versiondoc}
\fancyfoot[C]{\thepage}
\setlength{\emergencystretch}{3em}
\providecommand{\tightlist}{\setlength{\itemsep}{0pt}\setlength{\parskip}{0pt}}
\setcounter{secnumdepth}{-\maxdimen} % la numeración va en el título, como en el machote
\setcounter{tocdepth}{2}
\usepackage{bookmark}
\usepackage{xurl}
\urlstyle{same}
\hypersetup{pdflang={es}, hidelinks, pdftitle={IS1-Fase1-EquipoB-v\versiondoc}}

% ---------------------------------------------------------------------------
% Figuras. El machote pide número, título y ruta del archivo fuente.
%   \figura{archivo}{opciones de tamaño}{número}{título}{ruta fuente}
\DeclareUrlCommand\ruta{\urlstyle{tt}} % rutas del repositorio, con cortes de línea
\captionsetup{justification=raggedright,singlelinecheck=true}
\newcommand{\figura}[5]{%
  \begin{figure}[H]\centering
  \includegraphics[#2,keepaspectratio]{#1}
  \caption[Figura #3. #4]{\textbf{Figura #3.} #4. Fuente: \ruta{#5}}
  \end{figure}}

% Figura descrita: mientras no exista la imagen se muestra un recuadro neutro
% con su contenido y la ruta de la fuente. Se reemplaza por \figura cuando exista.
%   \figurapendiente{número}{título}{contenido de la figura}{ruta fuente}
\newcommand{\figurapendiente}[4]{%
  \begin{figure}[H]\centering
  \fcolorbox{black!40}{black!3}{\parbox{0.92\linewidth}{\small
    \textbf{Contenido:} #3\par\smallskip
    \textbf{Archivo fuente:} \ruta{#4}}}
  \caption[Figura #1. #2]{\textbf{Figura #1.} #2. Fuente: \ruta{#4}}
  \end{figure}}

% Dato por definir (sin corchetes ni color, para la versión de entrega).
\newcommand{\pendiente}[1]{por definir: #1}


\begin{document}

\begin{titlepage}
\centering
{\large Universidad del Istmo · Campus Ixtepec\par}
Ingeniería en Desarrollo de Software y Sistemas Inteligentes\par
Ingeniería de Software I (IDSSI-27) · Quinto semestre\par
\vspace{3cm}
{\LARGE\bfseries Documento de Ingeniería del Proyecto -- Fase 1\par}
\vspace{1cm}
{\Large Sistema de Órdenes de Servicio y Citas\par}
\vspace{0.5cm}
Servicio Automotriz Hermanos Ruiz, S.A. de C.V.\par
\vfill
\begin{tabular}{rl}
\textbf{Asignatura:} & Ingeniería de Software I (IDSSI-27)\\
\textbf{Documento:} & Documento de Ingeniería del Proyecto -- Fase 1\\
\textbf{Equipo:} & B (5 integrantes)\\
\textbf{Cliente:} & Servicio Automotriz Hermanos Ruiz\\
\textbf{Producto:} & Sistema de Órdenes de Servicio y Citas\\
\textbf{Versión del documento:} & v\versiondoc\\
\textbf{Fecha:} & \fechadoc\\
\textbf{Repositorio:} & \url{https://github.com/AbdiasChavez/Florentino}\\
\end{tabular}
\end{titlepage}

\section*{Control de versiones}

La versión sigue la tabla de entrega acumulativa del machote (0.2): v0.1 en la Review 0, v0.2 en la
Review 1, v1.0 con el SRS firmado en la Review 2, y así sucesivamente.

{\small
\begin{longtable}[]{@{}
  >{\raggedright\arraybackslash}p{0.07\linewidth}
  >{\raggedright\arraybackslash}p{0.12\linewidth}
  >{\raggedright\arraybackslash}p{0.12\linewidth}
  >{\raggedright\arraybackslash}p{0.43\linewidth}
  >{\raggedright\arraybackslash}p{0.16\linewidth}@{}}
\toprule
\textbf{Versión} & \textbf{Fecha} & \textbf{Autor(es)} & \textbf{Descripción del cambio} & \textbf{Revisó} \\
\midrule
v0.1 & 29/10/2026 & Equipo B & Primera entrega del documento, revisada en la Review 1 (incluía avance de secciones de cortes posteriores). & \pendiente{} \\
v0.2 & \fechadoc & Equipo B & Atención a las observaciones de la Review 1. & \pendiente{} \\
\bottomrule
\end{longtable}
}

\section*{Aprobación del cliente}

Se llena en la Review 2 (firma del SRS) y en la Demo final (aceptación del Release v0.1.0).

{\small
\begin{longtable}[]{@{}
  >{\raggedright\arraybackslash}p{0.2\linewidth}
  >{\raggedright\arraybackslash}p{0.2\linewidth}
  >{\raggedright\arraybackslash}p{0.22\linewidth}
  >{\raggedright\arraybackslash}p{0.16\linewidth}
  >{\raggedright\arraybackslash}p{0.1\linewidth}@{}}
\toprule
\textbf{Hito} & \textbf{Resultado (aceptado / aceptado con observaciones / rechazado)} & \textbf{Observaciones} & \textbf{Firma del cliente} & \textbf{Fecha} \\
\midrule
SRS v1.0 (Review 2) & & & & \\
Release v0.1.0 (Demo final) & & & & \\
\bottomrule
\end{longtable}
}

\clearpage
\tableofcontents
\clearpage

\section{1. Información general del
proyecto}\label{informaciuxf3n-general-del-proyecto}

\subsection{1.1 Integrantes y roles}\label{integrantes-y-roles}

{\def\LTcaptype{none} % do not increment counter
\begin{longtable}[]{@{}
  >{\raggedright\arraybackslash}p{(\linewidth - 8\tabcolsep) * \real{0.2456}}
  >{\raggedright\arraybackslash}p{(\linewidth - 8\tabcolsep) * \real{0.1282}}
  >{\raggedright\arraybackslash}p{(\linewidth - 8\tabcolsep) * \real{0.2012}}
  >{\raggedright\arraybackslash}p{(\linewidth - 8\tabcolsep) * \real{0.1686}}
  >{\raggedright\arraybackslash}p{(\linewidth - 8\tabcolsep) * \real{0.2564}}@{}}
\toprule\noalign{}
\begin{minipage}[b]{\linewidth}\raggedright
\textbf{Nombre completo}
\end{minipage} & \begin{minipage}[b]{\linewidth}\raggedright
\textbf{Matrícula}
\end{minipage} & \begin{minipage}[b]{\linewidth}\raggedright
\textbf{Usuario GitHub}
\end{minipage} & \begin{minipage}[b]{\linewidth}\raggedright
\textbf{Rol(es) y sprints}
\end{minipage} & \begin{minipage}[b]{\linewidth}\raggedright
\textbf{Áreas principales de contribución}
\end{minipage} \\
\midrule\noalign{}
\endhead
\bottomrule\noalign{}
\endlastfoot
Ángel Yael Avendaño Guzmán & 0224170023 & AyagDan & Product Owner
(S0--S6) & Enlace con el cliente, levantamiento de requerimientos en
entrevistas, priorización del Product Backlog y aceptación de
entregables. \\
Abdías Daniel Chávez Cruz & 0224170024 & absdan & Scrum Master /
Tech Lead (S0--S6) & Facilitación de Scrum, arquitectura del backend en
FastAPI, diseño de la base de datos relacional y gestión técnica de Pull
Requests. \\
Dulce María Estudillo Rosado & 0224170028 & DulceRosado & QA Lead /
Developer (S0--S6) & Diseño del plan de pruebas, ejecución de casos de
aceptación, validación de calidad y realización del QA cruzado al otro
equipo. \\
Isis Alhelí González Maldonado & 0224170030 & \ruta{IsisAlheliGonzalezMaldonado} &
Developer (S0--S6) & Redacción de la especificación de requisitos de
software (SRS), historias de usuario y modelado de procesos del
negocio. \\
Isis Yamilet Meléndez Hernández & 0224170037 & \ruta{IsisMelendezHernandez} &
Developer (S0--S6): especialista en interfaz, datos y documentación &
Integración y compilación del Documento de Ingeniería (migración a
LaTeX/Word), configuración de flujos de CI en GitHub Actions, despliegue
y apoyo en el desarrollo de los componentes de la interfaz (PWA). \\
\end{longtable}
}

\subsection{1.2 Enlaces del proyecto}\label{enlaces-del-proyecto}

{\def\LTcaptype{none} % do not increment counter
\begin{longtable}[]{@{}
  >{\raggedright\arraybackslash}p{(\linewidth - 2\tabcolsep) * \real{0.3631}}
  >{\raggedright\arraybackslash}p{(\linewidth - 2\tabcolsep) * \real{0.6369}}@{}}
\toprule\noalign{}
\begin{minipage}[b]{\linewidth}\raggedright
\textbf{Recurso}
\end{minipage} & \begin{minipage}[b]{\linewidth}\raggedright
\textbf{Enlace}
\end{minipage} \\
\midrule\noalign{}
\endhead
\bottomrule\noalign{}
\endlastfoot
Repositorio & \url{https://github.com/AyagDan/SistemaDeServiciosYCitas} \\
Tablero (GitHub Projects / Jira) & Historias como issues del repositorio:
\url{https://github.com/AyagDan/SistemaDeServiciosYCitas/issues} \\
Prototipo de UI (Figma / Penpot) & Se agrega con los prototipos de baja
fidelidad (4.10, Review 2). \\
Ambiente de pruebas (staging) & Desde el Sprint 4. \\
Release v0.1.0 & Se publica en GitHub Releases en la Demo final
(Sprint 6). \\
Canal de comunicación del equipo & Grupo de WhatsApp del equipo
(mensajes y llamadas); reuniones presenciales en la escuela. \\
\end{longtable}
}

\subsection{1.3 Resumen ejecutivo}\label{resumen-ejecutivo}

Servicio Automotriz Hermanos Ruiz opera hoy con cuadernos, órdenes en
papel y llamadas. Esto provoca disputas por cobros que el cliente dice
no haber autorizado, reclamos por daños sin evidencia del estado en que
llegó el vehículo, muchas llamadas para preguntar por el avance y poca
visibilidad de qué vehículo ocupa cada una de las 4 bahías.

El Sistema de Órdenes de Servicio y Citas digitalizará el ciclo de
atención: la cita, la recepción del vehículo con checklist y
fotografías, el diagnóstico y la asignación de bahía, la cotización que
el cliente autoriza desde su celular con un código de verificación, y el
cobro y la entrega. El cliente consultará el estado de su vehículo con un
enlace seguro, sin instalar aplicaciones.

Estado actual (v\versiondoc): además de la visión y la viabilidad
(sección 3), ya está en borrador la especificación de requisitos
(sección 4): el glosario con los términos del taller, los usuarios, los
requisitos funcionales, las 30 historias de usuario con su prioridad, los
requisitos de calidad medibles y las reglas de negocio, incluidos los
procedimientos para cuando falle la luz o internet. En esta Review se
recorre con el cliente para validarla; la versión firmada se entrega en
la Review 2. La Fase 2 agregará reportes gerenciales, consulta de
existencias y de garantías, y el tablero simplificado de bahías.

\section{2. Gestión del proyecto con
Scrum}\label{gestiuxf3n-del-proyecto-con-scrum}

\subsection{2.1 Acuerdos de trabajo}\label{acuerdos-de-trabajo}

Los acuerdos completos están en CONTRIBUTING.md del repositorio. En
resumen:

\textbf{• Comunicación:} grupo de WhatsApp del equipo para lo
urgente y GitHub (issues y Pull Requests) para todo lo técnico. Tiempo
máximo de respuesta: 24 horas en días hábiles.

\textbf{• Reuniones:} se hacen de manera presencial o por llamada en el
grupo de WhatsApp: Daily breve de lunes a viernes y los eventos de Scrum
al inicio y al cierre de cada sprint de dos semanas. Lo que se decide en
una reunión se registra en GitHub (issue, minuta o documento).

\textbf{• Decisiones:} por consenso; si no se logra, decide el Product
Owner en temas de alcance y prioridad, y el Scrum Master en temas de
proceso.

\textbf{• Ramas:} seguimos GitHub Flow. main está protegida y cada
historia o tarea se trabaja en su propia rama, creada desde main, con la
convención \texttt{tipo/número-descripción}, donde el tipo es
\texttt{feat}, \texttt{fix}, \texttt{docs}, \texttt{test} o \texttt{ci} y
el número es el issue que atiende (por ejemplo,
\texttt{fix/9-version-stakeholders} para la observación B11). Cada rama
se integra a main mediante un Pull Request que cierra su issue y tiene al
menos una revisión de otro integrante; después se borra. Las ramas por
integrante (\texttt{rama/<NOMBRE>}) se usaron hasta la v0.2 y ya no
reciben trabajo nuevo.

\textbf{• Commits:} en español y en presente, con el formato
«Rol: mensaje» y referencia a la historia cuando aplique (por ejemplo,
«Developers: implementa recepción con checklist (HU-01)»).

\textbf{• Documentos:} el Documento de Ingeniería se escribe en LaTeX,
un archivo por sección en \texttt{docs/documento/secciones/}, para que
git pueda fusionar los cambios de varios roles. El entregable se publica en
\texttt{docs/entregas/} en PDF y en Word (.docx) con el nombre
\texttt{IS1-Fase1-EquipoB-v<versión>} (por ejemplo, IS1-Fase1-EquipoB-v0.1.pdf).

\textbf{• Incumplimiento:} se registra en la retrospectiva; a la segunda
falta en el mismo sprint, el Scrum Master lo informa al profesor.

\subsection{2.2 Definición de Listo (DoR) y Definición de Terminado
(DoD)}\label{definiciuxf3n-de-listo-dor-y-definiciuxf3n-de-terminado-dod}

{\def\LTcaptype{none} % do not increment counter
\begin{longtable}[]{@{}
  >{\raggedright\arraybackslash}p{(\linewidth - 2\tabcolsep) * \real{0.5001}}
  >{\raggedright\arraybackslash}p{(\linewidth - 2\tabcolsep) * \real{0.4999}}@{}}
\toprule\noalign{}
\begin{minipage}[b]{\linewidth}\raggedright
\textbf{Definición de Listo}
\end{minipage} & \begin{minipage}[b]{\linewidth}\raggedright
\textbf{Definición de Terminado}
\end{minipage} \\
\midrule\noalign{}
\endhead
\bottomrule\noalign{}
\endlastfoot
• La historia sigue el formato Como\ldots{} quiero\ldots{} para\ldots{}
y tiene ID e issue en el tablero.

• Tiene al menos 2 criterios de aceptación en formato
Dado/Cuando/Entonces.

• El equipo la estimó en puntos y cabe en un sprint (8 puntos o menos).

• Sus reglas de negocio (RN) y requisitos (RF/RNF) relacionados están
identificados.

• Sus dependencias (otras historias, datos, wireframe) están resueltas o
planificadas.

• El Product Owner la priorizó (MoSCoW). & • El código está en main
mediante un Pull Request revisado por otro integrante.

• Cumple todos sus criterios de aceptación y el Product Owner lo
verificó.

• Tiene pruebas automatizadas y la integración continua pasa sin
errores.

• Se probó solo con datos ficticios; no hay datos personales reales ni
contraseñas en el repositorio.

• La documentación afectada (SRS, modelos, README) está actualizada.

• Está desplegada en el ambiente de pruebas (desde el Sprint 4). \\
\end{longtable}
}

\subsection{2.3 Resumen de sprints}\label{resumen-de-sprints}

{\def\LTcaptype{none} % do not increment counter
\begin{longtable}[]{@{}
  >{\raggedright\arraybackslash}p{(\linewidth - 10\tabcolsep) * \real{0.0856}}
  >{\raggedright\arraybackslash}p{(\linewidth - 10\tabcolsep) * \real{0.1495}}
  >{\raggedright\arraybackslash}p{(\linewidth - 10\tabcolsep) * \real{0.3270}}
  >{\raggedright\arraybackslash}p{(\linewidth - 10\tabcolsep) * \real{0.1601}}
  >{\raggedright\arraybackslash}p{(\linewidth - 10\tabcolsep) * \real{0.1389}}
  >{\raggedright\arraybackslash}p{(\linewidth - 10\tabcolsep) * \real{0.1389}}@{}}
\toprule\noalign{}
\begin{minipage}[b]{\linewidth}\raggedright
\textbf{Sprint}
\end{minipage} & \begin{minipage}[b]{\linewidth}\raggedright
\textbf{Fechas}
\end{minipage} & \begin{minipage}[b]{\linewidth}\raggedright
\textbf{Objetivo del sprint}
\end{minipage} & \begin{minipage}[b]{\linewidth}\raggedright
\textbf{¿Se cumplió? (sí / parcial / no)}
\end{minipage} & \begin{minipage}[b]{\linewidth}\raggedright
\textbf{Puntos comprometidos}
\end{minipage} & \begin{minipage}[b]{\linewidth}\raggedright
\textbf{Puntos terminados}
\end{minipage} \\
\midrule\noalign{}
\endhead
\bottomrule\noalign{}
\endlastfoot
S0 & \pendiente{} & Conocer al cliente, levantar los requisitos
iniciales y preparar el repositorio y el documento base. &
Sí & --- & --- \\
S1 & \pendiente{} & SRS en borrador (4.1--4.7) para validarlo con el
cliente. & Sí: el borrador se valida con el cliente en la Review 1 &
0 & 0 \\
\end{longtable}
}

Los sprints 0 y 1 fueron de especificación: su objetivo no incluyó
historias de usuario, así que no hay puntos comprometidos ni terminados
y las 30 historias siguen en estado «Pendiente» (4.5). La velocidad se
calcula a partir del Sprint 2, con los puntos de las historias
terminadas según la DoD.

\subsubsection{Evidencia de los eventos}\label{evidencia-de-los-eventos}

{\def\LTcaptype{none} % do not increment counter
\begin{longtable}[]{@{}
  >{\raggedright\arraybackslash}p{(\linewidth - 8\tabcolsep) * \real{0.0962}}
  >{\raggedright\arraybackslash}p{(\linewidth - 8\tabcolsep) * \real{0.2136}}
  >{\raggedright\arraybackslash}p{(\linewidth - 8\tabcolsep) * \real{0.2137}}
  >{\raggedright\arraybackslash}p{(\linewidth - 8\tabcolsep) * \real{0.3162}}
  >{\raggedright\arraybackslash}p{(\linewidth - 8\tabcolsep) * \real{0.1604}}@{}}
\toprule\noalign{}
\begin{minipage}[b]{\linewidth}\raggedright
\textbf{Sprint}
\end{minipage} & \begin{minipage}[b]{\linewidth}\raggedright
\textbf{Minuta de Review}
\end{minipage} & \begin{minipage}[b]{\linewidth}\raggedright
\textbf{Retrospectiva}
\end{minipage} & \begin{minipage}[b]{\linewidth}\raggedright
\textbf{Acciones de mejora acordadas}
\end{minipage} & \begin{minipage}[b]{\linewidth}\raggedright
\textbf{¿Se aplicaron?}
\end{minipage} \\
\midrule\noalign{}
\endhead
\bottomrule\noalign{}
\endlastfoot
S0 & \pendiente{} & \pendiente{} & --- & --- \\
S1 & \ruta{docs/minutas/2026-11-10-review-1.md} & \pendiente{} & --- & --- \\
\end{longtable}
}

\subsection{2.4 Registro de riesgos}\label{registro-de-riesgos}

{\def\LTcaptype{none} % do not increment counter
\begin{longtable}[]{@{}
  >{\raggedright\arraybackslash}p{(\linewidth - 10\tabcolsep) * \real{0.0856}}
  >{\raggedright\arraybackslash}p{(\linewidth - 10\tabcolsep) * \real{0.2990}}
  >{\raggedright\arraybackslash}p{(\linewidth - 10\tabcolsep) * \real{0.1282}}
  >{\raggedright\arraybackslash}p{(\linewidth - 10\tabcolsep) * \real{0.1069}}
  >{\raggedright\arraybackslash}p{(\linewidth - 10\tabcolsep) * \real{0.2735}}
  >{\raggedright\arraybackslash}p{(\linewidth - 10\tabcolsep) * \real{0.1067}}@{}}
\toprule\noalign{}
\begin{minipage}[b]{\linewidth}\raggedright
\textbf{ID}
\end{minipage} & \begin{minipage}[b]{\linewidth}\raggedright
\textbf{Riesgo}
\end{minipage} & \begin{minipage}[b]{\linewidth}\raggedright
\textbf{Probabilidad}
\end{minipage} & \begin{minipage}[b]{\linewidth}\raggedright
\textbf{Impacto}
\end{minipage} & \begin{minipage}[b]{\linewidth}\raggedright
\textbf{Respuesta y responsable}
\end{minipage} & \begin{minipage}[b]{\linewidth}\raggedright
\textbf{Estado}
\end{minipage} \\
\midrule\noalign{}
\endhead
\bottomrule\noalign{}
\endlastfoot
R-01 & Un integrante abandona el curso a mitad del semestre & Media &
Alto & Mitigar: conocimiento compartido mediante revisión cruzada de PR;
responsable: SM & Abierto \\
R-02 & La operación sin conexión y la sincronización posterior (RNF-08,
RNF-09) resultan más complejas de lo estimado & Alta & Alto & Mitigar:
prototipo técnico en el Sprint 1 y limitar lo offline a recepción y
diagnóstico; responsable: equipo de desarrollo & Abierto \\
R-03 & El cliente no está disponible para validar requisitos o asistir a
las Reviews & Media & Alto & Mitigar: agendar sesiones con anticipación
y validar por correo cuando no pueda asistir; responsable: PO &
Abierto \\
R-04 & Cambios de alcance frecuentes por parte del cliente & Media &
Medio & Mitigar: registrar cada solicitud de cambio (2.5) con análisis
de impacto; responsable: PO & Abierto \\
R-05 & El servicio gratuito de correo limita el envío de enlaces y
códigos 2FA & Media & Alto & Mitigar: probar los límites del plan
gratuito en el Sprint 1 y dejar SMS como opción (RF-02.3); responsable:
equipo de desarrollo & Abierto \\
R-06 & Las fotografías de recepción saturan el almacenamiento del
hospedaje gratuito & Media & Medio & Mitigar: comprimir las imágenes en
el dispositivo y fijar un tamaño máximo por foto; responsable: equipo de
desarrollo & Abierto \\
R-07 & Poca experiencia del equipo con FastAPI y PWA & Alta & Medio &
Mitigar: capacitación en S0--S1 y trabajo en pares en las primeras
historias; responsable: SM & Abierto \\
R-08 & Conflictos al editar en paralelo el documento .docx en ramas
distintas & Alta & Medio & Evitar: el documento se migró a LaTeX con un
archivo por sección; responsable: SM & Cerrado \\
R-09 & Se publican datos personales reales en el repositorio público &
Media & Alto & Evitar: solo datos ficticios y revisión de datos
personales en cada PR; .gitignore excluye .env y credenciales;
responsable: todos & Abierto \\
R-10 & Código o texto generado con IA que el equipo no entiende o no
revisó & Media & Alto & Mitigar: declarar el uso de IA en el documento, revisión de
cada PR por otro integrante y pruebas automatizadas como evidencia;
responsable: SM & Abierto \\
R-11 & Diferencias entre SQLite (pruebas) y PostgreSQL (staging) ocultan
defectos & Media & Medio & Mitigar: cargar el esquema en PostgreSQL
en la integración continua desde que exista el backend (S2);
responsable: QA Lead & Abierto \\
R-12 & Sin ambiente de pruebas no se puede mostrar el flujo núcleo al
cliente & Media & Alto & Mitigar: decidir el hospedaje del ambiente de pruebas antes del Sprint 4;
responsable: Developers & Abierto \\
\end{longtable}
}

\subsection{2.6 Declaración de uso de inteligencia
artificial}\label{declaracion-de-uso-de-ia}

El equipo usó un asistente de inteligencia artificial para pasar el
documento a LaTeX, atender las observaciones de la Review 1 y organizar el
repositorio. La tabla indica en qué se usó, cómo se verificó el resultado
y quién responde por él.

{\small
\begin{longtable}[]{@{}
  >{\raggedright\arraybackslash}p{0.13\linewidth}
  >{\raggedright\arraybackslash}p{0.21\linewidth}
  >{\raggedright\arraybackslash}p{0.20\linewidth}
  >{\raggedright\arraybackslash}p{0.22\linewidth}
  >{\raggedright\arraybackslash}p{0.13\linewidth}@{}}
\toprule
\textbf{Herramienta} & \textbf{Sección o artefacto} & \textbf{Para qué se usó} & \textbf{Cómo se verificó} & \textbf{Responsable} \\
\midrule
\endhead
\bottomrule
\endlastfoot
Claude Code (Anthropic) & Documento en LaTeX (\ruta{docs/documento/}):
migración desde el Word v0.1, plantilla, preámbulo y
\ruta{generar-docx.py} & Pasar a LaTeX el contenido que el equipo ya
había redactado en Word, con un archivo por sección & El documento
compila con pdflatex sin errores & Abdías Daniel Chávez Cruz \\
Claude Code (Anthropic) & Correcciones B11 y B5: 3.4, Figura 3.1,
portada, control de versiones, CONTRIBUTING.md, README, 2.1 y Anexo F &
Redactar los cambios que pide cada observación & Compilación sin errores
y comprobación de que no quedan los stakeholders retirados ni las ramas
por rol & Abdías Daniel Chávez Cruz \\
Claude Code (Anthropic) & Corrección B3: Anexo A (guion y minuta de la
entrevista) y fuentes de las reglas de 4.7 & Redactar el guion y la
minuta, y ligar cada regla con su hallazgo & Cada regla de entrevista de
4.7 cita un hallazgo de la minuta. La minuta está pendiente de validar
por el Product Owner contra sus notas de la entrevista & Abdías Daniel
Chávez Cruz \\
Claude Code (Anthropic) & Corrección B4: issues de las 30 historias,
números en 4.5, 2.3, minuta de la Review 1 y 12.2 & Registrar los issues
a partir de 4.5, resumir la retroalimentación del profesor y contar
commits, PR y revisiones & Datos comprobados contra GitHub y contra
\ruta{informacion/correcionesContexto.md} & Abdías Daniel Chávez Cruz \\
Claude Code (Anthropic) & Propuesta de alcance del Release v0.1.0
(observación B2) & Redactar la propuesta a partir de 4.5 y del
anteproyecto & Sumas de historias y puntos comprobadas contra 4.5.
Pendiente de revisión por el Product Owner y de negociación con el
cliente & Abdías Daniel Chávez Cruz \\
Claude Code (Anthropic) & Versiones v0.2.1 a v0.2.3, etiquetas de Git,
issues, ramas, Pull Requests, carpeta \ruta{informacion/} y esta
declaración & Aplicar los acuerdos de CONTRIBUTING.md en el repositorio &
Estado del repositorio comprobado en GitHub después de cada cambio &
Abdías Daniel Chávez Cruz \\
\pendiente{} & Contenido original de las secciones 1 a 4, redactado en
Word en octubre de 2026 & \pendiente{} & \pendiente{} & Cada autor declara
su parte \\
\end{longtable}
}

\textbf{Datos personales.} No se ingresaron datos de clientes en la
herramienta: el proyecto solo usa datos ficticios. El asistente trabajó
sobre el contenido del repositorio, que incluye los nombres y las
matrículas de los integrantes (sección 1.1).

\textbf{Concentración del trabajo.} Todo el uso del asistente estuvo a
cargo de un solo integrante. Para repartir ese conocimiento, el Tech Lead
explicará al equipo los cambios hechos en una sesión interna con minuta
(observación B6, fecha límite 18/11/2026).

\section{3. Visión y viabilidad del
producto}\label{visiuxf3n-y-viabilidad-del-producto}

\subsection{3.1 Problema y contexto del
cliente}\label{problema-y-contexto-del-cliente}

Servicio Automotriz Hermanos Ruiz, S.A. de C.V. es un taller que atiende
vehículos ligeros (automóviles particulares, motocicletas y camionetas)
de clientes particulares y de pequeñas flotillas empresariales. Tiene 4
bahías de trabajo y un estacionamiento para 5 vehículos en espera. Hoy
las citas se apuntan en un cuaderno, las órdenes se llenan en papel y la
comunicación con el cliente es por teléfono o WhatsApp (ver el proceso
actual en 3.5).

Según el anteproyecto y la entrevista con el cliente del 19/10/2026, los
problemas principales son:

\textbf{• Disputas por cobros:} las autorizaciones son verbales o por
mensaje y no queda registro de qué se aprobó ni por cuánto («dijo que
sí, pero cobra otra cosa»).

\textbf{• Reclamos por daños:} no hay checklist ni fotografías del
estado en que llega el vehículo.

\textbf{• Exceso de llamadas:} el cliente llama para saber el avance y
la recepcionista interrumpe su trabajo para averiguarlo.

\textbf{• Poca visibilidad:} la carga de trabajo y las bahías se asignan
de palabra, y hay vehículos detenidos sin justificación.

\textbf{• Historial disperso:} las órdenes se archivan en carpetas y es
difícil consultar reparaciones y garantías anteriores.

Hallazgos de la entrevista que no estaban en el anteproyecto:

• El cliente quiere una doble verificación con código de 6 dígitos al
autorizar cotizaciones, en forma remota o presencial (RN-02).

• El cliente puede aprobar o rechazar por separado los conceptos
opcionales de una cotización (RN-03).

• Si una refacción tarda varios días, el vehículo pasa al
estacionamiento para liberar la bahía (RN-06).

• Las citas pedidas en línea no apartan lugar hasta que la recepcionista
las confirma (RN-08).

• Varios mecánicos pueden colaborar en una orden, pero siempre hay un
responsable (RN-09).

Además, hay cortes de luz e internet, por lo que el personal necesita
seguir capturando sin conexión, y el jefe de mecánicos trabajará con 1 o
2 tabletas en un ambiente con ruido, polvo y manos sucias (4.3).

\subsection{3.2 Objetivos del producto e
indicadores}\label{objetivos-del-producto-e-indicadores}

{\def\LTcaptype{none} % do not increment counter
\begin{longtable}[]{@{}
  >{\raggedright\arraybackslash}p{(\linewidth - 8\tabcolsep) * \real{0.2349}}
  >{\raggedright\arraybackslash}p{(\linewidth - 8\tabcolsep) * \real{0.2030}}
  >{\raggedright\arraybackslash}p{(\linewidth - 8\tabcolsep) * \real{0.1710}}
  >{\raggedright\arraybackslash}p{(\linewidth - 8\tabcolsep) * \real{0.1708}}
  >{\raggedright\arraybackslash}p{(\linewidth - 8\tabcolsep) * \real{0.2202}}@{}}
\toprule\noalign{}
\begin{minipage}[b]{\linewidth}\raggedright
\textbf{Objetivo}
\end{minipage} & \begin{minipage}[b]{\linewidth}\raggedright
\textbf{Indicador}
\end{minipage} & \begin{minipage}[b]{\linewidth}\raggedright
\textbf{Situación actual}
\end{minipage} & \begin{minipage}[b]{\linewidth}\raggedright
\textbf{Meta}
\end{minipage} & \begin{minipage}[b]{\linewidth}\raggedright
\textbf{Cómo se medirá}
\end{minipage} \\
\midrule\noalign{}
\endhead
\bottomrule\noalign{}
\endlastfoot
Eliminar las disputas por cobros no autorizados & \% de órdenes cerradas
cuya cotización tiene autorización digital registrada (fecha, hora, IP y
monto) & Autorizaciones verbales o por mensaje, sin registro & 100 \% de
las órdenes cerradas & Reporte mensual del sistema: órdenes cerradas con
autorización confirmada / total de órdenes cerradas \\
Reducir las llamadas de seguimiento a recepción & Llamadas por día de
clientes que preguntan por el avance de su vehículo & Por medir:
bitácora de llamadas de una semana antes de liberar el portal del
cliente & Al menos 50 \% menos que la medición inicial & Bitácora de
llamadas de recepción, una semana antes y una semana después de la
liberación \\
Contar con evidencia del estado del vehículo al recibirlo & \% de
recepciones con checklist completo y las 4 fotografías obligatorias &
Órdenes en papel sin checklist ni fotografías & 100 \% de las
recepciones & Reporte mensual del sistema de recepciones completas \\
\end{longtable}
}

\subsection{3.3 Alcance}\label{alcance}

{\def\LTcaptype{none} % do not increment counter
\begin{longtable}[]{@{}
  >{\raggedright\arraybackslash}p{(\linewidth - 4\tabcolsep) * \real{0.3333}}
  >{\raggedright\arraybackslash}p{(\linewidth - 4\tabcolsep) * \real{0.3333}}
  >{\raggedright\arraybackslash}p{(\linewidth - 4\tabcolsep) * \real{0.3333}}@{}}
\toprule\noalign{}
\begin{minipage}[b]{\linewidth}\raggedright
\textbf{Dentro del alcance de la Fase 1}
\end{minipage} & \begin{minipage}[b]{\linewidth}\raggedright
\textbf{Dentro del alcance de la Fase 2}
\end{minipage} & \begin{minipage}[b]{\linewidth}\raggedright
\textbf{Fuera de alcance}
\end{minipage} \\
\midrule\noalign{}
\endhead
\bottomrule\noalign{}
\endlastfoot
• Acceso del personal por rol (HU-27)

• Citas en línea y agenda con confirmación manual (HU-03, HU-04, HU-06)

• Recepción con checklist y fotografías, también sin conexión (HU-01,
HU-25, HU-26)

• Asignación de bahía y mecánico, y diagnóstico (HU-07 a HU-09)

• Cotización desglosada y autorización del cliente con código 2FA
(HU-02, HU-12 a HU-16)

• Solicitud de refacciones y proveedores (HU-10, HU-17, HU-18)

• Aviso de fin de trabajos, cobro, entrega y crédito a flotillas (HU-20
a HU-23)

• Portal del cliente responsivo y aviso de privacidad (HU-28, HU-30) & •
Búsqueda rápida de clientes y vehículos (HU-05)

• Tablero simplificado de bahías para tableta (HU-11)

• Consulta de existencias de refacciones (HU-19)

• Consulta de garantías por el cliente (HU-24)

• Panel y reportes gerenciales (HU-29, RF-11)

• Envío del enlace por SMS (RF-02.3) & • Facturación electrónica (CFDI)
y contabilidad

• Pagos en línea con tarjeta

• Aplicación móvil nativa (iOS o Android)

• Vehículos pesados, tráileres y maquinaria (RN-07)

• Nómina y control de asistencia del personal \\
\end{longtable}
}

\subsection{3.4 Stakeholders}\label{stakeholders}

{\def\LTcaptype{none} % do not increment counter
\begin{longtable}[]{@{}
  >{\raggedright\arraybackslash}p{(\linewidth - 8\tabcolsep) * \real{0.1922}}
  >{\raggedright\arraybackslash}p{(\linewidth - 8\tabcolsep) * \real{0.2031}}
  >{\raggedright\arraybackslash}p{(\linewidth - 8\tabcolsep) * \real{0.2349}}
  >{\raggedright\arraybackslash}p{(\linewidth - 8\tabcolsep) * \real{0.1496}}
  >{\raggedright\arraybackslash}p{(\linewidth - 8\tabcolsep) * \real{0.2202}}@{}}
\toprule\noalign{}
\begin{minipage}[b]{\linewidth}\raggedright
\textbf{Stakeholder}
\end{minipage} & \begin{minipage}[b]{\linewidth}\raggedright
\textbf{Rol en el negocio}
\end{minipage} & \begin{minipage}[b]{\linewidth}\raggedright
\textbf{Interés en el sistema}
\end{minipage} & \begin{minipage}[b]{\linewidth}\raggedright
\textbf{Poder / interés}
\end{minipage} & \begin{minipage}[b]{\linewidth}\raggedright
\textbf{Estrategia de participación}
\end{minipage} \\
\midrule\noalign{}
\endhead
\bottomrule\noalign{}
\endlastfoot
Dueños del taller & Propietarios; deciden la inversión y aprueban el
proyecto. & Reducir disputas por cobros, cuidar la rentabilidad y la
reputación, y no pagar licencias. & Alto poder -- alto interés &
Gestionar de cerca mediante entrevistas y revisiones de avance en cada
entrega, validación de requisitos y aprobación formal del alcance. \\
Gerente & Supervisa la operación y autoriza entregas a crédito a
flotillas. & Consultar tableros de control, controlar saldos y autorizar
créditos con respaldo. & Alto poder -- alto interés & Gestionar de cerca
mediante sesiones de definición de reglas de crédito y de reportes, y
pruebas de aceptación de su flujo. \\
Recepcionista & Registra ingresos, confirma citas y atiende a clientes.
& Reducir llamadas de seguimiento y tener el registro de recepción y
fotografías en un solo lugar. & Bajo poder -- alto interés & Mantener
informada e involucrada, participando en las pruebas de usabilidad,
capacitación y retroalimentación continua. \\
Jefe de mecánicos & Centraliza diagnóstico, cotización y asignación de
bahías en la tableta. & Flujos cortos, botones grandes y operación sin
conexión. & Alto poder -- alto interés & Gestionar de cerca los
prototipos en el entrono real (tableta, pruebas con manos sucias) y
ajustes por iteración. \\
Mecánicos de bahía y de apoyo & Realizan reparaciones, inspecciones y
cubrimientos de turnos. & Recibir órdenes claras y enfocarse en el
trabajo físico & Bajo poder -- alto interés & Mantener informado con
capacitaciones breves y un canal para reportar dificultades. \\
Encargado de almacén & Controla requisiciones, proveedores, costos e
inventario de consumibles. & Directorio de proveedores, registro de
costos y fechas, y tickets de compra por orden & Bajo poder -- alto
interés & Mantener informado e involucrado mediante entrevistas sobre su
flujo de compras y pruebas del módulo M6. \\
Clientes particulares & Dueños de autos, motocicletas y vehículos. & Ver
el estado de su vehículo, autorizar presupuestos y recibir un cobro
transparente & Bajo poder -- alto interés & Mantener informado con las
pruebas de usabilidad con algunos clientes y alguna encuesta breve de
satisfacción al finalizar el servicio. \\
Clientes de flotillas empresariales & Empresas con vehículos que pueden
recibir entrega a crédito. & Condiciones de crédito claras y seguimiento
de varios vehículos & Poder medio -- interés medio & Mantener
satisfechos al definir con el gerente las reglas del crédito y validar
el flujo con un representante. \\
Equipo de desarrollo & Diseñan, construyen y prueban la plataforma. &
Entregar un sistema que cumpla los requisitos y la norma ISO/IEC 25010 &
Alto poder -- alto interés & Gestionar de cerca con reuniones
periódicas, control de versiones y revisiones internas de calidad. \\
Docente evaluador & Revisa el documento y el producto en curso. & Que el
proyecto cumpla los criterios y entregas del curso & Alto poder -- alto
interés & Mantener satisfecho entregando a tiempo, pidiendo asesoría en
los avances y pidiendo aclaraciones de dudas. \\
\end{longtable}
}

\figura{figuras/mapa-poder-interes.png}{width=0.8\linewidth}{3.1}{Mapa de poder/interés de los
stakeholders de la tabla anterior. Elaboración propia}{docs/03-vision/mapa-poder-interes.py}

\subsection{3.5 Proceso actual (as-is)}\label{proceso-actual-as-is}

\figura{figuras/proceso-as-is.png}{height=0.8\textheight}{3.2}{Diagrama de actividades del proceso actual (as-is) de
atención del taller, con los problemas identificados en cada etapa.
Elaboración propia}{docs/03-vision/proceso-as-is.puml}

El diagrama de actividades muestra el proceso actual de atención del
taller antes de implementar la solución tecnológica. Se organiza en
cuatro carriles (cliente, recepcionista, mecánico y registro manual) y
recorre el flujo completo: solicitud de cita, recepción del vehículo,
diagnóstico, cotización, autorización, reparación, entrega y archivo del
historial. Los recuadros rojos señalan los problemas detectados en cada
etapa, como el riesgo de errores en el registro de citas, los datos
incompletos en las órdenes en papel, la información dispersa e ilegible,
las demoras en la cotización manual, el exceso de llamadas y un
historial difícil de consultar.

\subsection{3.6 Estudio de viabilidad}\label{estudio-de-viabilidad}

{\def\LTcaptype{none} % do not increment counter
\begin{longtable}[]{@{}
  >{\raggedright\arraybackslash}p{(\linewidth - 4\tabcolsep) * \real{0.1708}}
  >{\raggedright\arraybackslash}p{(\linewidth - 4\tabcolsep) * \real{0.5513}}
  >{\raggedright\arraybackslash}p{(\linewidth - 4\tabcolsep) * \real{0.2779}}@{}}
\toprule\noalign{}
\begin{minipage}[b]{\linewidth}\raggedright
\textbf{Tipo}
\end{minipage} & \begin{minipage}[b]{\linewidth}\raggedright
\textbf{Análisis y evidencia}
\end{minipage} & \begin{minipage}[b]{\linewidth}\raggedright
\textbf{Conclusión (viable / viable con condiciones / no viable)}
\end{minipage} \\
\midrule\noalign{}
\endhead
\bottomrule\noalign{}
\endlastfoot
Técnica & La pila propuesta (Python con FastAPI, PostgreSQL y una PWA
con Service Workers e IndexedDB) es de código abierto y está bien
documentada. El mayor riesgo técnico es la operación sin conexión con
sincronización posterior (RNF-08, RNF-09, R-02), que se validará con un
prototipo en el Sprint 1. & Viable con condiciones: validar la
sincronización sin conexión en el Sprint 1. \\
Económica & Costo de licencias de \$0: todas las herramientas son de
código abierto o de planes educativos (RNF-07). El taller usa las
computadoras que ya tiene en recepción y almacén, y 1 o 2 tabletas para
el jefe de mecánicos. Los planes gratuitos de correo y almacenamiento
tienen límites (R-05, R-06). & Viable con condiciones: el taller aporta
las tabletas y el sistema respeta los límites de los planes
gratuitos. \\
Operativa & El personal tiene experiencia tecnológica baja a media
(4.3.1). Se atiende con flujos de pocos pasos y botones grandes (RNF-02,
RNF-03), capacitación breve y pruebas de usabilidad con el personal. Los
dueños y el gerente apoyan el proyecto porque reduce disputas y llamadas
(3.4). & Viable con condiciones: capacitar y probar con el personal
antes de la Demo final. \\
Legal & Se tratan datos de identificación y contacto de los clientes
(nombre, teléfono y correo) y datos de sus vehículos, incluidas
fotografías; no se tratan datos sensibles. Las finalidades y medidas de
seguridad (control de acceso por rol, HTTPS, tokens temporales y aviso
de privacidad) se detallan en 4.8 (Review 2), conforme a la Ley Federal de Protección de
Datos Personales en Posesión de los Particulares (DOF 20-03-2025). En
desarrollo y pruebas solo se usan datos ficticios (RNF-05). & Viable \\
De calendario & El semestre tiene 7 sprints (S0 a S6) hasta la Demo
final de la semana 16. El backlog tiene 30 historias; el release v0.1.0
incluye el flujo núcleo y las historias Should o Could pueden pasar a la
Fase 2 si la velocidad del equipo no alcanza. & Viable con condiciones:
priorizar las historias Must. \\
\end{longtable}
}

\textbf{Conclusión general:} el proyecto es viable con condiciones:
validar en el Sprint 1 la operación sin conexión, que el taller aporte
las tabletas, respetar los límites de los planes gratuitos y priorizar
las historias Must del flujo núcleo.

\section{4. Especificación de requisitos de software
(SRS)}\label{especificaciuxf3n-de-requisitos-de-software-srs}

Estructura adaptada de ISO/IEC/IEEE 29148:2018. Cada requisito tiene un
identificador único que se utiliza para la trazabilidad.

\subsection{4.1 Propósito, alcance y audiencia del
SRS}\label{propuxf3sito-alcance-y-audiencia-del-srs}

\textbf{Propósito:} este SRS describe qué debe hacer el Sistema de
Órdenes de Servicio y Citas y con qué calidad, para que el cliente lo
valide y lo firme en la Review 2, y para que el equipo lo use como base
del diseño, la implementación y las pruebas.

\textbf{Alcance:} cubre el ciclo de atención del taller (citas,
recepción, diagnóstico, cotización y autorización, refacciones, cobro y
entrega) y el portal web del cliente, según el alcance de la sección
3.3.

\textbf{Audiencia:} los dueños y el gerente del taller (aprobación), el
personal que usará el sistema (validación de sus flujos), el equipo de
desarrollo y el docente evaluador.

\textbf{Convenciones:} cada requisito tiene un identificador único: RF
(funcional), RNF (no funcional), RN (regla de negocio), HU (historia de
usuario) y CU (caso de uso). Estos identificadores se usan en la matriz
de trazabilidad (5.3).

\subsection{4.2 Glosario del dominio}\label{glosario-del-dominio}

{\def\LTcaptype{none} % do not increment counter
\begin{longtable}[]{@{}
  >{\raggedright\arraybackslash}p{(\linewidth - 4\tabcolsep) * \real{0.2349}}
  >{\raggedright\arraybackslash}p{(\linewidth - 4\tabcolsep) * \real{0.5342}}
  >{\raggedright\arraybackslash}p{(\linewidth - 4\tabcolsep) * \real{0.2309}}@{}}
\toprule\noalign{}
\begin{minipage}[b]{\linewidth}\raggedright
\textbf{Término}
\end{minipage} & \begin{minipage}[b]{\linewidth}\raggedright
\textbf{Definición}
\end{minipage} & \begin{minipage}[b]{\linewidth}\raggedright
\textbf{Sinónimos que usa el cliente}
\end{minipage} \\
\midrule\noalign{}
\endhead
\bottomrule\noalign{}
\endlastfoot
Autorización & Aceptación de una cotización por el cliente; registra
fecha, hora, IP, monto y el código de verificación. & Aprobación, firma
digital \\
Bahía & Espacio de trabajo donde se repara un vehículo. El taller tiene
4 y cada una aloja un vehículo a la vez. & --- \\
Checklist de inspección & Lista de verificación del estado de
carrocería, interiores y accesorios al recibir el vehículo. & Lista de
inspección \\
Cita & Solicitud para atender un vehículo en una fecha y hora; queda
«Pendiente» hasta que la recepcionista la confirma. & --- \\
Código 2FA & Código de 6 dígitos, de un solo uso y vigente 5 minutos,
que se envía al correo del cliente para confirmar una autorización. &
Código de verificación, doble verificación \\
Compra & Adquisición de refacciones a un proveedor, respaldada por un
ticket. & Ticket de compra \\
Concepto & Renglón de una cotización: mano de obra, refacción o insumo.
Puede ser opcional. & Rubro \\
Cotización & Costo desglosado de mano de obra, refacciones e insumos de
una orden; vigente 5 días naturales. & Presupuesto \\
Diagnóstico & Descripción técnica de la falla y de las reparaciones
necesarias que registra el jefe de mecánicos. & --- \\
Enlace seguro & Liga web única que recibe el cliente para consultar su
vehículo y autorizar cotizaciones sin usuario ni contraseña. & Token de
acceso \\
Entrega a crédito & Entrega del vehículo con saldo pendiente; solo para
flotillas registradas y con autorización del gerente. & Crédito \\
Flotilla & Empresa cliente con varios vehículos que puede recibir
entrega a crédito. & Flotilla empresarial \\
Folio & Identificador único de una orden de servicio. & --- \\
Fotografía de recepción & Evidencia digital del estado del vehículo al
llegar. Son obligatorias cuatro: frente, trasera, costado izquierdo y
costado derecho. & --- \\
Garantía & Cobertura de 30 días sobre la mano de obra, contados desde la
entrega del vehículo. & --- \\
Insumo & Material consumible que se usa en la reparación, como aceite o
líquidos. & Consumible \\
Jefe de mecánicos & Líder técnico: diagnostica, cotiza, asigna bahías y
mecánicos y solicita refacciones desde la tableta. & Líder técnico \\
Mecánico responsable & Mecánico titular de una orden; otros mecánicos
pueden colaborar sin cambiar la titularidad. & Mecánico principal \\
Orden de servicio & Registro de todo el trabajo sobre un vehículo, desde
la recepción hasta la entrega. & Orden \\
Proveedor & Refaccionaria local a la que el almacén compra piezas. &
Refaccionaria \\
Recepción vehicular & Registro de la llegada del vehículo: kilometraje,
nivel de combustible, checklist y fotografías. & Check-in, registro de
entrada \\
Re-cotización & Nueva versión de la cotización cuando los trabajos
adicionales aumentan el monto más de 10 \%. & --- \\
Refacción & Pieza que se instala en el vehículo. & Pieza, repuesto \\
Saldo pendiente & Diferencia entre el total autorizado de la orden y lo
que el cliente ha pagado. & --- \\
Solicitud de refacción & Pedido de piezas que el jefe de mecánicos hace
al almacén para una orden. & Requisición \\
\end{longtable}
}

\subsection{4.3 Descripción general}\label{descripciuxf3n-general}

\subsubsection{4.3.1 Usuarios y sus
características}\label{usuarios-y-sus-caracteruxedsticas}

El sistema será utilizado por distintos perfiles dentro de la estructura
operativa del negocio y la parte gerencial del taller mecánico, además
de los clientes finales del negocio. La siguiente tabla detalla las
características, nivel técnico y contexto de uso:

{\def\LTcaptype{none} % do not increment counter
\begin{longtable}[]{@{}
  >{\raggedright\arraybackslash}p{(\linewidth - 6\tabcolsep) * \real{0.2138}}
  >{\raggedright\arraybackslash}p{(\linewidth - 6\tabcolsep) * \real{0.2990}}
  >{\raggedright\arraybackslash}p{(\linewidth - 6\tabcolsep) * \real{0.2138}}
  >{\raggedright\arraybackslash}p{(\linewidth - 6\tabcolsep) * \real{0.2734}}@{}}
\toprule\noalign{}
\begin{minipage}[b]{\linewidth}\raggedright
\textbf{Tipo de usuario}
\end{minipage} & \begin{minipage}[b]{\linewidth}\raggedright
\textbf{Tareas principales}
\end{minipage} & \begin{minipage}[b]{\linewidth}\raggedright
\textbf{Experiencia tecnológica}
\end{minipage} & \begin{minipage}[b]{\linewidth}\raggedright
\textbf{Dispositivo y contexto de uso}
\end{minipage} \\
\midrule\noalign{}
\endhead
\bottomrule\noalign{}
\endlastfoot
Recepcionista & Registro de clientes/vehículos, check-in con fotos,
apertura de órdenes, confirmación de citas web, cobranza y entrega. &
Media: Manejo fluido de navegación web, correo electrónico y paquetes de
oficina. & Computadora de escritorio en mostrador de recepción. Entorno
de oficina con teclado/mouse. \\
Jefe de Mecánicos (Líder Técnico) & Diagnóstico técnico, elaboración de
cotizaciones, asignación de bahías/mecánicos, requisición de refacciones
y sellado final. & Baja--Media: Familiarizado con smartphones
personales; poca experiencia con sistemas de gestión. & Tabletas
táctiles exclusivas de taller (1 o 2 dispositivos). Entorno con ruido,
polvo y manipulación con manos sucias/grasosas. \\
Mecánico en Bahía & Ejecución de reparaciones físicas, notificación de
hallazgos adicionales o solicitud de refacciones al Jefe de Mecánicos. &
Baja: Uso de teléfonos inteligentes para uso personal cotidiano. & Sin
dispositivo asignado directo; interactúa a través del Jefe de Mecánicos
para evitar deterioro de hardware \\
Encargado de Almacén & Atender solicitudes de piezas, consultar
directorio de proveedores locales, registrar compras/tickets e historial
de insumos. & Baja--Media\textbf{:} Uso básico de computadora para
consulta de inventarios. & Computadora de escritorio o terminal dedicada
en el área de almacén/refacciones. \\
Gerente (Dueño) & Monitoreo del Dashboard de KPIs (ocupación, ingresos),
aprobación de entregas a crédito a flotillas y excepciones. & Media:
Manejo de smartphones y computadoras para administración del negocio. &
Laptop personal o dispositivo móvil (smartphone/tableta) con acceso
remoto/local \\
Cliente & Solicitud de citas en línea, seguimiento en vivo del servicio
y aprobación/rechazo de cotizaciones mediante código 2FA. & Variable
(Baja a Alta): Desde usuarios sin conocimientos técnicos hasta usuarios
avanzados. & Dispositivo móvil personal (iOS/Android) vía navegador web
progresivo (PWA) sin instalar aplicaciones. \\
\end{longtable}
}

\subsubsection{4.3.2 Restricciones, supuestos y
dependencias}\label{restricciones-supuestos-y-dependencias}

\begin{enumerate}
\def\labelenumi{\arabic{enumi}.}
\item
  \textbf{Restricciones:}
\end{enumerate}

\begin{itemize}
\item
  Presupuesto y Licenciamiento: El desarrollo debe realizarse
  estrictamente con herramientas, lenguajes y frameworks de código
  abierto (\emph{Open Source} como Python y FastAPI), garantizando cero
  costos de licencias o rentas de software.
\item
  Infraestructura Física: El taller cuenta únicamente con 4 bahías de
  trabajo físicas (capacidad estricta de 1 vehículo por bahía) y un
  estacionamiento con capacidad máxima de 5 vehículos en espera.
\item
  Dispositivos Hardware: La captura digital en taller se limita a máximo
  2 tabletas táctiles centralizadas por el Jefe de Mecánicos para
  proteger la inversión y evitar daños.
\end{itemize}

\begin{itemize}
\item
  Alcance del Vehículo\textbf{:} La plataforma atenderá exclusivamente
  vehículos ligeros (automóviles particulares, motocicletas y
  camionetas), excluyendo maquinaria pesada o industrial.
\end{itemize}

\begin{enumerate}
\def\labelenumi{\arabic{enumi}.}
\setcounter{enumi}{1}
\item
  \textbf{Supuestos:}
\end{enumerate}

\begin{itemize}
\item
  Se asume que el taller mantiene activa una conexión de red local
  (Wi-Fi) para la sincronización entre computadoras de mostrador y
  tabletas de taller en condiciones normales.
\item
  Se asume que los clientes cuentan con un correo electrónico para
  recibir el enlace de seguimiento y los códigos de verificación de 6
  dígitos (2FA); el número celular se usa para el envío opcional del
  enlace por SMS (RF-02.3).
\item
  Se asume que los datos ingresados durante la fase de desarrollo y
  pruebas son 100\% ficticios para dar cumplimiento a la Ley Federal de
  Protección de Datos Personales.
\end{itemize}

\begin{enumerate}
\def\labelenumi{\arabic{enumi}.}
\setcounter{enumi}{2}
\item
  \textbf{Dependencias:}
\end{enumerate}

\begin{itemize}
\item
  Infraestructura Eléctrica y Red: Operación dependiente de baterías en
  laptops/tabletas y arquitectura \emph{Local-First / PWA} (Service
  Workers e IndexedDB) para permitir la captura de datos \emph{offline}
  durante cortes de luz o fallas de internet.
\item
  Proveedores Externos de Comunicación: Dependencia del servicio de
  correo electrónico para entregar enlaces y códigos 2FA al cliente, y
  de un servicio de SMS para el envío opcional del enlace.
\item
  Disponibilidad de Refaccionarias Locales: La adquisición de piezas
  fuera de stock depende del directorio de proveedores locales e
  historial de compras del almacén.
\end{itemize}

\subsection{4.4 Requisitos funcionales}\label{requisitos-funcionales}

{\def\LTcaptype{none} % do not increment counter
\begin{longtable}[]{@{}
  >{\raggedright\arraybackslash}p{(\linewidth - 8\tabcolsep) * \real{0.0962}}
  >{\raggedright\arraybackslash}p{(\linewidth - 8\tabcolsep) * \real{0.4444}}
  >{\raggedright\arraybackslash}p{(\linewidth - 8\tabcolsep) * \real{0.1174}}
  >{\raggedright\arraybackslash}p{(\linewidth - 8\tabcolsep) * \real{0.1497}}
  >{\raggedright\arraybackslash}p{(\linewidth - 8\tabcolsep) * \real{0.1923}}@{}}
\toprule\noalign{}
\begin{minipage}[b]{\linewidth}\raggedright
\textbf{ID}
\end{minipage} & \begin{minipage}[b]{\linewidth}\raggedright
\textbf{Requisito funcional}
\end{minipage} & \begin{minipage}[b]{\linewidth}\raggedright
\textbf{Módulo}
\end{minipage} & \begin{minipage}[b]{\linewidth}\raggedright
\textbf{Prioridad (MoSCoW)}
\end{minipage} & \begin{minipage}[b]{\linewidth}\raggedright
\textbf{Historias relacionadas}
\end{minipage} \\
\midrule\noalign{}
\endhead
\bottomrule\noalign{}
\endlastfoot
RF-01 & El sistema deberá permitir autenticar mediante credenciales
seguras al personal del taller según su rol (Gerente, Recepcionista,
Jefe de Mecánicos y Encargado de Almacén). & M1 -- Autenticación y
Seguridad & Must & HU-27 \\
RF-02.1 & El sistema deberá permitir al cliente consultar el estado de
su vehículo y autorizar cotizaciones mediante un enlace web seguro
enviado por correo electrónico o por SMS (RF-02.3), sin instalar
aplicaciones. & M1 -- Autenticación y Seguridad & Must & HU-02, HU-28 \\
RF-02.2 & El sistema deberá de expirar el enlace del cliente de manera
automática una vez que el vehículo sea entregado y saldado. & M1 --
Autenticación y Seguridad & Should & HU-02, HU-23 \\
RF-02.3 & El sistema deberá de permitir enviar el enlace del cliente
también por SMS. & M1 -- Autenticación y Seguridad & Could & HU-02 \\
RF-03.1 & El sistema deberá registrar, buscar y actualizar la
información de clientes y vehículos (placas, marca, modelo, año y VIN
(opcional)). & M3 -- Recepción E Historial Inicial & Must & HU-05,
HU-06 \\
RF-03.2 & El sistema deberá mantener el historial de servicios de cada
vehículo, con fotografías de recepción, diagnósticos técnicos y
garantías aplicables (30 días en mano de obra). & M7 -- Portal del
Cliente, Cobro y Gerencia & Should & HU-23, HU-24 \\
RF-04 & El sistema deberá registrar en la recepción el kilometraje, el
nivel de combustible y un checklist de carrocería, y exigir fotografías
digitales del estado inicial de la unidad. & M3 -- Recepción E Historial
Inicial & Must & HU-01, HU-25, HU-26 \\
RF-05 & El sistema deberá permitir abrir órdenes de servicio asociadas a
una de las 4 bahías y a un mecánico responsable, controlando los estados
operativos del taller. & M5 -- Taller, Bahías y Expediente Técnico &
Must & HU-07, HU-08, HU-09, HU-11 \\
RF-06 & El sistema deberá generar cotizaciones desglosadas (mano de obra
y refacciones/insumos) con vigencia de 5 días naturales y requerir firma
digital validada con un código enviado al correo del cliente, de forma
remota o presencial. & M4 Cotización y Contratos 2FA & Must & HU-02,
HU-12, HU-13, HU-14, HU-15 \\
RF-07 & El sistema deberá bloquear el avance del trabajo y exigir una
nueva autorización con verificación cuando los trabajos adicionales
incrementen el monto inicial en más de 10 \%. & M4 Cotización y
Contratos 2FA & Must & HU-16 \\
RF-08.1 & El sistema deberá permitir al jefe de mecánicos solicitar al
almacén las piezas que requieran los mecánicos, vinculadas a una orden
de servicio. & M6 -- Almacén y Refacciones & Should & HU-10, HU-19 \\
RF-08.2 & El sistema deberá permitir al encargado de almacén mantener un
directorio de proveedores locales (teléfonos o correos) y registrar
costos y fechas de llegada de las piezas. & M6 -- Almacén y Refacciones
& Should & HU-17, HU-18 \\
RF-08.3 & El sistema deberá de emitir un ticket e historial de compras
asociado a cada orden de servicio. & M6 -- Almacén y Refacciones & Could
& HU-18 \\
RF-09 & El sistema deberá ofrecer una agenda digital por día y bahía,
donde las citas solicitadas en línea quedan ``Pendientes'' hasta que la
recepcionista verifique disponibilidad y las confirme manualmente & M2
-- Citas y Solicitudes & Should & HU-03, HU-04, HU-06 \\
RF-10 & El sistema deberá impedir la entrega física del vehículo con
saldo pendiente, a menos que el gerente autorice la entrega a crédito
para vehículos de flotillas empresariales previamente registradas. & M7
-- Portal del Cliente, Cobro y Gerencia & Must & HU-21, HU-22, HU-23 \\
RF-11 & El sistema deberá ofrecer al gerente un panel con las órdenes
abiertas y su estado, la ocupación de las 4 bahías, los tiempos de
estancia y los ingresos acumulados. & M7 -- Portal del Cliente, Cobro y
Gerencia & Should & HU-29 \\
\end{longtable}
}

\subsection{4.5 Historias de usuario (Product
Backlog)}\label{historias-de-usuario-product-backlog}

{\def\LTcaptype{none} % do not increment counter
\begin{longtable}[]{@{}
  >{\raggedright\arraybackslash}p{(\linewidth - 10\tabcolsep) * \real{0.1389}}
  >{\raggedright\arraybackslash}p{(\linewidth - 10\tabcolsep) * \real{0.4124}}
  >{\raggedright\arraybackslash}p{(\linewidth - 10\tabcolsep) * \real{0.1068}}
  >{\raggedright\arraybackslash}p{(\linewidth - 10\tabcolsep) * \real{0.0855}}
  >{\raggedright\arraybackslash}p{(\linewidth - 10\tabcolsep) * \real{0.1282}}
  >{\raggedright\arraybackslash}p{(\linewidth - 10\tabcolsep) * \real{0.1282}}@{}}
\toprule\noalign{}
\begin{minipage}[b]{\linewidth}\raggedright
\textbf{ID / issue}
\end{minipage} & \begin{minipage}[b]{\linewidth}\raggedright
\textbf{Historia (Como\ldots{} quiero\ldots{} para\ldots)}
\end{minipage} & \begin{minipage}[b]{\linewidth}\raggedright
\textbf{Prioridad}
\end{minipage} & \begin{minipage}[b]{\linewidth}\raggedright
\textbf{Puntos}
\end{minipage} & \begin{minipage}[b]{\linewidth}\raggedright
\textbf{Release}
\end{minipage} & \begin{minipage}[b]{\linewidth}\raggedright
\textbf{Estado}
\end{minipage} \\
\midrule\noalign{}
\endhead
\bottomrule\noalign{}
\endlastfoot
HU-01 / \#31 & Como recepcionista del taller, quiero registrar la
recepción del vehículo adjuntando kilometraje, nivel de combustible,
checklist de carrocería y fotografías del estado inicial, para contar
con un respaldo digital inalterable que evite reclamos infundados por
daños o faltantes al entregar la unidad. & Must & 3 & v0.1.0 &
Pendiente \\
HU-02 / \#32 & Como cliente del taller, quiero recibir un enlace
web seguro con la cotización detallada (mano de obra y refacciones) y
botones claros de aprobación o rechazo, para revisar y autorizar los
trabajos desde mi teléfono sin necesidad de realizar llamadas ni acudir
físicamente al taller & Must & 3 & v0.1.0 & Pendiente \\
HU-03 / \#33 & Como cliente del taller, quiero solicitar una cita
en línea seleccionando el día y la hora disponible, para apartar un
espacio sin necesidad de realizar llamadas telefónicas. & Must & 3 &
v0.1.0 & Pendiente \\
HU-04 / \#34 & Como recepcionista del taller, quiero consultar la
agenda digital por día y por bahía, para verificar las solicitudes
pendientes de los clientes y realizar su confirmación manual. & Must & 3
& v0.1.0 & Pendiente \\
HU-05 / \#35 & Como recepcionista del taller, quiero buscar a un
cliente o vehículo registrado mediante su número telefónico o placas,
para agilizar el proceso de recepción en visitas recurrentes. & Should &
2 & Fase 2 & Pendiente \\
HU-06 / \#36 & Como recepcionista del taller, quiero registrar un
vehículo con cita previa seleccionando su folio de origen, para que el
sistema precargue automáticamente los datos del cliente y del auto. &
Should & 2 & v0.1.0 & Pendiente \\
HU-07 / \#37 & Como recepcionista del taller, quiero ver la alerta
de "Sin Bahías Disponibles" cuando las 4 bahías físicas estén ocupadas,
para ubicar temporalmente el vehículo en el estacionamiento de
recepción. & Must & 3 & v0.1.0 & Pendiente \\
HU-08 / \#38 & Como jefe de mecánicos, quiero asignar la orden de
servicio a una de las 4 bahías físicas disponibles y a un mecánico
responsable, para iniciar formalmente los trabajos en el taller. & Must
& 5 & v0.1.0 & Pendiente \\
HU-09 / \#39 & Como jefe de mecánicos, quiero registrar el
diagnóstico técnico y las observaciones iniciales en la tableta, para
cambiar el estado de la orden a "En Diagnóstico". & Must & 3 & v0.1.0 &
Pendiente \\
HU-10 / \#40 & Como jefe de mecánicos, quiero registrar las
refacciones e insumos que los mecánicos detecten durante la revisión,
para que el área de almacén gestione su cotización o pedido. & Should &
3 & v0.1.0 & Pendiente \\
HU-11 / \#41 & Como jefe de mecánicos, quiero visualizar el
estatus actual de las 4 bahías de trabajo en una interfaz simplificada
para tabletas, para conocer rápidamente qué unidades están listas o en
proceso. & Should & 3 & Fase 2 & Pendiente \\
HU-12 / \#42 & Como recepcionista o jefe de mecánicos, quiero
generar un presupuesto desglosado de mano de obra y refacciones con
vigencia de 5 días naturales, para enviarlo de forma transparente al
cliente. & Must & 5 & v0.1.0 & Pendiente \\
HU-13 / \#43 & Como cliente del taller, quiero seleccionar o
deseleccionar servicios opcionales dentro de mi cotización antes de
aprobarla, para adaptar el presupuesto a mis necesidades financieras. &
Should & 3 & v0.1.0 & Pendiente \\
HU-14 / \#44 & Como cliente del taller, quiero realizar la firma
digital de aprobación mediante un código de verificación de 2 pasos
enviado a mi correo, para autorizar los trabajos de forma segura ya sea
de forma presencial o remota. & Must & 5 & v0.1.0 & Pendiente \\
HU-15 / \#45 & Como cliente del taller, quiero presionar un botón
de rechazo en la cotización digital, para cancelar la orden y coordinar
la recolección del vehículo sin autorizar reparaciones. & Should & 2 &
v0.1.0 & Pendiente \\
HU-16 / \#46 & Como cliente del taller, quiero que se me pida una
nueva autorización con código 2FA cuando los costos imprevistos superen
el 10 \% del presupuesto autorizado, para no pagar trabajos que no
aprobé. & Must & 5 & v0.1.0 & Pendiente \\
HU-17 / \#47 & Como encargado de almacén, quiero consultar y
administrar un directorio de proveedores locales con teléfonos y correos
electrónicos, para solicitar refacciones especiales rápidamente. &
Should & 3 & v0.1.0 & Pendiente \\
HU-18 / \#48 & Como encargado de almacén, quiero registrar los
costos, fechas de llegada y tickets de compra de las refacciones
asociadas a una orden, para mantener un control financiero estricto. &
Should & 3 & v0.1.0 & Pendiente \\
HU-19 / \#49 & Como jefe de mecánicos, quiero verificar el stock
de refacciones comunes en el inventario interno, para saber si se
dispone de la pieza de forma inmediata. & Could & 2 & Fase 2 &
Pendiente \\
HU-20 / \#50 & Como recepcionista, quiero recibir una notificación
automática cuando los trabajos del vehículo finalicen, para avisar al
cliente que su unidad está lista para entrega. & Should & 2 & v0.1.0 &
Pendiente \\
HU-21 / \#51 & Como gerente del taller, quiero que el sistema
impida registrar la entrega de un vehículo cuya orden tenga saldo
pendiente, para evitar pérdidas financieras al taller. & Must & 3 &
v0.1.0 & Pendiente \\
HU-22 / \#52 & Como gerente del taller, quiero autorizar de forma
excepcional la entrega a crédito únicamente para vehículos de flotillas
empresariales registradas, para cumplir con los convenios comerciales. &
Should & 3 & v0.1.0 & Pendiente \\
HU-23 / \#53 & Como recepcionista, quiero registrar el pago total
del servicio y cerrar formalmente la orden, para liberar la bahía de
trabajo y archivar el historial del vehículo. & Must & 3 & v0.1.0 &
Pendiente \\
HU-24 / \#54 & Como cliente, quiero consultar la vigencia de la
garantía de 30 días en mano de obra asociada a mi servicio anterior,
para reclamar ajustes en caso de incidencias. & Could & 2 & Fase 2 &
Pendiente \\
HU-25 / \#55 & Como recepcionista o jefe de mecánicos, quiero
registrar la llegada de autos, tomar fotografías y hacer diagnósticos
sin conexión a internet usando la batería del equipo, para que la
operación no se detenga ante fallas eléctricas. & Must & 8 & v0.1.0 &
Pendiente \\
HU-26 / \#56 & Como recepcionista o jefe de mecánicos, quiero que
los registros y fotografías que capturé sin conexión se sincronicen
automáticamente cuando vuelva la red, para no perder ni duplicar
información. & Must & 8 & v0.1.0 & Pendiente \\
HU-27 / \#57 & Como gerente del taller, quiero dar de alta al
personal y asignarle su rol (Gerente, Recepcionista, Jefe de Mecánicos o
Encargado de Almacén), para que cada quien acceda solo a la información
que le corresponde. & Must & 3 & v0.1.0 & Pendiente \\
HU-28 / \#58 & Como cliente, quiero visualizar la plataforma web
desde mi teléfono celular de forma responsiva y sin instalar
aplicaciones adicionales, para consultar mi estatus con facilidad. &
Must & 5 & v0.1.0 & Pendiente \\
HU-29 / \#59 & Como gerente, quiero generar reportes mensuales de
ingresos, servicios realizados y tiempos de ocupación de las 4 bahías,
para evaluar el rendimiento general del taller. & Could & 5 & Fase 2 &
Pendiente \\
HU-30 / \#60 & Como cliente del taller, quiero conocer el aviso de
privacidad antes de proporcionar mis datos, para saber para qué se
usarán conforme a la LFPDPPP. & Must & 3 & v0.1.0 & Pendiente \\
\end{longtable}
}

Detalle de las historias de usuario

\textbf{HU-01: Recepción e Inspección con Evidencia Fotográfica}

{\def\LTcaptype{none} % do not increment counter
\begin{longtable}[]{@{}
  >{\raggedright\arraybackslash}p{(\linewidth - 2\tabcolsep) * \real{0.2778}}
  >{\raggedright\arraybackslash}p{(\linewidth - 2\tabcolsep) * \real{0.7222}}@{}}
\toprule\noalign{}
\begin{minipage}[b]{\linewidth}\raggedright
\textbf{Campo}
\end{minipage} & \begin{minipage}[b]{\linewidth}\raggedright
\textbf{Contenido}
\end{minipage} \\
\midrule\noalign{}
\endhead
\bottomrule\noalign{}
\endlastfoot
ID e issue & HU-01 \\
Historia & Como recepcionista del taller, quiero registrar la recepción
del vehículo adjuntando kilometraje, nivel de combustible, checklist de
carrocería y fotografías del estado inicial, para contar con un respaldo
digital inalterable que evite reclamos infundados por daños o faltantes
al entregar la unidad. \\
Criterio de aceptación 1 & Dado que la recepcionista se encuentra en el
módulo de recepción vehicular, cuando intenta registrar una unidad
nueva, entonces el sistema exige obligatoriamente adjuntar al menos 4
fotografías (frente, parte trasera, costado izquierdo y costado
derecho). \\
Criterio de aceptación 2 & Dado que se capturan las evidencias, cuando
se finaliza el formulario de inspección, entonces el sistema genera una
firma visual o sello digital con fecha y hora exacta de recepción,
bloqueando el avance al estado "En Diagnóstico" si falta completar algún
punto del checklist obligatorio. \\
Reglas de negocio & RN-12 (Evidencia fotográfica de recepción), RN-07
(Clasificación de vehículos autorizados). \\
Wireframe & W-01 (4.10, Review 2) \\
Casos de prueba & CPA-01 a CPA-04 \\
\end{longtable}
}

\textbf{HU-02: Autorización Digital de Cotización por Parte del Cliente}

{\def\LTcaptype{none} % do not increment counter
\begin{longtable}[]{@{}
  >{\raggedright\arraybackslash}p{(\linewidth - 2\tabcolsep) * \real{0.2778}}
  >{\raggedright\arraybackslash}p{(\linewidth - 2\tabcolsep) * \real{0.7222}}@{}}
\toprule\noalign{}
\begin{minipage}[b]{\linewidth}\raggedright
\textbf{Campo}
\end{minipage} & \begin{minipage}[b]{\linewidth}\raggedright
\textbf{Contenido}
\end{minipage} \\
\midrule\noalign{}
\endhead
\bottomrule\noalign{}
\endlastfoot
ID e issue & HU-02 \\
Historia & Como cliente del taller, quiero recibir un enlace web seguro
con la cotización detallada (mano de obra y refacciones) y botones
claros de aprobación o rechazo, para revisar y autorizar los trabajos
desde mi teléfono sin necesidad de realizar llamadas ni acudir
físicamente al taller \\
Criterio de aceptación 1 & Dado que la cotización ha sido generada por
el taller, cuando el sistema envía el aviso, entonces el cliente recibe
un enlace por correo electrónico (o por SMS cuando se habilite RF-02.3)
que no requiere registro previo ni contraseña, operando mediante un
token temporal firmado. \\
Criterio de aceptación 2 & Dado que el cliente revisa el desglose en su
dispositivo móvil, cuando presiona el botón "Autorizar" e ingresa un
código 2FA válido (HU-14), entonces el sistema registra de forma
inalterable la dirección IP, fecha, hora y el monto final aceptado,
cambiando automáticamente el estado de la orden a "Autorizada" y
notificando al mecánico responsable. \\
Reglas de negocio & RN-01 (Expiración y vigencia de enlaces), RN-02
(Doble verificación en autorizaciones), RN-11 (Vigencia de
cotizaciones). \\
Wireframe & W-04 (4.10, Review 2) \\
Casos de prueba & CPA-10, CPA-30 \\
\end{longtable}
}

\textbf{HU-03: Solicitud de Cita en Línea}

{\def\LTcaptype{none} % do not increment counter
\begin{longtable}[]{@{}
  >{\raggedright\arraybackslash}p{(\linewidth - 2\tabcolsep) * \real{0.2778}}
  >{\raggedright\arraybackslash}p{(\linewidth - 2\tabcolsep) * \real{0.7222}}@{}}
\toprule\noalign{}
\begin{minipage}[b]{\linewidth}\raggedright
\textbf{Campo}
\end{minipage} & \begin{minipage}[b]{\linewidth}\raggedright
\textbf{Contenido}
\end{minipage} \\
\midrule\noalign{}
\endhead
\bottomrule\noalign{}
\endlastfoot
ID e issue & HU-03 \\
Historia & Como cliente del taller, quiero solicitar una cita en línea
seleccionando el día y la hora disponible, para apartar un espacio sin
necesidad de realizar llamadas telefónicas. \\
Criterio de aceptación 1 & Dado que el cliente accede al portal web,
cuando selecciona fecha, hora y tipo de servicio, entonces el sistema
valida que exista disponibilidad en la agenda del taller. \\
Criterio de aceptación 2 & Dado que se envía la solicitud, cuando se
registra con éxito, entonces queda en estado "Pendiente" a la espera de
la confirmación manual de la recepcionista. \\
Reglas de negocio & RN-08 (Confirmación manual de citas web). \\
Wireframe & W-05 (4.10, Review 2) \\
Casos de prueba & CPA-31, CPA-34 \\
\end{longtable}
}

\textbf{HU-04: Consulta y Confirmación Manual de Citas}

{\def\LTcaptype{none} % do not increment counter
\begin{longtable}[]{@{}
  >{\raggedright\arraybackslash}p{(\linewidth - 2\tabcolsep) * \real{0.2778}}
  >{\raggedright\arraybackslash}p{(\linewidth - 2\tabcolsep) * \real{0.7222}}@{}}
\toprule\noalign{}
\begin{minipage}[b]{\linewidth}\raggedright
\textbf{Campo}
\end{minipage} & \begin{minipage}[b]{\linewidth}\raggedright
\textbf{Contenido}
\end{minipage} \\
\midrule\noalign{}
\endhead
\bottomrule\noalign{}
\endlastfoot
ID e issue & HU-04 \\
Historia & Como recepcionista del taller, quiero consultar la agenda
digital por día y por bahía, para verificar las solicitudes pendientes
de los clientes y realizar su confirmación manual. \\
Criterio de aceptación 1 & Dado que la recepcionista accede al módulo de
agenda, cuando filtra los registros por el día actual, entonces el
sistema despliega todas las citas con estatus "Pendiente". \\
Criterio de aceptación 2 & Dado que la recepcionista revisa una cita
solicitada, cuando presiona el botón "Confirmar", entonces el estatus
cambia a "Confirmada" y se habilita la asignación de la bahía
correspondiente. \\
Reglas de negocio & RN-08 (Confirmación manual de citas web), RN-05
(Capacidad física del taller). \\
Wireframe & W-05 (4.10, Review 2) \\
Casos de prueba & CPA-32 \\
\end{longtable}
}

\textbf{HU-05: Búsqueda Rápida de Clientes y Vehículos}

{\def\LTcaptype{none} % do not increment counter
\begin{longtable}[]{@{}
  >{\raggedright\arraybackslash}p{(\linewidth - 2\tabcolsep) * \real{0.2778}}
  >{\raggedright\arraybackslash}p{(\linewidth - 2\tabcolsep) * \real{0.7222}}@{}}
\toprule\noalign{}
\begin{minipage}[b]{\linewidth}\raggedright
\textbf{Campo}
\end{minipage} & \begin{minipage}[b]{\linewidth}\raggedright
\textbf{Contenido}
\end{minipage} \\
\midrule\noalign{}
\endhead
\bottomrule\noalign{}
\endlastfoot
ID e issue & HU-05 \\
Historia & Como recepcionista del taller, quiero buscar a un cliente o
vehículo registrado mediante su número telefónico o placas, para
agilizar el proceso de recepción en visitas recurrentes. \\
Criterio de aceptación 1 & Dado que la recepcionista se encuentra en el
módulo de búsqueda, cuando ingresa los dígitos de las placas o el
teléfono, entonces el sistema filtra instantáneamente la base de datos
de clientes recurrentes. \\
Criterio de aceptación 2 & Dado que se localiza el expediente, cuando el
usuario lo selecciona, entonces el sistema precarga los datos generales
del cliente y del auto para evitar duplicidad de registros. \\
Reglas de negocio & Sin regla de negocio específica; ver RF-03.1. \\
Wireframe & Se define en 4.10 (Review 2) \\
Casos de prueba & Se definen en 5.2 (Review 2) \\
\end{longtable}
}

\textbf{HU-06: Registro de Vehículos con Cita Previa}

{\def\LTcaptype{none} % do not increment counter
\begin{longtable}[]{@{}
  >{\raggedright\arraybackslash}p{(\linewidth - 2\tabcolsep) * \real{0.2778}}
  >{\raggedright\arraybackslash}p{(\linewidth - 2\tabcolsep) * \real{0.7222}}@{}}
\toprule\noalign{}
\begin{minipage}[b]{\linewidth}\raggedright
\textbf{Campo}
\end{minipage} & \begin{minipage}[b]{\linewidth}\raggedright
\textbf{Contenido}
\end{minipage} \\
\midrule\noalign{}
\endhead
\bottomrule\noalign{}
\endlastfoot
ID e issue & HU-06 \\
Historia & Como recepcionista del taller, quiero registrar un vehículo
con cita previa seleccionando su folio de origen, para que el sistema
precargue automáticamente los datos del cliente y del auto. \\
Criterio de aceptación 1 & Dado que el cliente llega al taller con cita
agendada, cuando la recepcionista introduce el número de folio de la
cita, entonces el sistema recupera la información asociada. \\
Criterio de aceptación 2 & Dado que los datos están precargados, cuando
se procede con el ingreso, entonces el sistema vincula la cita
directamente con la nueva orden de servicio abierta. \\
Reglas de negocio & RN-08 (Confirmación manual de citas web). \\
Wireframe & W-01 (4.10, Review 2) \\
Casos de prueba & CPA-33 \\
\end{longtable}
}

\textbf{HU-07: Alerta de Bahías Ocupadas}

{\def\LTcaptype{none} % do not increment counter
\begin{longtable}[]{@{}
  >{\raggedright\arraybackslash}p{(\linewidth - 2\tabcolsep) * \real{0.2778}}
  >{\raggedright\arraybackslash}p{(\linewidth - 2\tabcolsep) * \real{0.7222}}@{}}
\toprule\noalign{}
\begin{minipage}[b]{\linewidth}\raggedright
\textbf{Campo}
\end{minipage} & \begin{minipage}[b]{\linewidth}\raggedright
\textbf{Contenido}
\end{minipage} \\
\midrule\noalign{}
\endhead
\bottomrule\noalign{}
\endlastfoot
ID e issue & HU-07 \\
Historia & Como recepcionista del taller, quiero ver la alerta de "Sin
Bahías Disponibles" cuando las 4 bahías físicas estén ocupadas, para
ubicar temporalmente el vehículo en el estacionamiento de recepción. \\
Criterio de aceptación 1 & Dado que las 4 bahías físicas del taller
registran una orden activa, cuando la recepcionista intenta registrar la
entrada de un nuevo vehículo, entonces el sistema emite una advertencia
visual en pantalla indicando "Sin Bahías Disponibles". \\
Criterio de aceptación 2 & Dado que se activa la alerta de saturación,
cuando el usuario acepta la notificación, entonces el sistema permite
guardar el registro temporal asignando la unidad al estacionamiento de
espera. \\
Reglas de negocio & RN-05 (Capacidad física del taller). \\
Wireframe & W-01 (4.10, Review 2) \\
Casos de prueba & CPA-05, CPA-06 \\
\end{longtable}
}

\textbf{HU-08: Asignación de Bahía y Mecánico}

{\def\LTcaptype{none} % do not increment counter
\begin{longtable}[]{@{}
  >{\raggedright\arraybackslash}p{(\linewidth - 2\tabcolsep) * \real{0.2778}}
  >{\raggedright\arraybackslash}p{(\linewidth - 2\tabcolsep) * \real{0.7222}}@{}}
\toprule\noalign{}
\begin{minipage}[b]{\linewidth}\raggedright
\textbf{Campo}
\end{minipage} & \begin{minipage}[b]{\linewidth}\raggedright
\textbf{Contenido}
\end{minipage} \\
\midrule\noalign{}
\endhead
\bottomrule\noalign{}
\endlastfoot
ID e issue & HU-08 \\
Historia & Como jefe de mecánicos, quiero asignar la orden de servicio a
una de las 4 bahías físicas disponibles y a un mecánico responsable,
para iniciar formalmente los trabajos en el taller. \\
Criterio de aceptación 1 & Dado que la orden se encuentra abierta,
cuando el jefe de mecánicos selecciona una de las 4 bahías físicas
libres, entonces el sistema actualiza el estado de la bahía a
"Ocupada". \\
Criterio de aceptación 2 & Dado que las 4 bahías están ocupadas, cuando
se intenta una asignación, entonces el sistema emite una alerta
impidiendo el cruce de unidades en la misma bahía. \\
Reglas de negocio & RN-05 (Capacidad física del taller), RN-09
(Asignación y colaboración de mecánicos). \\
Wireframe & W-02 (4.10, Review 2) \\
Casos de prueba & CPA-07, CPA-08 \\
\end{longtable}
}

\textbf{HU-09: Registro de Diagnóstico en Tableta}

{\def\LTcaptype{none} % do not increment counter
\begin{longtable}[]{@{}
  >{\raggedright\arraybackslash}p{(\linewidth - 2\tabcolsep) * \real{0.2778}}
  >{\raggedright\arraybackslash}p{(\linewidth - 2\tabcolsep) * \real{0.7222}}@{}}
\toprule\noalign{}
\begin{minipage}[b]{\linewidth}\raggedright
\textbf{Campo}
\end{minipage} & \begin{minipage}[b]{\linewidth}\raggedright
\textbf{Contenido}
\end{minipage} \\
\midrule\noalign{}
\endhead
\bottomrule\noalign{}
\endlastfoot
ID e issue & HU-09 \\
Historia & Como jefe de mecánicos, quiero registrar el diagnóstico
técnico y las observaciones iniciales en la tableta, para cambiar el
estado de la orden a "En Diagnóstico". \\
Criterio de aceptación 1 & Dado que el jefe de mecánicos se encuentra
autenticado en la tableta, cuando selecciona una orden con vehículo en
bahía, entonces visualiza un formulario táctil optimizado con campos
amplios para el dictamen técnico. \\
Criterio de aceptación 2 & Dado que se ingresan las observaciones de la
revisión, cuando se presiona "Guardar Diagnóstico", entonces el sistema
actualiza automáticamente el estado de la orden a "En Diagnóstico" y
habilita la creación de cotizaciones. \\
Reglas de negocio & RN-12 (Evidencia fotográfica de recepción).
Usabilidad en tableta: RNF-02 y RNF-03. \\
Wireframe & W-02 (4.10, Review 2) \\
Casos de prueba & CPA-02, CPA-03 \\
\end{longtable}
}

\textbf{HU-10: Solicitud de Refacciones e Insumos}

{\def\LTcaptype{none} % do not increment counter
\begin{longtable}[]{@{}
  >{\raggedright\arraybackslash}p{(\linewidth - 2\tabcolsep) * \real{0.2778}}
  >{\raggedright\arraybackslash}p{(\linewidth - 2\tabcolsep) * \real{0.7222}}@{}}
\toprule\noalign{}
\begin{minipage}[b]{\linewidth}\raggedright
\textbf{Campo}
\end{minipage} & \begin{minipage}[b]{\linewidth}\raggedright
\textbf{Contenido}
\end{minipage} \\
\midrule\noalign{}
\endhead
\bottomrule\noalign{}
\endlastfoot
ID e issue & HU-10 \\
Historia & Como jefe de mecánicos, quiero registrar las refacciones e
insumos que los mecánicos detecten durante la revisión, para que el área
de almacén gestione su cotización o pedido. \\
Criterio de aceptación 1 & Dado que un mecánico reporta piezas
necesarias en la bahía, cuando el jefe de mecánicos las añade desde la
tableta y las vincula a la orden, entonces la solicitud se envía
automáticamente al módulo de almacén. \\
Criterio de aceptación 2 & Dado que el almacén recibe el requerimiento,
cuando se procesa, entonces el sistema refleja el estatus de la
refacción como "Solicitada" o "Pedida". \\
Reglas de negocio & RN-04 (Re-cotización obligatoria por sobrecostos),
RN-06 (Reubicación por espera de refacciones). \\
Wireframe & Se define en 4.10 (Review 2) \\
Casos de prueba & Se definen en 5.2 (Review 2) \\
\end{longtable}
}

\textbf{HU-11: Visualización del Estatus de Bahías}

{\def\LTcaptype{none} % do not increment counter
\begin{longtable}[]{@{}
  >{\raggedright\arraybackslash}p{(\linewidth - 2\tabcolsep) * \real{0.2778}}
  >{\raggedright\arraybackslash}p{(\linewidth - 2\tabcolsep) * \real{0.7222}}@{}}
\toprule\noalign{}
\begin{minipage}[b]{\linewidth}\raggedright
\textbf{Campo}
\end{minipage} & \begin{minipage}[b]{\linewidth}\raggedright
\textbf{Contenido}
\end{minipage} \\
\midrule\noalign{}
\endhead
\bottomrule\noalign{}
\endlastfoot
ID e issue & HU-11 \\
Historia & Como jefe de mecánicos, quiero visualizar el estatus actual
de las 4 bahías de trabajo en una interfaz simplificada para tabletas,
para conocer rápidamente qué unidades están listas o en proceso. \\
Criterio de aceptación 1 & Dado que el jefe de mecánicos accede al
tablero de control, cuando visualiza la sección de bahías, entonces
identifica mediante códigos de color las unidades activas, en
diagnóstico o listas. \\
Criterio de aceptación 2 & Dado que se actualiza el estado de una
reparación, cuando cambia de fase, entonces el tablero refresca la
información en tiempo real para optimizar la supervisión del taller. \\
Reglas de negocio & RN-05 (Capacidad física del taller). Usabilidad en
tableta: RNF-02. \\
Wireframe & W-02 (4.10, Review 2) \\
Casos de prueba & Se definen en 5.2 (Review 2) \\
\end{longtable}
}

\textbf{HU-12: Generación de Presupuesto Desglosado}

{\def\LTcaptype{none} % do not increment counter
\begin{longtable}[]{@{}
  >{\raggedright\arraybackslash}p{(\linewidth - 2\tabcolsep) * \real{0.2778}}
  >{\raggedright\arraybackslash}p{(\linewidth - 2\tabcolsep) * \real{0.7222}}@{}}
\toprule\noalign{}
\begin{minipage}[b]{\linewidth}\raggedright
\textbf{Campo}
\end{minipage} & \begin{minipage}[b]{\linewidth}\raggedright
\textbf{Contenido}
\end{minipage} \\
\midrule\noalign{}
\endhead
\bottomrule\noalign{}
\endlastfoot
ID e issue & HU-12 \\
Historia & Como recepcionista o jefe de mecánicos, quiero generar un
presupuesto desglosado de mano de obra y refacciones con vigencia de 5
días naturales, para enviarlo de forma transparente al cliente. \\
Criterio de aceptación 1 & Dado que el diagnóstico ha concluido, cuando
el personal selecciona la opción de cotizar, entonces el sistema
presenta apartados separados para costos de mano de obra y
refacciones. \\
Criterio de aceptación 2 & Dado que se emite la cotización, entonces el
sistema calcula automáticamente el total con IVA e integra una fecha de
expiración exactamente a los 5 días naturales posteriores a su
emisión. \\
Reglas de negocio & RN-11 (Vigencia de cotizaciones), RN-03 (Selección
granular de rubros). \\
Wireframe & W-03 (4.10, Review 2) \\
Casos de prueba & CPA-10 \\
\end{longtable}
}

\textbf{HU-13: Selección de Servicios Opcionales en Cotización}

{\def\LTcaptype{none} % do not increment counter
\begin{longtable}[]{@{}
  >{\raggedright\arraybackslash}p{(\linewidth - 2\tabcolsep) * \real{0.2778}}
  >{\raggedright\arraybackslash}p{(\linewidth - 2\tabcolsep) * \real{0.7222}}@{}}
\toprule\noalign{}
\begin{minipage}[b]{\linewidth}\raggedright
\textbf{Campo}
\end{minipage} & \begin{minipage}[b]{\linewidth}\raggedright
\textbf{Contenido}
\end{minipage} \\
\midrule\noalign{}
\endhead
\bottomrule\noalign{}
\endlastfoot
ID e issue & HU-13 \\
Historia & Como cliente del taller, quiero seleccionar o deseleccionar
servicios opcionales dentro de mi cotización antes de aprobarla, para
adaptar el presupuesto a mis necesidades financieras. \\
Criterio de aceptación 1 & Dado que el cliente revisa el presupuesto
desglosado en su celular, cuando marca o desmarca las casillas de
servicios opcionales o sugeridos, entonces el sistema recalcula el costo
total en tiempo real. \\
Criterio de aceptación 2 & Dado que se ajustan los opcionales, cuando el
cliente procede a la firma, entonces el sistema guarda únicamente los
conceptos seleccionados como parte oficial de la orden. \\
Reglas de negocio & RN-03 (Selección granular de rubros). \\
Wireframe & W-04 (4.10, Review 2) \\
Casos de prueba & CPA-12 \\
\end{longtable}
}

\textbf{HU-14: Firma Digital y Autorización por 2FA}

{\def\LTcaptype{none} % do not increment counter
\begin{longtable}[]{@{}
  >{\raggedright\arraybackslash}p{(\linewidth - 2\tabcolsep) * \real{0.2778}}
  >{\raggedright\arraybackslash}p{(\linewidth - 2\tabcolsep) * \real{0.7222}}@{}}
\toprule\noalign{}
\begin{minipage}[b]{\linewidth}\raggedright
\textbf{Campo}
\end{minipage} & \begin{minipage}[b]{\linewidth}\raggedright
\textbf{Contenido}
\end{minipage} \\
\midrule\noalign{}
\endhead
\bottomrule\noalign{}
\endlastfoot
ID e issue & HU-14 \\
Historia & Como cliente del taller, quiero realizar la firma digital de
aprobación mediante un código de verificación de 2 pasos enviado a mi
correo, para autorizar los trabajos de forma segura ya sea de forma
presencial o remota. \\
Criterio de aceptación 1 & Dado que el cliente accede al enlace web de
cotización, when solicita el código de autorización, entonces el sistema
envía un token numérico de 6 dígitos al correo electrónico registrado
del cliente. \\
Criterio de aceptación 2 & Dado que el cliente ingresa el código
recibido en la interfaz, cuando presiona "Autorizar Presupuesto",
entonces el sistema valida el token, registra la dirección IP y la marca
de tiempo, cambiando la orden a "Autorizada". \\
Reglas de negocio & RN-02 (Doble verificación en autorizaciones). No
repudio: RNF-04. \\
Wireframe & W-04 (4.10, Review 2) \\
Casos de prueba & CPA-11, CPA-13 a CPA-16 \\
\end{longtable}
}

\textbf{HU-15: Rechazo de Cotización y Cancelación}

{\def\LTcaptype{none} % do not increment counter
\begin{longtable}[]{@{}
  >{\raggedright\arraybackslash}p{(\linewidth - 2\tabcolsep) * \real{0.2778}}
  >{\raggedright\arraybackslash}p{(\linewidth - 2\tabcolsep) * \real{0.7222}}@{}}
\toprule\noalign{}
\begin{minipage}[b]{\linewidth}\raggedright
\textbf{Campo}
\end{minipage} & \begin{minipage}[b]{\linewidth}\raggedright
\textbf{Contenido}
\end{minipage} \\
\midrule\noalign{}
\endhead
\bottomrule\noalign{}
\endlastfoot
ID e issue & HU-15 \\
Historia & Como cliente del taller, quiero presionar un botón de rechazo
en la cotización digital, para cancelar la orden y coordinar la
recolección del vehículo sin autorizar reparaciones. \\
Criterio de aceptación 1 & Dado que el cliente evalúa la cotización en
su dispositivo, cuando presiona el botón "Rechazar", entonces el sistema
solicita una confirmación de la acción. \\
Criterio de aceptación 2 & Dado que se confirma el rechazo, cuando se
procesa, entonces la orden cambia a estatus "Rechazada" y notifica a
recepción para coordinar la entrega física del auto sin costo de
reparación. \\
Reglas de negocio & RN-02 (Doble verificación en autorizaciones). Flujo
de rechazo: CU-11 (A1). \\
Wireframe & W-04 (4.10, Review 2) \\
Casos de prueba & CPA-17 \\
\end{longtable}
}

\textbf{HU-16: Re-cotización por Incremento \textgreater10\%}

{\def\LTcaptype{none} % do not increment counter
\begin{longtable}[]{@{}
  >{\raggedright\arraybackslash}p{(\linewidth - 2\tabcolsep) * \real{0.2778}}
  >{\raggedright\arraybackslash}p{(\linewidth - 2\tabcolsep) * \real{0.7222}}@{}}
\toprule\noalign{}
\begin{minipage}[b]{\linewidth}\raggedright
\textbf{Campo}
\end{minipage} & \begin{minipage}[b]{\linewidth}\raggedright
\textbf{Contenido}
\end{minipage} \\
\midrule\noalign{}
\endhead
\bottomrule\noalign{}
\endlastfoot
ID e issue & HU-16 \\
Historia & Como cliente del taller, quiero que se me pida una nueva
autorización con código 2FA cuando los costos imprevistos superen el 10
\% del presupuesto autorizado, para no pagar trabajos que no aprobé. \\
Criterio de aceptación 1 & Dado que surgen refacciones imprevistas,
cuando el incremento total supera el 10\% del presupuesto aprobado,
entonces el sistema bloquea el flujo y solicita una nueva autorización
con código 2FA. \\
Criterio de aceptación 2 & Dado que el cliente recibe la notificación
del sobrecosto, cuando aprueba mediante el código digital, entonces se
actualiza el monto total e historial de la orden. \\
Reglas de negocio & RN-04 (Re-cotización obligatoria por sobrecostos),
RN-02 (Doble verificación en autorizaciones). \\
Wireframe & W-03 (4.10, Review 2) \\
Casos de prueba & CPA-18, CPA-19 \\
\end{longtable}
}

\textbf{HU-17: Directorio y Gestión de Proveedores Locales}

{\def\LTcaptype{none} % do not increment counter
\begin{longtable}[]{@{}
  >{\raggedright\arraybackslash}p{(\linewidth - 2\tabcolsep) * \real{0.2778}}
  >{\raggedright\arraybackslash}p{(\linewidth - 2\tabcolsep) * \real{0.7222}}@{}}
\toprule\noalign{}
\begin{minipage}[b]{\linewidth}\raggedright
\textbf{Campo}
\end{minipage} & \begin{minipage}[b]{\linewidth}\raggedright
\textbf{Contenido}
\end{minipage} \\
\midrule\noalign{}
\endhead
\bottomrule\noalign{}
\endlastfoot
ID e issue & HU-17 \\
Historia & Como encargado de almacén, quiero consultar y administrar un
directorio de proveedores locales con teléfonos y correos electrónicos,
para solicitar refacciones especiales rápidamente. \\
Criterio de aceptación 1 & Dado que el almacenista ingresa al módulo de
refacciones, cuando abre el directorio de proveedores, entonces
visualiza la lista con datos de contacto actualizados (teléfonos y
correos). \\
Criterio de aceptación 2 & Dado que se requiere dar de alta un nuevo
proveedor, cuando el almacenista completa el formulario de registro,
entonces el sistema lo almacena para futuras solicitudes de piezas. \\
Reglas de negocio & RN-06 (Reubicación por espera de refacciones). \\
Wireframe & Se define en 4.10 (Review 2) \\
Casos de prueba & Se definen en 5.2 (Review 2) \\
\end{longtable}
}

\textbf{HU-18: Registro de Costos e Historial de Compras}

{\def\LTcaptype{none} % do not increment counter
\begin{longtable}[]{@{}
  >{\raggedright\arraybackslash}p{(\linewidth - 2\tabcolsep) * \real{0.2778}}
  >{\raggedright\arraybackslash}p{(\linewidth - 2\tabcolsep) * \real{0.7222}}@{}}
\toprule\noalign{}
\begin{minipage}[b]{\linewidth}\raggedright
\textbf{Campo}
\end{minipage} & \begin{minipage}[b]{\linewidth}\raggedright
\textbf{Contenido}
\end{minipage} \\
\midrule\noalign{}
\endhead
\bottomrule\noalign{}
\endlastfoot
ID e issue & HU-18 \\
Historia & Como encargado de almacén, quiero registrar los costos,
fechas de llegada y tickets de compra de las refacciones asociadas a una
orden, para mantener un control financiero estricto. \\
Criterio de aceptación 1 & Dado que se adquiere una pieza externa con un
proveedor, cuando el almacenista registra el ticket, costo unitario y
fecha estimada de entrega, entonces el sistema lo vincula a la orden de
servicio. \\
Criterio de aceptación 2 & Dado que se ingresa la compra, cuando se
consolida el registro, entonces el sistema actualiza el historial de
adquisiciones de la refacción para auditoría interna. \\
Reglas de negocio & Sin regla de negocio específica; ver RF-08.2 y
RF-08.3. \\
Wireframe & Se define en 4.10 (Review 2) \\
Casos de prueba & Se definen en 5.2 (Review 2) \\
\end{longtable}
}

\textbf{HU-19: Verificación de Stock Interno}

{\def\LTcaptype{none} % do not increment counter
\begin{longtable}[]{@{}
  >{\raggedright\arraybackslash}p{(\linewidth - 2\tabcolsep) * \real{0.2778}}
  >{\raggedright\arraybackslash}p{(\linewidth - 2\tabcolsep) * \real{0.7222}}@{}}
\toprule\noalign{}
\begin{minipage}[b]{\linewidth}\raggedright
\textbf{Campo}
\end{minipage} & \begin{minipage}[b]{\linewidth}\raggedright
\textbf{Contenido}
\end{minipage} \\
\midrule\noalign{}
\endhead
\bottomrule\noalign{}
\endlastfoot
ID e issue & HU-19 \\
Historia & Como jefe de mecánicos, quiero verificar el stock de
refacciones comunes en el inventario interno, para saber si se dispone
de la pieza de forma inmediata. \\
Criterio de aceptación 1 & Dado que el jefe de mecánicos busca un insumo
estándar, cuando consulta el inventario en el sistema, entonces
visualiza las existencias disponibles en almacén. \\
Criterio de aceptación 2 & Dado que hay existencia en stock, cuando se
selecciona la pieza, entonces el sistema descuenta automáticamente la
cantidad del inventario interno. \\
Reglas de negocio & Sin regla de negocio específica; ver RF-08.1. \\
Wireframe & Se define en 4.10 (Review 2) \\
Casos de prueba & Se definen en 5.2 (Review 2) \\
\end{longtable}
}

\textbf{HU-20: Notificación Automática de Fin de Trabajos}

{\def\LTcaptype{none} % do not increment counter
\begin{longtable}[]{@{}
  >{\raggedright\arraybackslash}p{(\linewidth - 2\tabcolsep) * \real{0.2778}}
  >{\raggedright\arraybackslash}p{(\linewidth - 2\tabcolsep) * \real{0.7222}}@{}}
\toprule\noalign{}
\begin{minipage}[b]{\linewidth}\raggedright
\textbf{Campo}
\end{minipage} & \begin{minipage}[b]{\linewidth}\raggedright
\textbf{Contenido}
\end{minipage} \\
\midrule\noalign{}
\endhead
\bottomrule\noalign{}
\endlastfoot
ID e issue & HU-20 \\
Historia & Como recepcionista, quiero recibir una notificación
automática cuando los trabajos del vehículo finalicen, para avisar al
cliente que su unidad está lista para entrega. \\
Criterio de aceptación 1 & Dado que el mecánico marca la orden como
terminada en bahía, cuando se actualiza el estatus, entonces el sistema
envía una alerta interna a la recepcionista. \\
Criterio de aceptación 2 & Dado que la recepción recibe la alerta,
cuando se verifica el estado, entonces se habilita el aviso automatizado
al cliente de que su vehículo está listo. \\
Reglas de negocio & Sin regla de negocio específica; ver CU-07. \\
Wireframe & W-06 (4.10, Review 2) \\
Casos de prueba & Se definen en 5.2 (Review 2) \\
\end{longtable}
}

\textbf{HU-21: Bloqueo de Entrega por Saldo Pendiente}

{\def\LTcaptype{none} % do not increment counter
\begin{longtable}[]{@{}
  >{\raggedright\arraybackslash}p{(\linewidth - 2\tabcolsep) * \real{0.2778}}
  >{\raggedright\arraybackslash}p{(\linewidth - 2\tabcolsep) * \real{0.7222}}@{}}
\toprule\noalign{}
\begin{minipage}[b]{\linewidth}\raggedright
\textbf{Campo}
\end{minipage} & \begin{minipage}[b]{\linewidth}\raggedright
\textbf{Contenido}
\end{minipage} \\
\midrule\noalign{}
\endhead
\bottomrule\noalign{}
\endlastfoot
ID e issue & HU-21 \\
Historia & Como gerente del taller, quiero que el sistema impida
registrar la entrega de un vehículo cuya orden tenga saldo pendiente,
para evitar pérdidas financieras al taller. \\
Criterio de aceptación 1 & Dado que un vehículo se encuentra en estado
"Lista para Entrega", cuando la recepcionista intenta registrar la
salida y el balance de pago es mayor a cero, entonces el sistema bloquea
la acción de liberación. \\
Criterio de aceptación 2 & Dado que se activa el bloqueo por saldo
pendiente, cuando el sistema detecta que cuenta con un convenio de
crédito autorizado, entonces permite omitir la restricción
exclusivamente bajo permisos del rol Gerente. \\
Reglas de negocio & RN-10 (Restricción de entrega y excepción a
flotillas). \\
Wireframe & W-06 (4.10, Review 2) \\
Casos de prueba & CPA-21 \\
\end{longtable}
}

\textbf{HU-22: Autorización Excepcional de Crédito a Flotillas}

{\def\LTcaptype{none} % do not increment counter
\begin{longtable}[]{@{}
  >{\raggedright\arraybackslash}p{(\linewidth - 2\tabcolsep) * \real{0.2778}}
  >{\raggedright\arraybackslash}p{(\linewidth - 2\tabcolsep) * \real{0.7222}}@{}}
\toprule\noalign{}
\begin{minipage}[b]{\linewidth}\raggedright
\textbf{Campo}
\end{minipage} & \begin{minipage}[b]{\linewidth}\raggedright
\textbf{Contenido}
\end{minipage} \\
\midrule\noalign{}
\endhead
\bottomrule\noalign{}
\endlastfoot
ID e issue & HU-22 \\
Historia & Como gerente del taller, quiero autorizar de forma
excepcional la entrega a crédito únicamente para vehículos de flotillas
empresariales registradas, para cumplir con los convenios
comerciales. \\
Criterio de aceptación 1 & Dado que un vehículo pertenece a una empresa
con cuenta de flotilla registrada, cuando se intenta liberar sin pago en
caja, entonces el sistema solicita la clave de autorización
gerencial. \\
Criterio de aceptación 2 & Dado que el gerente ingresa su contraseña de
aprobación, cuando valida el convenio comercial, entonces el sistema
permite liberar la unidad y factura el adeudo a crédito empresarial. \\
Reglas de negocio & RN-10 (Restricción de entrega y excepción a
flotillas). \\
Wireframe & W-06 (4.10, Review 2) \\
Casos de prueba & CPA-24, CPA-25 \\
\end{longtable}
}

\textbf{HU-23: Registro de Pago Total y Cierre de Orden}

{\def\LTcaptype{none} % do not increment counter
\begin{longtable}[]{@{}
  >{\raggedright\arraybackslash}p{(\linewidth - 2\tabcolsep) * \real{0.2778}}
  >{\raggedright\arraybackslash}p{(\linewidth - 2\tabcolsep) * \real{0.7222}}@{}}
\toprule\noalign{}
\begin{minipage}[b]{\linewidth}\raggedright
\textbf{Campo}
\end{minipage} & \begin{minipage}[b]{\linewidth}\raggedright
\textbf{Contenido}
\end{minipage} \\
\midrule\noalign{}
\endhead
\bottomrule\noalign{}
\endlastfoot
ID e issue & HU-23 \\
Historia & Como recepcionista, quiero registrar el pago total del
servicio y cerrar formalmente la orden, para liberar la bahía de trabajo
y archivar el historial del vehículo. \\
Criterio de aceptación 1 & Dado que el cliente cubre el monto total en
caja, cuando la recepcionista registra el método de pago y confirma la
transacción, entonces el sistema salda la cuenta a cero. \\
Criterio de aceptación 2 & Dado que el pago ha sido registrado con
éxito, cuando se ejecuta el cierre, entonces la orden cambia a estado
"Cerrada", liberando de inmediato la bahía física ocupada y generando el
comprobante digital. \\
Reglas de negocio & RN-10 (Restricción de entrega y excepción a
flotillas), RN-13 (Garantía de mano de obra), RN-01 (Expiración y
vigencia de enlaces). \\
Wireframe & W-06 (4.10, Review 2) \\
Casos de prueba & CPA-22, CPA-23 \\
\end{longtable}
}

\textbf{HU-24: Consulta de Vigencia de Garantía}

{\def\LTcaptype{none} % do not increment counter
\begin{longtable}[]{@{}
  >{\raggedright\arraybackslash}p{(\linewidth - 2\tabcolsep) * \real{0.2778}}
  >{\raggedright\arraybackslash}p{(\linewidth - 2\tabcolsep) * \real{0.7222}}@{}}
\toprule\noalign{}
\begin{minipage}[b]{\linewidth}\raggedright
\textbf{Campo}
\end{minipage} & \begin{minipage}[b]{\linewidth}\raggedright
\textbf{Contenido}
\end{minipage} \\
\midrule\noalign{}
\endhead
\bottomrule\noalign{}
\endlastfoot
ID e issue & HU-24 \\
Historia & Como cliente, quiero consultar la vigencia de la garantía de
30 días en mano de obra asociada a mi servicio anterior, para reclamar
ajustes en caso de incidencias. \\
Criterio de aceptación 1 & Dado que el cliente revisa el expediente de
su servicio pasado en el portal web, cuando consulta el apartado de
garantías, entonces visualiza la fecha exacta de expiración (30 días
naturales en mano de obra). \\
Criterio de aceptación 2 & Dado que la garantía se encuentra vigente,
cuando se solicita una revisión por el mismo concepto, entonces el
sistema valida el estatus sin generar cargos adicionales de mano de
obra. \\
Reglas de negocio & RN-13 (Garantía de mano de obra). \\
Wireframe & Se define en 4.10 (Review 2) \\
Casos de prueba & Se definen en 5.2 (Review 2) \\
\end{longtable}
}

\textbf{HU-25: Operación Offline del Personal Interno}

{\def\LTcaptype{none} % do not increment counter
\begin{longtable}[]{@{}
  >{\raggedright\arraybackslash}p{(\linewidth - 2\tabcolsep) * \real{0.2778}}
  >{\raggedright\arraybackslash}p{(\linewidth - 2\tabcolsep) * \real{0.7222}}@{}}
\toprule\noalign{}
\begin{minipage}[b]{\linewidth}\raggedright
\textbf{Campo}
\end{minipage} & \begin{minipage}[b]{\linewidth}\raggedright
\textbf{Contenido}
\end{minipage} \\
\midrule\noalign{}
\endhead
\bottomrule\noalign{}
\endlastfoot
ID e issue & HU-25 \\
Historia & Como recepcionista o jefe de mecánicos, quiero registrar la
llegada de autos, tomar fotografías y hacer diagnósticos sin conexión a
internet usando la batería del equipo, para que la operación no se
detenga ante fallas eléctricas. \\
Criterio de aceptación 1 & Dado que ocurre un corte de energía o caída
de internet en el taller, cuando el personal opera mediante las baterías
de laptops o tabletas, entonces el sistema local permite continuar
capturando registros, checklists y fotografías. \\
Criterio de aceptación 2 & Dado que se trabaja de forma local (offline),
cuando se guardan los datos, entonces el almacenamiento temporal
mediante Service Workers e IndexedDB asegura que no exista pérdida de
información. \\
Reglas de negocio & Sin regla de negocio específica; ver RNF-06 y
RNF-08. \\
Wireframe & Se define en 4.10 (Review 2) \\
Casos de prueba & Se definen en 5.2 (Review 2) \\
\end{longtable}
}

\textbf{HU-26: Sincronización Automática en la Nube}

{\def\LTcaptype{none} % do not increment counter
\begin{longtable}[]{@{}
  >{\raggedright\arraybackslash}p{(\linewidth - 2\tabcolsep) * \real{0.2778}}
  >{\raggedright\arraybackslash}p{(\linewidth - 2\tabcolsep) * \real{0.7222}}@{}}
\toprule\noalign{}
\begin{minipage}[b]{\linewidth}\raggedright
\textbf{Campo}
\end{minipage} & \begin{minipage}[b]{\linewidth}\raggedright
\textbf{Contenido}
\end{minipage} \\
\midrule\noalign{}
\endhead
\bottomrule\noalign{}
\endlastfoot
ID e issue & HU-26 \\
Historia & Como recepcionista o jefe de mecánicos, quiero que los
registros y fotografías que capturé sin conexión se sincronicen
automáticamente cuando vuelva la red, para no perder ni duplicar
información. \\
Criterio de aceptación 1 & Dado que se restablece la conexión a internet
en el taller, cuando el sistema detecta conectividad, entonces inicia de
forma automática la sincronización de los datos almacenados
localmente. \\
Criterio de aceptación 2 & Dado que finaliza el proceso de
sincronización, cuando se actualiza la base de datos central en la nube,
entonces la información multimedia y de auditoría se consolida sin
duplicidades. \\
Reglas de negocio & Sin regla de negocio específica; ver RNF-09. \\
Wireframe & Se define en 4.10 (Review 2) \\
Casos de prueba & Se definen en 5.2 (Review 2) \\
\end{longtable}
}

\textbf{HU-27: Gestión de Roles y Permisos de Seguridad}

{\def\LTcaptype{none} % do not increment counter
\begin{longtable}[]{@{}
  >{\raggedright\arraybackslash}p{(\linewidth - 2\tabcolsep) * \real{0.2778}}
  >{\raggedright\arraybackslash}p{(\linewidth - 2\tabcolsep) * \real{0.7222}}@{}}
\toprule\noalign{}
\begin{minipage}[b]{\linewidth}\raggedright
\textbf{Campo}
\end{minipage} & \begin{minipage}[b]{\linewidth}\raggedright
\textbf{Contenido}
\end{minipage} \\
\midrule\noalign{}
\endhead
\bottomrule\noalign{}
\endlastfoot
ID e issue & HU-27 \\
Historia & Como gerente del taller, quiero dar de alta al personal y
asignarle su rol (Gerente, Recepcionista, Jefe de Mecánicos o Encargado
de Almacén), para que cada quien acceda solo a la información que le
corresponde. \\
Criterio de aceptación 1 & Dado que un usuario intenta acceder a la
plataforma, cuando introduce sus credenciales, entonces el sistema
autentica su identidad y restringe los menús según su rol asignado
(Gerente, Recepcionista, Jefe de Mecánicos o Encargado de Almacén). \\
Criterio de aceptación 2 & Dado que un usuario sin privilegios intenta
ingresar a una ruta protegida (como reportes gerenciales), entonces el
sistema deniega el acceso y muestra una pantalla de error por permisos
insuficientes. \\
Reglas de negocio & Sin regla de negocio específica; ver RNF-05 y 4.8
(control de acceso por rol). \\
Wireframe & Se define en 4.10 (Review 2) \\
Casos de prueba & CPA-28, CPA-29 \\
\end{longtable}
}

\textbf{HU-28: Accesibilidad PWA Responsiva para Clientes}

{\def\LTcaptype{none} % do not increment counter
\begin{longtable}[]{@{}
  >{\raggedright\arraybackslash}p{(\linewidth - 2\tabcolsep) * \real{0.2778}}
  >{\raggedright\arraybackslash}p{(\linewidth - 2\tabcolsep) * \real{0.7222}}@{}}
\toprule\noalign{}
\begin{minipage}[b]{\linewidth}\raggedright
\textbf{Campo}
\end{minipage} & \begin{minipage}[b]{\linewidth}\raggedright
\textbf{Contenido}
\end{minipage} \\
\midrule\noalign{}
\endhead
\bottomrule\noalign{}
\endlastfoot
ID e issue & HU-28 \\
Historia & Como cliente, quiero visualizar la plataforma web desde mi
teléfono celular de forma responsiva y sin instalar aplicaciones
adicionales, para consultar mi estatus con facilidad. \\
Criterio de aceptación 1 & Dado que el cliente abre el enlace recibido
por SMS o correo en un navegador móvil (iOS/Android), cuando carga la
página, entonces el diseño se adapta responsivamente a la pantalla sin
requerir descargas de tiendas de aplicaciones. \\
Criterio de aceptación 2 & Dado que el cliente navega por el portal web,
cuando consulta el progreso de su vehículo o galería de fotos, entonces
los elementos táctiles responden de manera fluida y adaptada a
dispositivos móviles. \\
Reglas de negocio & RN-01 (Expiración y vigencia de enlaces).
Portabilidad: RNF-01. \\
Wireframe & W-04 (4.10, Review 2) \\
Casos de prueba & Se definen en 5.2 (Review 2) \\
\end{longtable}
}

\textbf{HU-29: Generación de Reportes Gerenciales}

{\def\LTcaptype{none} % do not increment counter
\begin{longtable}[]{@{}
  >{\raggedright\arraybackslash}p{(\linewidth - 2\tabcolsep) * \real{0.2778}}
  >{\raggedright\arraybackslash}p{(\linewidth - 2\tabcolsep) * \real{0.7222}}@{}}
\toprule\noalign{}
\begin{minipage}[b]{\linewidth}\raggedright
\textbf{Campo}
\end{minipage} & \begin{minipage}[b]{\linewidth}\raggedright
\textbf{Contenido}
\end{minipage} \\
\midrule\noalign{}
\endhead
\bottomrule\noalign{}
\endlastfoot
ID e issue & HU-29 \\
Historia & Como gerente, quiero generar reportes mensuales de ingresos,
servicios realizados y tiempos de ocupación de las 4 bahías, para
evaluar el rendimiento general del taller. \\
Criterio de aceptación 1 & Dado que el gerente ingresa al módulo de
reportes, cuando selecciona el intervalo de fechas mensual, entonces el
sistema genera métricas consolidadas de ingresos y servicios
completados. \\
Criterio de aceptación 2 & Dado que se consultan los tiempos operativos,
cuando el sistema procesa las órdenes, entonces muestra gráficos de
ocupación de las 4 bahías para medir la eficiencia del taller. \\
Reglas de negocio & Sin regla de negocio específica; ver RF-11. \\
Wireframe & Se define en 4.10 (Review 2) \\
Casos de prueba & Se definen en 5.2 (Review 2) \\
\end{longtable}
}

\textbf{HU-30: Cumplimiento de Privacidad y Uso de Datos Ficticios}

{\def\LTcaptype{none} % do not increment counter
\begin{longtable}[]{@{}
  >{\raggedright\arraybackslash}p{(\linewidth - 2\tabcolsep) * \real{0.2778}}
  >{\raggedright\arraybackslash}p{(\linewidth - 2\tabcolsep) * \real{0.7222}}@{}}
\toprule\noalign{}
\begin{minipage}[b]{\linewidth}\raggedright
\textbf{Campo}
\end{minipage} & \begin{minipage}[b]{\linewidth}\raggedright
\textbf{Contenido}
\end{minipage} \\
\midrule\noalign{}
\endhead
\bottomrule\noalign{}
\endlastfoot
ID e issue & HU-30 \\
Historia & Como cliente del taller, quiero conocer el aviso de
privacidad antes de proporcionar mis datos, para saber para qué se
usarán conforme a la LFPDPPP. \\
Criterio de aceptación 1 & Dado que la recepcionista registra a un
cliente nuevo, cuando captura sus datos personales, entonces el sistema
muestra el aviso de privacidad simplificado (4.8) antes de guardar el
registro. \\
Criterio de aceptación 2 & Dado que el cliente abre su enlace seguro por
primera vez, cuando carga la página, entonces el sistema muestra el
enlace al aviso de privacidad integral y cómo ejercer sus derechos
ARCO. \\
Reglas de negocio & Sin regla de negocio específica; ver 4.8 y
RNF-05. \\
Wireframe & Se define en 4.10 (Review 2) \\
Casos de prueba & CPA-31 \\
\end{longtable}
}

\figura{figuras/story-map.png}{width=\linewidth}{4.1}{Story map con el corte del release v0.1.0: actividades del flujo de atención en columnas, historias ordenadas por prioridad y, bajo la línea, las que pasan a la Fase 2. Elaboración propia}{docs/04-requisitos/story-map.py}

\subsection{4.6 Requisitos no
funcionales}\label{requisitos-no-funcionales}

{\def\LTcaptype{none} % do not increment counter
\begin{longtable}[]{@{}
  >{\raggedright\arraybackslash}p{(\linewidth - 8\tabcolsep) * \real{0.0962}}
  >{\raggedright\arraybackslash}p{(\linewidth - 8\tabcolsep) * \real{0.2136}}
  >{\raggedright\arraybackslash}p{(\linewidth - 8\tabcolsep) * \real{0.2735}}
  >{\raggedright\arraybackslash}p{(\linewidth - 8\tabcolsep) * \real{0.2137}}
  >{\raggedright\arraybackslash}p{(\linewidth - 8\tabcolsep) * \real{0.2031}}@{}}
\toprule\noalign{}
\begin{minipage}[b]{\linewidth}\raggedright
\textbf{ID}
\end{minipage} & \begin{minipage}[b]{\linewidth}\raggedright
\textbf{Característica ISO/IEC 25010:2023}
\end{minipage} & \begin{minipage}[b]{\linewidth}\raggedright
\textbf{Requisito}
\end{minipage} & \begin{minipage}[b]{\linewidth}\raggedright
\textbf{Métrica y umbral}
\end{minipage} & \begin{minipage}[b]{\linewidth}\raggedright
\textbf{Cómo se verificará}
\end{minipage} \\
\midrule\noalign{}
\endhead
\bottomrule\noalign{}
\endlastfoot
RNF-01 & Flexibilidad (instalabilidad y adaptabilidad) & El portal del
cliente deberá ser una PWA responsiva que funcione en navegadores
móviles iOS y Android sin instalar software ni conceder permisos
especiales. & Funciona sin errores bloqueantes en las 2 últimas
versiones de Safari (iOS) y Chrome (Android); 0 permisos solicitados al
cliente; sin desplazamiento horizontal en anchos de 360 a 430 px. &
Pruebas manuales en 2 móviles reales (iOS y Android) con una lista de
verificación de las pantallas del cliente y revisión con las
herramientas de desarrollo del navegador. \\
RNF-02 & Capacidad de interacción & La interfaz de la tableta del jefe
de Mecánicos deberá tener botones grandes, alto contraste y controles
táctiles adaptados al uso con manos sucias y en ambiente ruidoso. &
Botones de al menos 48 × 48 px con separación mínima de 8 px; contraste
de texto de al menos 7:1; al menos 95 \% de aciertos en 20 toques con
guantes de trabajo. & Auditoría con Lighthouse o axe DevTools
(gratuitos), medición del contraste con un verificador y prueba en la
tableta real con guantes. \\
RNF-03 & Capacidad de interacción & La interfaz de la tableta deberá
ofrecer flujos de pocos pasos para las operaciones frecuentes del
taller. & Asignar bahía, actualizar el estado de una orden y sellar un
trabajo en máximo 4 toques y 30 s cada una. & Prueba de usabilidad
cronometrada con al menos 3 usuarios (o compañeros que simulen el rol),
repitiendo cada operación 3 veces. \\
RNF-04 & Seguridad (no repudio e integridad) & Toda autorización de
presupuesto o de sobrecosto mayor al 10 \% deberá registrar marca de
tiempo y metadatos inalterables, incluida la dirección IP, para evitar
disputas con el cliente. & 100 \% de las autorizaciones guardan fecha y
hora UTC, IP, identificador de orden y huella SHA-256 del documento; 100
\% de los intentos de alterar un registro se detectan al verificar la
huella. & Prueba con 20 autorizaciones, revisión de los registros
guardados e intento de modificar directamente la base de datos para
confirmar que la verificación lo detecta. \\
RNF-05 & Seguridad (confidencialidad y cumplimiento legal) & El manejo
de datos personales deberá apegarse a la Ley Federal de Protección de
Datos Personales en Posesión de los Particulares y usar solo datos
ficticios en las pruebas. & 100 \% de las conexiones con HTTPS (TLS 1.2
o superior); datos personales sensibles cifrados en reposo; aviso de
privacidad visible en el 100 \% de los formularios que capturan datos; 0
registros reales en el ambiente de pruebas. & Análisis con testssl.sh o
SSL Labs, lista de verificación de la ley revisada por el equipo y
revisión de los datos de prueba antes de cada entrega. \\
RNF-06 & Eficiencia de desempeño & El registro de recepción se deberá
guardar de forma local con rapidez, aun sin conexión. & Guardar un
registro con 5 fotografías en máximo 2 s (percentil 95) en la tableta o
laptop del taller. & 20 repeticiones cronometradas con la herramienta de
rendimiento del navegador, calculando el percentil 95. \\
RNF-07 & Mantenibilidad (modificabilidad) & El desarrollo deberá usar
únicamente herramientas de código abierto, frameworks de uso libre (como
FastAPI y Python) y planes educativos. & 100 \% de las dependencias con
licencia de código abierto o plan educativo vigente; costo de
licenciamiento de \$0 en esta fase. & Inventario de dependencias con
herramientas de revisión de licencias (por ejemplo, pip-licenses) y
registro de los planes educativos usados. \\
RNF-08 & Fiabilidad (tolerancia a fallos) & El personal interno deberá
poder operar sin internet o sin energía eléctrica, con equipos en
batería, bajo una arquitectura local con sincronización diferida. & Las
7 operaciones definidas (registro de entrada, checklist, fotografías,
diagnósticos, asignación de bahías, consulta de proveedores y
autorizaciones gerenciales) funcionan al 100 \% durante al menos 4 h sin
conexión; 0 datos perdidos. & Casos de prueba de las 7 operaciones con
la red desconectada y un equipo con batería, comparando los registros al
terminar. \\
RNF-09 & Fiabilidad (recuperabilidad) & Al restablecerse la conexión, el
sistema deberá sincronizar automáticamente datos, archivos multimedia y
registros de auditoría con la base de datos central, sin duplicados. &
100 \% de los registros sincronizados, 0 duplicados y 0 pérdidas; 200
registros con 50 fotografías en máximo 5 min. & Crear 200 registros sin
conexión, reconectar y comparar conteos y huellas con la base central;
repetir 3 veces e interrumpir la conexión a mitad de una
sincronización. \\
RNF-10 & Fiabilidad (tolerancia a fallos del cliente) & Ante una falla
de conectividad durante la consulta de estatus o la autorización 2FA, la
interfaz del cliente deberá informar el estado y reintentar
automáticamente. & Mensaje de estado en máximo 3 s; reintento automático
cada 30 s o menos hasta completar; 0 autorizaciones duplicadas o
perdidas. & Cortar el servidor durante la consulta y durante el envío
del código, cronometrar el mensaje y revisar la base de datos tras
restablecer el servicio. \\
\end{longtable}
}

\subsection{4.7 Reglas de negocio y requisitos de
dominio}\label{reglas-de-negocio-y-requisitos-de-dominio}

{\def\LTcaptype{none} % do not increment counter
\begin{longtable}[]{@{}
  >{\raggedright\arraybackslash}p{(\linewidth - 6\tabcolsep) * \real{0.0963}}
  >{\raggedright\arraybackslash}p{(\linewidth - 6\tabcolsep) * \real{0.4871}}
  >{\raggedright\arraybackslash}p{(\linewidth - 6\tabcolsep) * \real{0.2351}}
  >{\raggedright\arraybackslash}p{(\linewidth - 6\tabcolsep) * \real{0.1815}}@{}}
\toprule\noalign{}
\begin{minipage}[b]{\linewidth}\raggedright
\textbf{ID}
\end{minipage} & \begin{minipage}[b]{\linewidth}\raggedright
\textbf{Regla de negocio}
\end{minipage} & \begin{minipage}[b]{\linewidth}\raggedright
\textbf{Fuente (anteproyecto / entrevista con fecha)}
\end{minipage} & \begin{minipage}[b]{\linewidth}\raggedright
\textbf{Requisitos que afecta}
\end{minipage} \\
\midrule\noalign{}
\endhead
\bottomrule\noalign{}
\endlastfoot
RN-01 & Expiración y vigencia de enlaces: El enlace web de seguimiento
permanecerá activo durante la estancia del vehículo en el taller.
Expirará automáticamente al entregar y saldar la unidad. Si el cliente
requiere ver su historial posteriormente, deberá solicitar la generación
de un nuevo enlace. La vigencia del enlace es independiente de la
vigencia de la cotización (RN-11). & Entrevista con el cliente (19/10/2026), minuta A.2, E-03 & RF-02.1, RF-02.2 \\
RN-02 & Doble verificación en autorizaciones: Toda aprobación remota o
presencial de cotización exigirá una doble verificación (código 2FA de 6
dígitos, de un solo uso y vigente 5 minutos, enviado al correo
electrónico registrado del cliente) para evitar reclamos por trabajos no
autorizados. & Entrevista con el cliente (19/10/2026), minuta A.2, E-01 & RF-06, RNF-04 \\
RN-03 & Selección granular de rubros: El cliente puede aprobar o
rechazar de manera independiente los conceptos opcionales desglosados en
el presupuesto (mano de obra o refacciones), ejecutando únicamente lo
autorizado explícitamente. & Entrevista con el cliente (19/10/2026), minuta A.2, E-02 &
RF-06 \\
RN-04 & Re-cotización obligatoria por sobrecostos (\textgreater10\%):
Cualquier trabajo o repuesto imprevisto durante la reparación que
incremente el costo inicial más de un 10\% detendrá el avance hasta
contar con una nueva autorización validada por el cliente. &
Anteproyecto y propuesta técnica & RF-07, RNF-04 \\
RN-05 & Capacidad física del taller\textbf{:} La atención simultánea
está limitada estrictamente a 4 bahías físicas (1 vehículo por bahía) y
un área de estacionamiento de recepción/espera con capacidad máxima para
5 vehículos. & Anteproyecto y entrevista con el cliente (19/10/2026), minuta A.2, E-04 & RF-05,
RF-09 \\
RN-06 & Reubicación por espera de refacciones: Si un vehículo requiere
refacciones externas con demora de entrega de varios días, se liberará
la bahía de trabajo reubicando la unidad en el estacionamiento hasta el
arribo de las piezas. & Entrevista con el cliente (19/10/2026), minuta A.2, E-05 & RF-05,
RF-08.2 \\
RN-07 & Clasificación de vehículos autorizados: El sistema operará
exclusivamente con vehículos automotrices ligeros (automóviles
particulares, motocicletas y camionetas hasta de 3 toneladas). Se
excluyen bicicletas, tráilers y maquinaria pesada. & Entrevista con el cliente (19/10/2026), minuta A.2, E-06 & RF-03.1 \\
RN-08 & Confirmación manual de citas web: Las solicitudes de cita
agendadas desde el portal público ingresarán en estado "Pendiente" y no
reservarán espacio en la agenda hasta que la recepcionista las confirme
manualmente. & Entrevista con el cliente (19/10/2026), minuta A.2, E-07 & RF-09 \\
RN-09 & Asignación y colaboración de mecánicos\textbf{:} Cada orden
mantendrá un mecánico principal responsable. Otros mecánicos podrán
colaborar en el trabajo técnico sin alterar la titularidad de la orden,
salvo reasignación formal. & Entrevista con el cliente (19/10/2026), minuta A.2, E-08 &
RF-05 \\
RN-10 & Restricción de entrega y excepción a flotillas: Queda prohibida
la entrega física del vehículo con saldo pendiente. Como única
excepción, el Gerente podrá autorizar la entrega a crédito para
vehículos pertenecientes a flotillas empresariales registradas. &
Anteproyecto y entrevista con el cliente (19/10/2026), minuta A.2, E-09 & RF-10 \\
RN-11 & Vigencia de cotizaciones: Toda cotización tiene vigencia de 5
días naturales desde su emisión. Una cotización vencida no puede
autorizarse; el taller debe emitir una nueva versión. & Anteproyecto &
RF-06 \\
RN-12 & Evidencia fotográfica de recepción: Al recibir un vehículo se
exigen el checklist completo y al menos 4 fotografías (frente, trasera,
costado izquierdo y costado derecho); sin ellas la orden no puede pasar
a diagnóstico. & Anteproyecto; el mínimo de 4 fotografías debe validarse
con el cliente & RF-04 \\
RN-13 & Garantía de mano de obra: Los trabajos tienen garantía de 30
días naturales sobre la mano de obra, contados desde la entrega del
vehículo. & RF-03.2; confirmar la fuente con el cliente & RF-03.2 \\
RN-14 & Medidor de combustible descompuesto: Si el indicador de combustible
del tablero no funciona, el nivel no se deja en blanco ni se adivina: se
registra como «No verificable», junto con el nivel que declara el cliente,
el método de estimación usado (4.7.1) y la aceptación del cliente. El
taller no responde por diferencias de combustible en una recepción
marcada como «No verificable», salvo lo que él mismo haya medido. &
Propuesta del equipo; validar con el cliente & RF-04,
HU-01 \\
RN-15 & Operación de contingencia en papel: Si la falta de luz o de
internet dura más que la autonomía de los equipos (RNF-08), la operación
continúa en la Bitácora de Contingencia impresa (4.7.1) con folios
provisionales. Todo lo registrado en papel se captura en el sistema a más
tardar 24 horas después de restablecerse el servicio y conserva la hoja
física escaneada como evidencia. & Propuesta del equipo; amplía HU-25;
validar con el cliente & RNF-08, RNF-09, RF-04, RF-06,
RF-10 \\
\end{longtable}
}

\subsubsection{4.7.1 Procedimientos de
contingencia}\label{procedimientos-de-contingencia}

Estos procedimientos detallan RN-14 y RN-15. Son operativos: indican qué
hace el personal cuando falla el vehículo del cliente o la infraestructura
del taller, y qué tendrá que soportar el sistema para que esos casos queden
registrados.

\paragraph{A. Vehículo con gasolina pero con el medidor descompuesto
(RN-14).}

El nivel de combustible al ingreso protege a ambas partes ante reclamos
(«dejé medio tanque y me lo entregaron en reserva»). Cuando el indicador
del tablero no funciona, la recepcionista aplica, en este orden, el primer
método que esté disponible y anota cuál usó:

\begin{enumerate}
\tightlist
\item \textbf{Lectura por OBD-II.} Muchas veces solo falla la aguja o el
tablero, pero el sensor (flotador) del tanque sigue funcionando. Con el
escáner OBD-II del taller se lee el parámetro «Fuel Tank Level Input» (PID
\texttt{0x2F}), que da el porcentaje del tanque, y se convierte al selector
de 5 niveles (0--12\,\% Reserva, 13--37\,\% 1/4, 38--62\,\% 1/2,
63--87\,\% 3/4, 88--100\,\% Lleno). Se fotografía la pantalla del escáner.
\item \textbf{Computadora de viaje.} Si el vehículo muestra la autonomía
restante (km), se fotografía y se registra como referencia, aclarando que
es una estimación del propio vehículo.
\item \textbf{Medición física, solo si el trabajo lo requiere.} Si la
reparación implica bajar o vaciar el tanque (bomba de gasolina, sensor de
nivel, fugas), el combustible se extrae con la bomba de trasvase a bidones
graduados aprobados para combustible y se registran los \textbf{litros}
medidos, con fotografía del bidón. Al terminar, ese mismo volumen se
regresa al tanque. No se abre el sistema de combustible solo para medir:
es un riesgo de incendio y no forma parte del servicio autorizado.
\item \textbf{Declaración del cliente.} Si nada de lo anterior es posible,
se registra el nivel que dice el cliente con la leyenda «Declarado por el
cliente, no verificable».
\end{enumerate}

En todos los casos se fotografía el tablero encendido (evidencia de que el
indicador no funciona), se agrega al checklist el hallazgo «Indicador de
combustible inoperante» y se le ofrece al cliente diagnosticarlo como
servicio adicional. El cliente acepta lo registrado con su firma en la
recepción. Al entregar, el nivel se vuelve a medir con el mismo método para
que la comparación sea justa. La conversión del porcentaje OBD-II a los 5
niveles es una de las decisiones candidatas a modelarse en DMN (7.9).

\emph{Impacto en el sistema:} el selector de combustible del wireframe W-01
gana la opción «No verificable», que obliga a capturar el método de
estimación (OBD-II, computadora de viaje, medición en litros o declaración
del cliente), el valor obtenido y una fotografía del tablero.
\pendiente{agregar los campos \texttt{metodo\_combustible} y
\texttt{litros\_medidos} a RecepcionVehicular (6.2) y su caso de prueba en
la sección 5.}

\paragraph{B. Cortes prolongados de luz o de internet (RN-15).}

La operación sin conexión de la PWA (HU-25, RNF-08) cubre los cortes
cortos. Para los largos se definen tres niveles:

{\small
\begin{longtable}{@{}L{0.17\linewidth}L{0.33\linewidth}L{0.44\linewidth}@{}}
\toprule
\textbf{Nivel} & \textbf{Situación} & \textbf{Cómo se opera} \\
\midrule
\endhead
\bottomrule
\endlastfoot
1. Sin internet & Hay luz, pero no hay conexión. & La PWA trabaja con
IndexedDB y sincroniza al volver la conexión (RNF-09). El cliente no
recibe enlace ni puede autorizar en línea; se le avisa por teléfono. \\
2. Sin luz & Los equipos funcionan con batería. & Igual que el nivel 1. Se
apagan los equipos que no se usen y se reservan las baterías para
recepción y diagnóstico. \\
3. Contingencia en papel & El corte supera 4 horas o la batería de los
equipos baja del 20\,\%. & El Gerente declara la contingencia y anota la
hora de inicio en la Bitácora. Todo se registra en papel, como se describe
abajo. \\
\end{longtable}
}

La \textbf{Bitácora de Contingencia} es una libreta o carpeta con formatos
impresos y numerados de antemano, que se guarda en recepción. Contiene:

\begin{itemize}
\tightlist
\item \textbf{Hoja de recepción}: folio provisional \texttt{C-AAAAMMDD-NN}
(consecutivo del día, escrito a mano), datos del cliente y del vehículo,
kilometraje, nivel de combustible (con la opción «No verificable» de la
parte A), checklist de carrocería e interiores y firma del cliente. Las 4
fotografías obligatorias (RN-12) se toman con el celular del taller y se
nombran con el folio provisional.
\item \textbf{Hoja de diagnóstico y cotización}, con el desglose de mano de
obra y refacciones.
\item \textbf{Hoja de autorización}: como no es posible enviar el código
2FA por correo (RN-02), la única autorización válida en contingencia es la
\textbf{firma autógrafa del cliente en presencia del personal}, con fecha y
hora. No se aceptan autorizaciones por teléfono.
\item \textbf{Recibo provisional} numerado para los pagos (solo efectivo,
porque la terminal bancaria tampoco funciona). Sigue vigente RN-10: no se
entrega un vehículo con saldo pendiente.
\item \textbf{Registro de movimientos}: cambios de estado y de bahía, con
hora y nombre de quien los hizo.
\end{itemize}

\textbf{Al restablecerse el servicio}, el Gerente anota la hora de fin en
la Bitácora. En las 24 horas siguientes, la recepcionista captura cada hoja
en el sistema: la orden recibe su folio definitivo OS-AAAA-NNNN y conserva
el folio provisional como referencia, se adjuntan las fotografías y la hoja
escaneada, y las autorizaciones quedan registradas con medio «Firma en
papel». El Gerente revisa que no falte ningún folio provisional del
consecutivo y firma la Bitácora como conciliada. Las hojas físicas se
resguardan con el mismo cuidado que los datos personales del sistema (4.8).

\pendiente{agregar al modelo el medio de autorización «Firma en papel» y el
campo \texttt{folio\_provisional}; diseñar e imprimir los formatos de la
Bitácora; validar con el cliente el umbral de 4 horas y el plazo de 24
horas.}


\section{5. Plan de pruebas y
trazabilidad}\label{plan-de-pruebas-y-trazabilidad}

\subsection{5.1 Plan de pruebas}\label{plan-de-pruebas}

Estructura adaptada de ISO/IEC/IEEE 29119-3:2021. Responsable: QA Lead.

{\small
\begin{longtable}{@{}L{0.22\linewidth}L{0.74\linewidth}@{}}
\toprule
\textbf{Apartado} & \textbf{Contenido} \\
\midrule
\endhead
\bottomrule
\endlastfoot
Alcance & \textbf{Se prueba en la Fase 1:} el flujo núcleo del release v0.1.0 --- recepción con checklist y
fotografías, asignación de bahía y mecánico, diagnóstico, cotización con vigencia, autorización con código
2FA (remota y presencial), re-cotización por sobrecosto, pago, entrega, crédito a flotillas, citas en línea,
lista de órdenes abiertas, panel del gerente y control de acceso por rol.
\textbf{No se prueba en la Fase 1:} almacén y compras (RF-08.x), operación sin conexión y sincronización
(RNF-08, RNF-09; dependen de la PWA), reportes mensuales (HU-29) y envío real de correo o SMS (R-05). \\
Niveles de prueba & \textbf{Unitarias y de integración:} pruebas automatizadas con pytest sobre la API completa
(servicio + base de datos), en \texttt{backend/tests/}. \textbf{Sistema:} recorrido del flujo núcleo en el
ambiente de pruebas (desde el Sprint 4). \textbf{Aceptación:} los casos CPA de 5.2, ejecutados con el
Product Owner y validados con el cliente en las Reviews. \\
Enfoque & Automatizado primero: cada regla de negocio (RN) y cada requisito de acceso tiene al menos una
prueba automatizada; se incluyen valores límite (5 vs.\ 6 días, 10\,\% vs.\ 10.06\,\%, 5 min vs.\ 5 min 1 s).
Manual para usabilidad y dispositivos (RNF-01 a RNF-03). Herramientas: pytest, httpx/TestClient,
GitHub Actions; Lighthouse/axe para accesibilidad (Sprint 5). \\
Entornos & \textbf{Local:} SQLite en memoria, reloj controlable para probar vigencias.
\textbf{CI:} GitHub Actions ejecuta las pruebas, carga \texttt{db/schema.sql} en PostgreSQL 16 y compila este
documento. \textbf{Staging:} hospedaje y URL según ADR-004 (Sprint 4). Siempre datos ficticios (RNF-05). \\
Criterios de entrada & La historia cumple la DoR; existe su caso CPA en 5.2; el código está en un Pull
Request hacia \texttt{main}. \\
Criterios de salida & 100\,\% de las pruebas automatizadas pasan en CI; 0 defectos S1/S2 abiertos; todos
los requisitos Must tienen al menos un caso que pasa (matriz 5.3); el PO aceptó las historias. \\
Roles & \textbf{QA Lead:} diseña los casos, mantiene la matriz y la CI. \textbf{Developers:} escriben pruebas
junto con el código. \textbf{PO:} ejecuta las pruebas de aceptación. \textbf{Equipo A:} QA cruzado.
\textbf{Cliente:} valida en las Reviews. \\
Riesgos de las pruebas & Las pruebas con SQLite podrían ocultar diferencias con PostgreSQL (mitigación:
la CI carga el esquema real; pendiente agregar pruebas de integración contra PostgreSQL). El envío de
códigos se simula con una bandeja en memoria. Sin dispositivos reales todavía para RNF-01/02. \\
Calendario & S2: casos CPA del flujo núcleo y CI. S3: pruebas de citas y panel. S4: pruebas en staging y
de rendimiento (RNF-06). S5: usabilidad, accesibilidad y SUS. S6: ejecución final, QA cruzado y reporte
(sección 10). \\
\end{longtable}
}

\subsection{5.2 Casos de prueba de
aceptación}\label{casos-de-prueba-de-aceptaciuxf3n}

Cada caso está automatizado en \texttt{backend/tests/} con el prefijo de su ID (por ejemplo,
\texttt{test\_cpa14\_codigo\_expira\_a\_los\_cinco\_minutos}). Resultado de la ejecución del 29/10/2026:
35 de 35 pasan. Los datos son ficticios.

{\footnotesize
\begin{longtable}{@{}L{0.07\linewidth}L{0.13\linewidth}L{0.17\linewidth}L{0.27\linewidth}L{0.24\linewidth}L{0.07\linewidth}@{}}
\toprule
\textbf{ID} & \textbf{HU / CU (regla)} & \textbf{Precondiciones} & \textbf{Pasos} & \textbf{Resultado esperado} & \textbf{Res.} \\
\midrule
\endhead
\bottomrule
\endlastfoot
CPA-01 & HU-01, CU-02 & Sesión de Recepcionista; vehículo registrado & 1. Registrar recepción con kilometraje, combustible y checklist. & Se crea la orden con folio OS-AAAA-NNNN en estado Abierta y su primer registro en el historial. & Pasa \\
CPA-02 & HU-01, HU-09 (RN-12) & Orden con 3 fotografías & 1. Asignar bahía. 2. Registrar diagnóstico. 3. Subir la 4.ª foto y repetir. & Sin las 4 fotos se rechaza (RN-12); con ellas la orden pasa a Diagnosticada. & Pasa \\
CPA-03 & HU-01 (RN-12) & Orden con checklist incompleto & 1. Intentar diagnóstico. 2. Completar checklist. 3. Repetir. & Primero se rechaza; tras completar el checklist se acepta. & Pasa \\
CPA-04 & HU-01 & Orden abierta con foto Frente & 1. Subir otra foto Frente. 2. Subir un archivo que no es imagen. & Foto repetida rechazada (409); archivo no imagen rechazado (415). & Pasa \\
CPA-05 & HU-07 (RN-05) & Las 4 bahías ocupadas & 1. Registrar una nueva recepción. & La orden queda EnEsperaDeBahia y la respuesta trae la alerta «sin bahías». & Pasa \\
CPA-06 & HU-07 (RN-05) & 4 bahías ocupadas y 5 vehículos en espera & 1. Registrar otra recepción. & Se rechaza: estacionamiento lleno. & Pasa \\
CPA-07 & HU-08 (RN-05) & Bahía 2 ocupada & 1. Asignar la bahía 2 a otra orden. 2. Cancelar la primera orden. 3. Repetir. & Primero se rechaza; al cancelar se libera y la asignación procede. & Pasa \\
CPA-08 & HU-08 (RN-09) & Orden abierta & 1. Asignar como responsable a una recepcionista. & Se rechaza: el responsable debe ser mecánico activo. & Pasa \\
CPA-09 & CU-02 & Vehículo con orden abierta & 1. Abrir otra orden para el mismo vehículo. & Se rechaza. & Pasa \\
CPA-10 & HU-12, HU-02 & Orden Diagnosticada & 1. Generar cotización con mano de obra, refacción y un concepto opcional. & Total \$1,750.00; vence a los 5 días; el cliente recibe su enlace seguro. & Pasa \\
CPA-11 & HU-14 (RN-11) & Cotización emitida el día 0 & 1. Día 5: pedir código. 2. Día 6: confirmar. & Día 5 vigente; día 6 se rechaza por vencida. & Pasa \\
CPA-12 & HU-13 (RN-03) & Cotización pendiente & 1. Deseleccionar el concepto opcional. 2. Deseleccionar uno obligatorio. & Total baja a \$1,500.00; el obligatorio no se puede quitar. & Pasa \\
CPA-13 & HU-14, CU-11 (RN-02, RNF-04) & Cotización pendiente & 1. El cliente pide código desde su enlace. 2. Lo confirma. & Autorización Confirmada con IP, token y huella SHA-256 válida; orden Autorizada; el cambio no tiene usuario interno. & Pasa \\
CPA-14 & HU-14 (RN-02) & Código emitido & 1. Esperar 5 min 1 s. 2. Confirmar. & Se rechaza: código expirado. & Pasa \\
CPA-15 & HU-14 (RN-02) & Código emitido & 1. Confirmar con código incorrecto. 2. Con el correcto. 3. Repetir el correcto. & Incorrecto rechazado; correcto aceptado; segundo uso rechazado. & Pasa \\
CPA-16 & CU-03 & Cotización pendiente; cliente en mostrador & 1. Recepcionista pide el código. 2. Captura el que dicta el cliente. & Autorización Presencial confirmada. & Pasa \\
CPA-17 & HU-15 & Cotización inicial pendiente & 1. El cliente rechaza desde su enlace. & Cotización Rechazada y orden Cancelada. & Pasa \\
CPA-18 & HU-16 (RN-04) & Orden Autorizada por \$1,750.00 & 1. Re-cotizar a \$1,926.00 (+10.06\,\%). 2. Terminar trabajo. 3. Autorizar y repetir. & Queda Pendiente y bloquea el avance; tras autorizar se termina. & Pasa \\
CPA-19 & HU-16 (RN-04) & Orden Autorizada por \$1,750.00 & 1. Re-cotizar a \$1,925.00 (+10\,\% exacto). & Se autoriza sin detener el trabajo; total autorizado \$1,925.00. & Pasa \\
CPA-20 & RNF-04 & Autorización confirmada & 1. Modificar el monto directamente en la base. 2. Verificar la huella. & La verificación detecta la alteración. & Pasa \\
CPA-21 & HU-21 (RN-10) & Orden lista; pago parcial & 1. Registrar la entrega. & Se rechaza por saldo pendiente. & Pasa \\
CPA-22 & HU-23, CU-04 (RN-01, RN-13) & Orden lista & 1. Pagar el total. 2. Entregar. 3. Abrir el enlace del cliente. & Orden Cerrada, garantía de 30 días, bahía libre, enlace expirado. & Pasa \\
CPA-23 & HU-23 & Orden lista & 1. Pagar un centavo más que el saldo. & Se rechaza. & Pasa \\
CPA-24 & HU-22, CU-13 (RN-10) & Orden lista de un cliente particular & 1. Recepcionista autoriza crédito. 2. Gerente lo intenta. & Recepcionista: 403. Gerente: rechazado porque no es flotilla. & Pasa \\
CPA-25 & HU-22, CU-13 (RN-10) & Orden lista de flotilla con crédito & 1. Gerente autoriza crédito. 2. Entregar con saldo. & Entrega permitida; el enlace sigue activo mientras haya saldo. & Pasa \\
CPA-26 & RF-01 & Sin sesión & 1. Consultar órdenes sin token y con token falso. & 401 en ambos casos. & Pasa \\
CPA-27 & RF-01 & Usuarios registrados & 1. Entrar con contraseña incorrecta. 2. Entrar como mecánico de bahía. & 401 en ambos casos. & Pasa \\
CPA-28 & RF-01, HU-27 & Sesiones de cada rol & 1. Recepcionista asigna bahía. 2. Jefe cancela orden. 3. Recepcionista abre el panel. 4. Almacén lista órdenes. 5. Recepcionista crea un gerente. & 403 en todos los casos. & Pasa \\
CPA-29 & HU-27 & Sesión de Gerente & 1. Dar de alta a una recepcionista. 2. Iniciar sesión con ella. & Alta correcta (correo normalizado) y sesión con rol Recepcionista. & Pasa \\
CPA-30 & RF-02.1, CU-10 (OWASP A01) & Cliente con enlace & 1. Abrir su enlace. 2. Alterar un carácter. 3. Usar un folio como enlace. & Solo su orden es visible; 2 y 3 dan 404. & Pasa \\
CPA-31 & HU-03, HU-30 (RN-08) & Público & 1. Solicitar cita sin aceptar el aviso. 2. Aceptándolo. & Sin aviso se rechaza; con aviso queda Pendiente y sin bahía. & Pasa \\
CPA-32 & HU-04, CU-01 & Dos citas pendientes a la misma hora & 1. Confirmar ambas en la bahía 3. 2. La segunda en la bahía 4. & La segunda en bahía 3 se rechaza; en bahía 4 se confirma. & Pasa \\
CPA-33 & CU-02 (RN-08) & Cita pendiente & 1. Abrir orden desde la cita. 2. Confirmarla y repetir. & Primero se rechaza; confirmada sí abre la orden. & Pasa \\
CPA-34 & HU-03 & Placas registradas a otro cliente & 1. Solicitar cita con esas placas y otro contacto. & Se rechaza. & Pasa \\
CPA-35 & RF-11, flujo núcleo & Una orden diagnosticada y una cancelada & 1. Consultar órdenes abiertas. 2. Consultar el panel. & Solo aparece la diagnosticada, con bahía y responsable; panel con 1 bahía ocupada. & Pasa \\
\end{longtable}
}

\subsection{5.3 Matriz de trazabilidad}\label{matriz-de-trazabilidad}

Los módulos se refieren a \texttt{backend/app/servicios/} (lógica) y \texttt{backend/app/rutas/} (API).
Resultado de la ejecución automatizada del 29/10/2026.

{\footnotesize
\begin{longtable}{@{}L{0.09\linewidth}L{0.17\linewidth}L{0.12\linewidth}L{0.2\linewidth}L{0.24\linewidth}L{0.1\linewidth}@{}}
\toprule
\textbf{Requisito} & \textbf{Historias} & \textbf{Caso de uso} & \textbf{Casos de prueba} & \textbf{Componente / módulo} & \textbf{Resultado} \\
\midrule
\endhead
\bottomrule
\endlastfoot
RF-01 (Must) & HU-27 & --- & CPA-26 a CPA-29 & seguridad.py, rutas/personal.py & Pasa \\
RF-02.1 (Must) & HU-02, HU-28 & CU-10, CU-11 & CPA-10, CPA-13, CPA-30 & rutas/portal.py, servicios/cotizaciones.py & Pasa (API; falta PWA) \\
RF-02.2 (Should) & HU-02, HU-23 & CU-04 & CPA-22, CPA-25 & servicios/cobro.py & Pasa \\
RF-02.3 (Could) & HU-02 & --- & --- & --- & Fase 2 \\
RF-03.1 (Must) & HU-05, HU-06 & CU-02 & CPA-05, CPA-34 & rutas/personal.py & Pasa \\
RF-03.2 (Should) & HU-23, HU-24 & CU-07 & CPA-22 & servicios/cobro.py (garantía) & Pasa (parcial) \\
RF-04 (Must) & HU-01, HU-25, HU-26 & CU-02 & CPA-01 a CPA-04 & servicios/ordenes.py & Pasa (sin modo offline) \\
RF-05 (Must) & HU-07, HU-08, HU-09, HU-11 & CU-05 & CPA-05 a CPA-09, CPA-35 & servicios/ordenes.py & Pasa \\
RF-06 (Must) & HU-02, HU-12 a HU-15 & CU-03, CU-05, CU-11 & CPA-10 a CPA-17 & servicios/cotizaciones.py & Pasa \\
RF-07 (Must) & HU-16 & CU-07 & CPA-18, CPA-19 & servicios/cotizaciones.py, cobro.py & Pasa \\
RF-08.1 (Should) & HU-10, HU-19 & CU-06 & \pendiente{} & --- & No iniciado \\
RF-08.2 (Should) & HU-17, HU-18 & CU-08 & \pendiente{} & --- & No iniciado \\
RF-08.3 (Could) & HU-18 & CU-08 & --- & --- & No iniciado \\
RF-09 (Should) & HU-03, HU-04, HU-06 & CU-01, CU-09 & CPA-31 a CPA-34 & servicios/citas.py & Pasa \\
RF-10 (Must) & HU-21, HU-22, HU-23 & CU-04, CU-13 & CPA-21 a CPA-25 & servicios/cobro.py & Pasa \\
RF-11 (Should) & HU-29 & CU-12 & CPA-35 & rutas/taller.py (panel) & Pasa (parcial) \\
RNF-04 & HU-14, HU-16 & CU-03, CU-11 & CPA-13, CPA-20 & servicios/cotizaciones.py (huella) & Pasa \\
RNF-05 & HU-30 & --- & CPA-31; datos ficticios en todas & seguridad.py & Pasa (parcial: falta HTTPS en staging) \\
\end{longtable}
}

Todos los requisitos \textbf{Must} tienen al menos una historia y un caso de prueba que pasa. Quedan sin
caso los requisitos del almacén (RF-08.x, Should/Could), planificados para el Sprint 4.

\section{6. Modelos de análisis}\label{modelos-de-anuxe1lisis}

Los modelos usan los nombres del glosario (4.2) y del SRS, y coinciden con el esquema
\texttt{db/schema.sql} y con el código de \texttt{backend/app/}. Notación UML 2.5.1.

\subsection{6.1 Diagramas de flujo de
datos}\label{diagramas-de-flujo-de-datos}

\figurapendiente{6.1}{Diagrama de flujo de datos de contexto}{el sistema como un único proceso
«Sistema de Órdenes de Servicio y Citas» y sus entidades externas: Cliente (envía solicitud de cita,
selección de conceptos, código 2FA, pago; recibe enlace seguro, cotización, código y estado), Recepcionista
(datos de recepción, checklist, fotografías, pagos, entrega), Jefe de Mecánicos (asignación, diagnóstico,
cotización, fin de trabajo), Gerente (alta de personal, crédito a flotillas; recibe panel) y Servicio de
correo/SMS (recibe mensajes salientes).}{docs/06-analisis/dfd/contexto.drawio}

\figurapendiente{6.2}{Diagrama de flujo de datos de nivel 0}{los procesos 1.0 Gestionar citas, 2.0
Recibir vehículo, 3.0 Asignar y diagnosticar, 4.0 Cotizar y autorizar, 5.0 Cobrar y entregar y 6.0
Consultar estado/panel; y los almacenes D1 Clientes y vehículos, D2 Citas, D3 Órdenes e historial, D4
Cotizaciones y autorizaciones, D5 Pagos y garantías, D6 Tokens de acceso. Los flujos con las entidades
externas deben ser los mismos que en la Figura 6.1 (balanceo).}{docs/06-analisis/dfd/nivel-0.drawio}

\figurapendiente{6.3}{Diagrama de flujo de datos de nivel 1 del proceso 4.0 Cotizar y
autorizar}{subprocesos 4.1 Generar cotización (lee D3, escribe D4), 4.2 Enviar enlace (escribe D6, envía
al servicio de correo), 4.3 Seleccionar conceptos opcionales, 4.4 Emitir código 2FA, 4.5 Validar código y
registrar autorización (escribe D4 y el historial en D3), 4.6 Rechazar cotización. Entradas y salidas
balanceadas con el proceso 4.0 de la Figura 6.2.}{docs/06-analisis/dfd/nivel-1-cotizar.drawio}

\subsection{6.2 Modelo de datos}\label{modelo-de-datos}

El modelo relacional está en 3FN: cada tabla describe una sola entidad, los atributos derivados del
diagrama de clases (/total, /saldoPendiente, /fechaVencimiento, /fechaFin) no se almacenan sino que se
calculan en vistas o columnas generadas, y la especialización \{disjoint, complete\} de Cliente usa una
tabla por subclase. El esquema completo, con sus restricciones y disparadores, está en
\texttt{db/schema.sql} y se valida en PostgreSQL 16 en cada ejecución de la integración continua.

\figurapendiente{6.4}{Modelo entidad-relación}{las 26 tablas de \texttt{db/schema.sql} con sus llaves
primarias y foráneas y la cardinalidad de cada relación; resaltar el núcleo cliente --- vehículo ---
orden\_servicio --- recepcion\_vehicular (checklist, fotografías) --- cotizacion (detalles,
autorizaciones, pagos) --- token\_acceso. Puede generarse con DBeaver o pgAdmin a partir del
esquema.}{db/schema.sql (exportar a docs/06-analisis/er/modelo-er.png)}

\textbf{Diccionario de datos} de las tablas del flujo núcleo (las demás se documentan en los
comentarios de \texttt{db/schema.sql}):

{\footnotesize
\begin{longtable}{@{}L{0.15\linewidth}L{0.17\linewidth}L{0.13\linewidth}L{0.06\linewidth}L{0.2\linewidth}L{0.21\linewidth}@{}}
\toprule
\textbf{Entidad} & \textbf{Atributo} & \textbf{Tipo} & \textbf{Nulo} & \textbf{Restricciones} & \textbf{Descripción} \\
\midrule
\endhead
\bottomrule
\endlastfoot
cliente & id\_cliente & INTEGER & No & PK, identidad & Identificador del cliente \\
 & tipo\_cliente & VARCHAR(10) & No & Particular o Flotilla & Subclase del cliente \\
 & telefono & VARCHAR(20) & No & & Contacto (dato personal) \\
 & correo & VARCHAR(160) & No & & Recibe enlace y códigos 2FA \\
vehiculo & id\_vehiculo & INTEGER & No & PK & Identificador \\
 & id\_cliente & INTEGER & No & FK cliente & Propietario \\
 & placas & VARCHAR(15) & No & Único & Placas en mayúsculas \\
 & vin & VARCHAR(17) & Sí & Único & Número de serie (opcional) \\
 & tipo & tipo\_vehiculo & No & Automovil, Motocicleta, Camioneta & Solo vehículos ligeros (RN-07) \\
orden\_servicio & folio & VARCHAR(20) & No & PK, formato OS-AAAA-NNNN & Folio de la orden \\
 & estado & estado\_orden & No & 7 valores (6.5) & Estado actual \\
 & numero\_bahia & SMALLINT & Sí & FK bahia; único entre órdenes activas & Bahía asignada (RN-05) \\
 & id\_usuario\_responsable & INTEGER & Sí & FK usuario & Mecánico responsable (RN-09) \\
 & entrega\_a\_credito & BOOLEAN & No & Requiere gerente y flotilla & Excepción de RN-10 \\
 & fecha\_entrega & TIMESTAMPTZ & Sí & $\geq$ fecha\_apertura & Momento de la entrega \\
recepcion\_vehicular & kilometraje & INTEGER & No & $\geq 0$ & Lectura del odómetro \\
 & nivel\_combustible & nivel\_combustible & No & 5 niveles & Combustible al ingreso \\
 & id\_cita & INTEGER & Sí & FK cita, único; solo cita Confirmada & Cita de origen \\
fotografia & posicion & posicion\_foto & No & Una por posición salvo «Otra» & Frente, Trasera, Izquierda, Derecha (RN-12) \\
 & archivo & VARCHAR(255) & No & & Ruta en el almacenamiento \\
cotizacion & version & INTEGER & No & $\geq 1$; único con folio & Versión de la cotización \\
 & fecha\_emision & DATE & No & & Inicio de vigencia \\
 & vigencia\_dias & INTEGER & No & $> 0$, por omisión 5 & RN-11 \\
 & estado & estado\_cotizacion & No & Pendiente, Autorizada, Rechazada, Vencida & Estado (6.5) \\
detalle\_cotizacion & tipo & tipo\_concepto & No & ManoDeObra, Refaccion, Insumo & Clase de concepto \\
 & precio\_unitario & NUMERIC(10,2) & No & $\geq 0$ & Precio en pesos \\
 & opcional / seleccionado & BOOLEAN & No & Solo un opcional puede deseleccionarse & RN-03 \\
autorizacion & medio & medio\_autorizacion & No & Remoto exige IP y token & Canal de la autorización \\
 & tipo & tipo\_autorizacion & No & Inicial o IncrementoMayor10 & RN-04 \\
 & monto\_aceptado & NUMERIC(10,2) & No & $\geq 0$ & Monto que el cliente aceptó \\
 & codigo\_verificacion\_hash & VARCHAR(255) & No & PBKDF2; nunca en claro & Código 2FA (RN-02) \\
 & huella\_sha256 & CHAR(64) & Sí & Obligatoria si Confirmada & Detecta alteraciones (RNF-04) \\
token\_acceso & token\_hash & CHAR(64) & No & PK; SHA-256 del token & El token en claro solo viaja en el enlace \\
 & fecha\_expiracion & TIMESTAMPTZ & Sí & & Se fija al entregar y saldar (RN-01) \\
pago & monto & NUMERIC(10,2) & No & $> 0$; solo sobre cotización Autorizada & Abono del cliente \\
\end{longtable}
}

\subsection{6.3 Diagrama de clases del
dominio}\label{diagrama-de-clases-del-dominio}

El diagrama de clases v4 tiene 26 clases del dominio (mínimo del Equipo B: 12), con atributos,
asociaciones, multiplicidades, composiciones, una generalización \{disjoint, complete\} (Cliente →
ClienteParticular | Flotilla), 15 enumeraciones y las restricciones de negocio como notas OCL
informales. Las clases marcadas «OfflineSync» (RecepcionVehicular, ChecklistInspeccion, Diagnostico) se
capturan sin conexión (RNF-08). Se divide en dos láminas por legibilidad; OrdenServicio, Usuario y
TokenAcceso aparecen en ambas como vista parcial.

\figura{../06-analisis/clases/diagrama-clases-v4-1.png}{angle=90,width=\linewidth,height=0.8\textheight}{6.5a}{Diagrama
de clases del dominio (1/2): clientes, citas, recepción y orden}{docs/06-analisis/clases/diagrama-clases-v4-1.png}

\figura{../06-analisis/clases/diagrama-clases-v4-2.png}{angle=90,width=\linewidth,height=0.8\textheight}{6.5b}{Diagrama
de clases del dominio (2/2): cotización, compras, pagos y garantía}{docs/06-analisis/clases/diagrama-clases-v4-2.png}

\subsection{6.4 Diagrama de objetos}\label{diagrama-de-objetos}

\figurapendiente{6.6}{Diagrama de objetos: orden autorizada por el cliente}{instantánea con datos
ficticios del caso CPA-13: \texttt{c1:ClienteParticular} (nombre = «Cliente Ficticio»), \texttt{v1:Vehiculo}
(placas = «PRB-001», Nissan Versa 2019), \texttt{os1:OrdenServicio} (folio = «OS-2026-0001», estado =
Autorizada, bahía 1), \texttt{r1:RecepcionVehicular} (kilometraje = 45210, combustible = Mitad) con su
\texttt{ch1:ChecklistInspeccion} (completo = true) y 4 \texttt{Fotografia}; \texttt{d1:Diagnostico};
\texttt{cot1:Cotizacion} (versión 1, estado = Autorizada, total = 1750.00) con 3 \texttt{DetalleCotizacion};
\texttt{a1:Autorizacion} (medio = Remoto, estado = Confirmada, monto = 1750.00) ligada a
\texttt{t1:TokenAcceso}; \texttt{m1:Usuario} (rol = Mecanico) como responsable.}{docs/06-analisis/objetos/orden-autorizada.puml}

\subsection{6.5 Diagramas de estados}\label{diagramas-de-estados}

Las transiciones de la orden están programadas en una sola tabla (\texttt{TRANSICIONES} en
\texttt{backend/app/servicios/ordenes.py}); cualquier transición que no aparezca aquí es rechazada por
el sistema y toda transición válida queda en HistorialEstado.

\figurapendiente{6.7}{Máquina de estados de la Orden de servicio}{los 7 estados (Abierta,
EnEsperaDeBahia, Diagnosticada, Autorizada, ListaParaEntrega, Cerrada, Cancelada) con las transiciones de
la tabla siguiente, cada una con evento, [condición] y / rol.}{docs/06-analisis/estados/orden-servicio.puml}

{\footnotesize
\begin{longtable}{@{}L{0.15\linewidth}L{0.19\linewidth}L{0.31\linewidth}L{0.14\linewidth}L{0.15\linewidth}@{}}
\toprule
\textbf{Estado origen} & \textbf{Evento} & \textbf{Condición} & \textbf{Estado destino} & \textbf{Rol autorizado} \\
\midrule
\endhead
\bottomrule
\endlastfoot
(inicio) & Registrar recepción & El vehículo no tiene otra orden abierta; si no hay bahía, hay lugar en el estacionamiento (RN-05) & Abierta & Recepcionista \\
Abierta & Recepción sin bahía libre & Las 4 bahías están ocupadas & EnEsperaDeBahia & Recepcionista \\
EnEsperaDeBahia & Asignar bahía & Bahía libre; responsable es mecánico activo (RN-09) & Abierta & Jefe de Mecánicos \\
Abierta & Registrar diagnóstico & Checklist completo y 4 fotografías (RN-12); bahía y responsable asignados & Diagnosticada & Jefe de Mecánicos \\
Diagnosticada & Confirmar código 2FA & Cotización vigente (RN-11); código correcto, de un solo uso y con menos de 5 minutos (RN-02) & Autorizada & Cliente (remoto) o Recepcionista (presencial) \\
Diagnosticada & Rechazar cotización & Cotización vigente & Cancelada & Cliente \\
Autorizada & Terminar trabajo & No hay re-cotización pendiente mayor al 10\,\% (RN-04) & ListaParaEntrega & Jefe de Mecánicos \\
ListaParaEntrega & Registrar entrega & Saldo pendiente igual a 0, o entrega a crédito autorizada (RN-10) & Cerrada & Recepcionista \\
Abierta, EnEsperaDeBahia, Diagnosticada & Cancelar & --- & Cancelada & Gerente \\
\end{longtable}
}

\figurapendiente{6.8}{Máquina de estados de la Cotización}{los estados Pendiente, Autorizada,
Rechazada y Vencida, con las transiciones de la tabla siguiente.}{docs/06-analisis/estados/cotizacion.puml}

{\footnotesize
\begin{longtable}{@{}L{0.15\linewidth}L{0.19\linewidth}L{0.31\linewidth}L{0.14\linewidth}L{0.15\linewidth}@{}}
\toprule
\textbf{Estado origen} & \textbf{Evento} & \textbf{Condición} & \textbf{Estado destino} & \textbf{Rol autorizado} \\
\midrule
\endhead
\bottomrule
\endlastfoot
(inicio) & Generar cotización & Orden Diagnosticada (inicial) o Autorizada (re-cotización); al menos un concepto; sin otra cotización pendiente vigente & Pendiente & Recepcionista o Jefe de Mecánicos \\
(inicio) & Re-cotizar & Orden Autorizada y nuevo total $\leq$ total inicial autorizado $\times$ 1.10 (RN-04) & Autorizada & Recepcionista o Jefe de Mecánicos \\
Pendiente & Cambiar selección & Solo conceptos opcionales (RN-03); invalida códigos ya emitidos & Pendiente & Cliente o Recepcionista \\
Pendiente & Confirmar código 2FA & Vigente y código válido (RN-02) & Autorizada & Cliente o Recepcionista \\
Pendiente & Rechazar & Vigente & Rechazada & Cliente \\
Pendiente & Pasa la fecha de vencimiento & Hoy $>$ fecha de emisión + 5 días (RN-11) & Vencida & Sistema \\
\end{longtable}
}

\subsection{6.6 Diagramas de actividades
(to-be)}\label{diagramas-de-actividades-to-be}

\figurapendiente{6.9}{Actividad to-be: recepción y diagnóstico}{carriles Cliente, Recepcionista,
Sistema y Jefe de mecánicos: búsqueda o registro del cliente con aviso de privacidad, captura de
kilometraje, combustible, checklist y 4 fotografías, apertura de la orden, decisión «¿hay bahía libre?»
(EnEsperaDeBahia con alerta), asignación de bahía y responsable, registro del diagnóstico y validación de
RN-12.}{docs/06-analisis/actividades/recepcion-diagnostico.puml}

\figurapendiente{6.10}{Actividad to-be: cotización y autorización con código 2FA}{carriles Jefe de
mecánicos, Sistema y Cliente: generación de la cotización con vigencia de 5 días, envío del enlace seguro,
selección de opcionales, decisión aceptar/rechazar, emisión y validación del código de 6 dígitos, registro
de la autorización con huella SHA-256, ejecución del trabajo y re-cotización bloqueante si el adicional
supera el 10\,\%.}{docs/06-analisis/actividades/cotizacion-autorizacion.puml}

\section{7. Arquitectura y diseño}\label{arquitectura-y-diseuxf1o}

\subsection{7.1 Estilo arquitectónico}\label{estilo-arquitectuxf3nico}

El sistema sigue un estilo \textbf{cliente-servidor con API REST}. El servidor es un \textbf{monolito
modular en capas} (ADR-001):

\begin{itemize}
\tightlist
\item \textbf{Rutas} (\texttt{app/rutas/}): traducen HTTP a llamadas de servicio, validan la entrada con
Pydantic y declaran qué rol puede usar cada operación.
\item \textbf{Servicios} (\texttt{app/servicios/}): toda la lógica y las reglas de negocio RN-xx. Es la
única capa que cambia el estado de una orden o de una cotización.
\item \textbf{Modelo} (\texttt{app/models.py}): entidades del dominio, iguales a las clases de 6.3 y a las
tablas de \texttt{db/schema.sql}.
\item \textbf{Datos}: PostgreSQL (ADR-003), que además refuerza las restricciones con CHECK, índices únicos
parciales y disparadores.
\end{itemize}

Los clientes son dos interfaces web que consumen la misma API: la del personal del taller (tableta o
laptop) y el portal del cliente, al que se entra con el enlace seguro (ADR-005).
Las interfaces web (PWA) se construyen en los Sprints 4 y 5; hoy la API se opera desde su
documentación interactiva en /docs.

{\small
\begin{longtable}{@{}L{0.14\linewidth}L{0.82\linewidth}@{}}
\toprule
\textbf{RNF} & \textbf{Decisión de arquitectura que lo atiende} \\
\midrule
\endhead
\bottomrule
\endlastfoot
RNF-04 & Una sola capa de servicios registra las autorizaciones: guarda fecha y hora UTC, IP, medio, monto
y una huella SHA-256 de lo aceptado, que se recalcula para detectar alteraciones (CPA-20). El historial de
estados y los diagnósticos son registros inmutables (disparador \texttt{fn\_registro\_inmutable}). \\
RNF-05 & Contraseñas con PBKDF2-SHA256; tokens del cliente y códigos 2FA guardados solo como hash; control
de acceso por rol declarado en cada ruta; el portal identifica la orden por el token y no por el folio
(OWASP A01, CPA-30); HTTPS lo provee el hospedaje (ADR-004). \\
RNF-07 & Toda la pila es de código abierto: Python, FastAPI, SQLAlchemy, PostgreSQL, pytest (ADR-002 y
ADR-003). \\
RNF-08 y RNF-09 & Las entidades «OfflineSync» llevan \texttt{uuid\_local} único, generado en el dispositivo:
reenviar un registro ya sincronizado no lo duplica (sincronización idempotente). La captura sin conexión
vivirá en la PWA con IndexedDB. Endpoint de sincronización por lotes: \pendiente{}. \\
RNF-10 & Las operaciones del portal son seguras de reintentar: pedir un código nuevo expira el anterior y
un código solo se acepta una vez (CPA-15), así que un reintento automático no duplica autorizaciones. \\
RNF-01 & El portal del cliente es web y responsivo; no requiere instalar nada ni conceder permisos. \\
\end{longtable}
}

\subsection{7.2 Diagramas C4}\label{diagramas-c4}

\figurapendiente{7.1}{C4 nivel 1: contexto del sistema}{el sistema y sus usuarios: Cliente (portal web
con enlace seguro), Recepcionista, Jefe de Mecánicos, Gerente y Encargado de Almacén; sistemas externos:
servicio de correo electrónico (enlaces y códigos 2FA) y, en Fase 2, servicio de SMS.}{docs/07-arquitectura/c4/contexto.puml}

\figurapendiente{7.2}{C4 nivel 2: contenedores}{PWA del personal (navegador en tableta/laptop, con
IndexedDB para operar sin conexión), portal del cliente (navegador móvil), API REST FastAPI (Python 3.12),
base de datos PostgreSQL 16, almacenamiento de fotografías y servicio de correo; protocolos HTTPS/JSON
entre navegador y API y SQL (psycopg) entre API y base de datos.}{docs/07-arquitectura/c4/contenedores.puml}

\figurapendiente{7.3}{C4 nivel 3: componentes de la API}{componentes Rutas del personal
(\texttt{rutas/personal.py}), Rutas del taller (\texttt{rutas/taller.py}), Portal (\texttt{rutas/portal.py}),
Seguridad (\texttt{seguridad.py}), Servicios de órdenes, citas, cotizaciones y cobro
(\texttt{servicios/}), Notificaciones (\texttt{notificaciones.py}) y Modelo/ORM
(\texttt{models.py}, \texttt{db.py}), con sus dependencias.}{docs/07-arquitectura/c4/componentes-api.puml}

\subsection{7.3 Registros de decisiones de arquitectura
(ADR)}\label{registros-de-decisiones-de-arquitectura-adr}

{\small
\begin{longtable}{@{}L{0.2\linewidth}L{0.76\linewidth}@{}}
\toprule
\textbf{Campo} & \textbf{Contenido} \\
\midrule
\endhead
\bottomrule
\endlastfoot
ID y título & \textbf{ADR-001 Estilo arquitectónico: monolito modular en capas con API REST} \\
Fecha y estado & 29/10/2026 · Aceptado (reemplaza la propuesta de microservicios de la Propuesta de
Solución Tecnológica) \\
Contexto & Equipo de 5 personas con poca experiencia en FastAPI (R-07), 6 sprints, hospedaje gratuito y
reglas de negocio que cruzan varias entidades en una misma transacción (autorizar cambia cotización,
orden e historial). \\
Alternativas evaluadas & 1. Microservicios (citas, órdenes, cotizaciones, cobro). 2. Monolito MVC con
páginas generadas en el servidor. 3. Monolito modular en capas que expone una API REST. \\
Decisión & Opción 3. Los módulos de servicio están separados por área y solo dependen hacia abajo
(7.5), de modo que podrían separarse después si hiciera falta. \\
Consecuencias & \textbf{+} Una sola transacción por operación; un solo despliegue; pruebas de extremo a
extremo sencillas. \textbf{+} La API sirve igual a la PWA del personal y al portal del cliente.
\textbf{--} Escala como una sola unidad (suficiente para un taller de 4 bahías). \textbf{--} Hay que vigilar
que las rutas no contengan lógica de negocio. \\
RNF relacionados & RNF-04, RNF-07, RNF-10 \\
\midrule
ID y título & \textbf{ADR-002 Pila tecnológica: Python 3.12 + FastAPI + SQLAlchemy 2} \\
Fecha y estado & 29/10/2026 · Aceptado (la interfaz web queda por decidir en el Sprint 3) \\
Contexto & La propuesta al cliente comprometió Python con FastAPI y herramientas gratuitas. Se necesita
validación de datos, documentación automática de la API y un ORM que funcione con PostgreSQL y SQLite. \\
Alternativas evaluadas & 1. Django + Django REST Framework. 2. Node.js + Express. 3. FastAPI + SQLAlchemy +
Pydantic. \\
Decisión & Opción 3: validación declarativa, documentación OpenAPI en \texttt{/docs} sin trabajo extra y
pruebas con \texttt{TestClient}. \\
Consecuencias & \textbf{+} Curva de aprendizaje corta para un equipo que ya usa Python. \textbf{--} FastAPI
no trae panel de administración ni autenticación integrada: se implementó la propia con la biblioteca
estándar (\texttt{seguridad.py}). \\
RNF relacionados & RNF-07, RNF-05 \\
\midrule
ID y título & \textbf{ADR-003 Base de datos: PostgreSQL 16} \\
Fecha y estado & 29/10/2026 · Aceptado \\
Contexto & El modelo tiene restricciones fuertes: especialización disjunta, una bahía con una sola orden
activa, una cita confirmada por bahía y horario, historial inmutable, atributos derivados. \\
Alternativas evaluadas & 1. MySQL 8. 2. SQLite. 3. PostgreSQL 16. \\
Decisión & PostgreSQL en el ambiente de pruebas y producción: tipos ENUM, índices únicos parciales,
columnas generadas y disparadores PL/pgSQL. SQLite solo para desarrollo local y pruebas automatizadas;
por eso las reglas también se validan en la capa de servicios. \\
Consecuencias & \textbf{+} La base de datos protege las reglas aunque alguien la modifique directamente.
\textbf{--} Dos motores (SQLite y PostgreSQL): la CI carga el esquema real en PostgreSQL para detectar
diferencias. \\
RNF relacionados & RNF-04, RNF-07, RNF-09 \\
\midrule
ID y título & \textbf{ADR-004 Hospedaje del ambiente de pruebas} \\
Fecha y estado & 29/10/2026 · Propuesto (decidir antes del Sprint 4) \\
Contexto & Se necesita un ambiente de pruebas con HTTPS, PostgreSQL y almacenamiento de fotografías, sin
costo (RNF-07). \\
Alternativas evaluadas & 1. Render (servicio web gratuito) + PostgreSQL gestionado gratuito (Neon o
Supabase). 2. Fly.io. 3. Servidor dentro del taller. \\
Decisión & \pendiente{}. Recomendación inicial: opción 1 por su HTTPS
automático y despliegue desde GitHub. \\
Consecuencias & Opción 1: \textbf{+} sin servidor que mantener; \textbf{--} el plan gratuito suspende el
servicio tras inactividad y limita el almacenamiento (R-06). \\
RNF relacionados & RNF-05, RNF-07 \\
\midrule
ID y título & \textbf{ADR-005 Almacenamiento de fotografías y acceso del cliente con enlace seguro} \\
Fecha y estado & 29/10/2026 · Aceptado (acceso del cliente); fotografías: aceptado para v0.1.0 \\
Contexto & El cliente debe consultar y autorizar sin cuenta ni contraseña, pero nadie debe poder ver la
orden de otro (OWASP A01). Cada recepción genera al menos 4 fotografías que no deben saturar el
hospedaje gratuito. \\
Alternativas evaluadas & \textbf{Acceso:} 1. Cuenta de cliente con contraseña. 2. Folio + teléfono.
3. Token aleatorio en el enlace + código 2FA para autorizar. \textbf{Fotos:} 1. Guardarlas en la base de
datos (BLOB). 2. Sistema de archivos del servidor. 3. Servicio de almacenamiento de objetos. \\
Decisión & \textbf{Acceso:} opción 3 --- token de 256 bits generado con \texttt{secrets}, guardado solo como
hash SHA-256, un enlace activo por orden, que expira al entregar y saldar (RN-01); autorizar exige además
el código de 6 dígitos enviado al correo (RN-02). \textbf{Fotos:} opción 2 en v0.1.0 (carpeta
\texttt{MEDIA\_DIR}, solo JPEG/PNG/WebP de hasta 5 MB, nombre aleatorio); la base guarda solo la ruta, así
que migrar a la opción 3 no cambia el modelo. \\
Consecuencias & \textbf{+} El folio no sirve para entrar (CPA-30); quien intercepte el enlace aún necesita
el correo para autorizar. \textbf{--} Si el cliente pierde el enlace, el taller debe reenviarlo. \textbf{--}
Los archivos locales se pierden si el hospedaje gratuito reinicia el disco: se confirma con
ADR-004 y, de ser necesario, se pasa a almacenamiento de objetos. \\
RNF relacionados & RNF-04, RNF-05, RNF-01 \\
\midrule
ID y título & \textbf{ADR-006 Reglas de negocio como tablas de decisión DMN con Camunda} \\
Fecha y estado & 02/11/2026 · Propuesto (diseño teórico; no implementado) \\
Contexto & Las reglas RN-xx viven hoy como código Python en \texttt{app/servicios/}. Varias son
\emph{decisiones} con umbrales que el cliente puede querer cambiar (10\,\% de re-cotización, 5 días de
vigencia, tipos de vehículo admitidos, crédito a flotillas, conversión del nivel de combustible de RN-14).
Cada cambio de umbral exige hoy modificar código, probar y desplegar. \\
Alternativas evaluadas & 1. Mantener las reglas en Python (situación actual). 2. Tablas de decisión
DMN 1.3 modeladas en Camunda Modeler y evaluadas por el motor de decisiones de Camunda (DMN Engine) a
través de su API REST. 3. Tablas de decisión DMN evaluadas dentro del mismo proceso Python con una
biblioteca DMN. \\
Decisión & Opción 2 como diseño objetivo (detallado en 7.9): las decisiones se modelan en DMN con Camunda
Modeler, se versionan en \texttt{docs/07-arquitectura/dmn/} y la capa de servicios las consulta mediante
un adaptador \texttt{servicios/decisiones.py}. Los procesos (BPMN) siguen en el código; Camunda solo se
usa para \emph{decisiones}. \\
Consecuencias & \textbf{+} Las reglas quedan como tablas legibles que el Gerente puede revisar y validar
sin leer código. \textbf{+} Cada evaluación devuelve la regla que se cumplió, lo que sirve como evidencia
de auditoría (RNF-04). \textbf{--} Agrega un servicio más que desplegar y mantener. \textbf{--} El
motor no está disponible sin conexión: el adaptador debe tener un respaldo local (7.9). \textbf{--}
Hay que confirmar la licencia de la edición de Camunda que se use contra RNF-07. \\
RNF relacionados & RNF-04, RNF-07, RNF-08 \\
\end{longtable}
}

\subsection{7.4 Diagramas de artefactos y de
despliegue}\label{diagramas-de-artefactos-y-de-despliegue}

\figurapendiente{7.4}{Diagrama de artefactos}{el artefacto \texttt{backend/} (paquete Python con
\texttt{app/} y \texttt{requirements.txt}), \texttt{db/schema.sql} y, cuando exista, el paquete estático de
la PWA; con sus dependencias de manifestación hacia los componentes de 7.3.}{docs/07-arquitectura/despliegue/artefactos.puml}

\figurapendiente{7.5}{Diagrama de despliegue del ambiente de pruebas}{nodo «dispositivo» tableta/laptop
del taller y celular del cliente (navegador), nodo «servidor» del hospedaje con el proceso uvicorn y la
API, nodo «base de datos» PostgreSQL 16 y almacenamiento de fotografías; conexiones HTTPS y TCP/5432, según el
hospedaje que se elija en ADR-004.}{docs/07-arquitectura/despliegue/despliegue.puml}

\subsection{7.5 Diagrama de paquetes}\label{diagrama-de-paquetes}

Dependencias entre paquetes de \texttt{backend/app} (todas en un solo sentido, sin ciclos):
\texttt{rutas} $\rightarrow$ \texttt{servicios} $\rightarrow$ \texttt{models} $\rightarrow$ \texttt{db};
\texttt{rutas} $\rightarrow$ \texttt{seguridad}, \texttt{schemas}; dentro de \texttt{servicios}:
\texttt{cobro} $\rightarrow$ \texttt{cotizaciones} $\rightarrow$ \texttt{ordenes}, y \texttt{citas} es
independiente.

\figurapendiente{7.6}{Diagrama de paquetes del backend}{los paquetes \texttt{app.rutas},
\texttt{app.servicios} (ordenes, citas, cotizaciones, cobro), \texttt{app.models}, \texttt{app.schemas},
\texttt{app.seguridad}, \texttt{app.notificaciones} y \texttt{app.db}, con las dependencias «import» descritas
arriba.}{docs/07-arquitectura/paquetes.puml}

\subsection{7.6 Diagramas de clases de la
solución}\label{diagramas-de-clases-de-la-soluciuxf3n}

\figurapendiente{7.7}{Clases de la solución del flujo núcleo}{las clases tal como existen en el código:
routers \texttt{personal}, \texttt{taller} y \texttt{portal}; módulos de servicio \texttt{ordenes}
(\texttt{abrir\_orden}, \texttt{asignar}, \texttt{registrar\_diagnostico}, \texttt{cambiar\_estado}),
\texttt{cotizaciones} (\texttt{crear\_cotizacion}, \texttt{solicitar\_codigo}, \texttt{confirmar\_codigo},
\texttt{huella}), \texttt{cobro} (\texttt{registrar\_pago}, \texttt{entregar}); entidades
\texttt{OrdenServicio}, \texttt{RecepcionVehicular}, \texttt{Cotizacion}, \texttt{Autorizacion},
\texttt{TokenAcceso}; y los esquemas Pydantic de entrada y salida.}{docs/07-arquitectura/clases-solucion.puml}

\subsection{7.7 Diagramas de secuencia}\label{diagramas-de-secuencia}

\figurapendiente{7.8}{Secuencia: recepción del vehículo (escenario principal del flujo núcleo)}{Recepcionista
→ \texttt{POST /ordenes} → \texttt{ordenes.abrir\_orden} (valida vehículo, cita y capacidad) → crea
OrdenServicio, HistorialEstado y RecepcionVehicular → respuesta con folio y alerta de bahías; luego
\texttt{POST /ordenes/\{folio\}/fotografias} por cada posición.}{docs/07-arquitectura/secuencias/recepcion.puml}

\figurapendiente{7.9}{Secuencia: autorización remota con código 2FA}{Cliente → \texttt{POST
/portal/\{token\}/codigo} → \texttt{token\_valido} → \texttt{solicitar\_codigo} → notificaciones (correo);
Cliente → \texttt{POST /portal/\{token\}/autorizacion} → \texttt{confirmar\_codigo} (vigencia, código,
huella) → \texttt{cambiar\_estado(Autorizada)}.}{docs/07-arquitectura/secuencias/autorizacion-remota.puml}

\figurapendiente{7.10}{Secuencia: re-cotización mayor al 10\,\%}{Jefe de mecánicos crea cotización
versión 2 → el servicio compara con el total inicial → queda Pendiente y se envía enlace →
\texttt{terminar} es rechazado (RN-04) hasta que el cliente autoriza.}{docs/07-arquitectura/secuencias/recotizacion.puml}

\figurapendiente{7.11}{Secuencia: pago y entrega}{Recepcionista → \texttt{POST pagos} →
\texttt{registrar\_pago} (no excede saldo) → \texttt{POST entrega} → \texttt{entregar} (saldo o crédito) →
Cerrada, Garantía, bahía libre y token expirado.}{docs/07-arquitectura/secuencias/entrega.puml}

\subsection{7.8 Diagrama de
comunicación}\label{diagrama-de-comunicaciuxf3n}

\figurapendiente{7.12}{Comunicación: asignación de bahía}{objetos \texttt{:RutaTaller},
\texttt{:ServicioOrdenes}, \texttt{os:OrdenServicio}, \texttt{b:Bahia}, \texttt{m:Usuario} y
\texttt{:HistorialEstado}, con los mensajes numerados 1 asignar, 1.1 obtener bahía, 1.2 validar
responsable, 1.3 ocupar bahía, 1.4 cambiarEstado(Abierta) si estaba en espera.}{docs/07-arquitectura/comunicacion-asignacion.puml}

\subsection{7.9 Modelado de decisiones con DMN y
Camunda}\label{modelado-de-decisiones-con-dmn-y-camunda}

Esta sección describe el diseño propuesto en ADR-006. Camunda todavía no forma parte del
código ni del despliegue; las reglas siguen implementadas en Python en \texttt{app/servicios/}.

\paragraph{Qué es DMN y para qué se usa aquí.} DMN (\emph{Decision Model and Notation}, estándar del
OMG, versión 1.3) describe decisiones de negocio como \textbf{tablas de decisión}: columnas de entrada,
columnas de salida y una fila por regla, con una \emph{política de acierto} (\emph{hit policy}) que dice qué
hacer si se cumplen varias filas. Las expresiones de cada celda se escriben en FEEL (\emph{Friendly
Enough Expression Language}), por ejemplo \texttt{> 10}, \texttt{[13..37]} o
\texttt{"Flotilla"}. \textbf{Camunda} aporta dos piezas: Camunda Modeler (aplicación de escritorio
gratuita para dibujar las tablas y guardarlas como archivos \texttt{.dmn}) y el motor de decisiones que
recibe los datos de entrada y devuelve el resultado.

En este proyecto Camunda se usa \textbf{solo para decisiones (DMN)}, no para orquestar procesos (BPMN): el
flujo de la orden (estados de 6.5) sigue controlado por \texttt{servicios/ordenes.py}, y únicamente le
pregunta al motor «¿qué dice la regla?» en los puntos donde hoy hay un \texttt{if} con umbrales de negocio.

\paragraph{Tablas de decisión propuestas.}

{\small
\begin{longtable}{@{}L{0.2\linewidth}L{0.14\linewidth}L{0.22\linewidth}L{0.18\linewidth}L{0.17\linewidth}@{}}
\toprule
\textbf{Decisión (clave)} & \textbf{Regla} & \textbf{Entradas} & \textbf{Salidas} & \textbf{Hit policy} \\
\midrule
\endhead
\bottomrule
\endlastfoot
\texttt{requiere\_recotizacion} & RN-04 & total inicial, total nuevo & requiere autorización (sí/no),
motivo & UNIQUE \\
\texttt{vigencia\_cotizacion} & RN-11 & días desde la emisión & vigente (sí/no) & UNIQUE \\
\texttt{vehiculo\_admitido} & RN-07 & tipo de vehículo, peso (t) & admitido (sí/no), motivo & FIRST \\
\texttt{permite\_entrega} & RN-10 & saldo pendiente, tipo de cliente, crédito autorizado por el Gerente &
entregar (sí/no), motivo & FIRST \\
\texttt{nivel\_combustible} & RN-14 & método de medición, porcentaje OBD-II & nivel (Reserva, 1/4, 1/2,
3/4, Lleno o No verificable) & UNIQUE \\
\texttt{nivel\_contingencia} & RN-15 & hay luz, hay internet, horas de corte, batería (\%) & nivel de
operación (1, 2 o 3) & FIRST \\
\end{longtable}
}

Como ejemplo, la tabla \texttt{permite\_entrega} (RN-10) queda así; se lee de arriba hacia abajo y gana la
primera fila que se cumpla (FIRST):

{\small
\begin{longtable}{@{}L{0.05\linewidth}L{0.17\linewidth}L{0.17\linewidth}L{0.2\linewidth}L{0.1\linewidth}L{0.21\linewidth}@{}}
\toprule
\textbf{\#} & \textbf{saldo} & \textbf{tipoCliente} & \textbf{creditoAutorizado} & \textbf{entregar} &
\textbf{motivo} \\
\midrule
\endhead
\bottomrule
\endlastfoot
1 & \texttt{<= 0} & \texttt{-} & \texttt{-} & \texttt{true} & \texttt{"Saldo cubierto"} \\
2 & \texttt{> 0} & \texttt{"Flotilla"} & \texttt{true} & \texttt{true} & \texttt{"Crédito a flotilla"} \\
3 & \texttt{> 0} & \texttt{-} & \texttt{-} & \texttt{false} & \texttt{"Saldo pendiente (RN-10)"} \\
\end{longtable}
}

\paragraph{Cómo se integra con el backend.}

\begin{enumerate}
\tightlist
\item \textbf{Modelado y versionado.} Cada tabla se dibuja en Camunda Modeler y se guarda en
\texttt{docs/07-arquitectura/dmn/<clave>.dmn}, en git, con revisión por Pull Request como el resto del
código. El Gerente valida la tabla impresa antes de publicarla.
\item \textbf{Despliegue.} Los archivos \texttt{.dmn} se despliegan al motor de Camunda, que corre como un
contenedor aparte junto a la API (un nodo más en el diagrama de despliegue de 7.4).
\item \textbf{Evaluación.} Un nuevo módulo \texttt{app/servicios/decisiones.py} expone una función
\texttt{evaluar(clave, datos)} que llama a la API REST del motor con las variables de entrada y devuelve
las salidas. Los demás servicios solo dependen de esa función; por ejemplo, \texttt{cobro.entregar} llama
a \texttt{evaluar("permite\_entrega", \{...\})} en lugar de comparar el saldo por su cuenta. Así se respeta
la regla de 7.1: la lógica de negocio sigue fuera de las rutas.
\item \textbf{Respaldo sin conexión.} Si el motor no responde (corte de internet, RNF-08, o contingencia
de RN-15), el adaptador evalúa una copia local de las mismas tablas con una biblioteca DMN de Python, o, en
su defecto, la implementación actual en Python. La PWA sin conexión no consulta al motor: las decisiones
que dependen de él se vuelven a evaluar al sincronizar (RNF-09).
\item \textbf{Auditoría.} Cada evaluación guarda en el historial de la orden la clave de la decisión, su
versión y la fila que se cumplió, como evidencia de por qué se permitió o bloqueó una operación
(RNF-04).
\item \textbf{Pruebas.} Cada fila de cada tabla tiene su caso en \texttt{backend/tests/}
(\texttt{test\_dmn\_<clave>.py}); los casos CPA existentes de RN-04, RN-07, RN-10 y RN-11 deben seguir
pasando sin cambios, lo que prueba que la migración no alteró el comportamiento.
\end{enumerate}

\paragraph{Plan de adopción.} Se migraría una decisión a la vez, empezando por \texttt{requiere\_recotizacion}
(RN-04), que es la más simple y la que el cliente tiene más probabilidad de ajustar. Antes de iniciar hay
que confirmar con el equipo y el docente: (1) la edición y versión de Camunda y que su licencia cumpla
RNF-07 para el uso que se le dará; (2) dónde se hospeda el motor (ADR-004); y (3) si el beneficio justifica
un servicio adicional en un taller de 4 bahías o si conviene la opción 3 de ADR-006 (DMN evaluado dentro de
la misma API).

\figurapendiente{7.13}{Diagrama de requisitos de decisión (DRD)}{las seis decisiones de la tabla de 7.9
como rectángulos de decisión, sus datos de entrada (orden, cotización, cliente, vehículo, recepción, estado
de la infraestructura) como óvalos y la fuente de conocimiento «Reglas de negocio 4.7» enlazada a cada
decisión; exportado de Camunda Modeler.}{docs/07-arquitectura/dmn/drd.dmn}

\section{8. Diseño de la interfaz de
usuario}\label{diseuxf1o-de-la-interfaz-de-usuario}

Esta sección corresponde al corte de la Review 5 (Sprints 4 y 5). La API
del flujo núcleo ya existe (sección 9); el prototipo de alta fidelidad,
la evaluación heurística y la prueba con usuarios se hacen sobre los
wireframes de 4.10 una vez validados con el cliente.

\subsection{8.1 Prototipo de alta
fidelidad}\label{prototipo-de-alta-fidelidad}

No aplica todavía: el prototipo navegable se elabora en el Sprint 4 a
partir de los wireframes de la sección 4.10.

\subsection{8.2 Accesibilidad}\label{accesibilidad}

No aplica todavía: los criterios de WCAG 2.2 (contraste, tamaño de
objetivos táctiles, textos alternativos y navegación por teclado) se
verifican sobre el prototipo en el Sprint 5.

\subsection{8.3 Evaluación heurística
recibida}\label{evaluaciuxf3n-heuruxedstica-recibida}

No aplica todavía: la evaluación con las 10 heurísticas de Nielsen la
realiza el otro equipo sobre el prototipo del Sprint 5.

\subsection{8.4 Prueba con usuarios y
SUS}\label{prueba-con-usuarios-y-sus}

No aplica todavía: la prueba con 3 a 5 personas ajenas al equipo y el
cálculo del SUS se realizan en el Sprint 5.

\section{9. Implementación del Release
v0.1.0}\label{implementaciuxf3n-del-release-v0.1.0}

\subsection{9.1 Pila tecnológica}\label{pila-tecnoluxf3gica}

{\small
\begin{longtable}{@{}L{0.27\linewidth}L{0.47\linewidth}L{0.2\linewidth}@{}}
\toprule
\textbf{Capa} & \textbf{Tecnología y versión} & \textbf{ADR que la justifica} \\
\midrule
\endhead
\bottomrule
\endlastfoot
Interfaz & PWA (\pendiente{}). Hoy: documentación interactiva OpenAPI de FastAPI
(\texttt{/docs}) & ADR-001, ADR-002 \\
Lógica de negocio / API & Python 3.12, FastAPI 0.115+, Pydantic 2, SQLAlchemy 2.0, Uvicorn & ADR-002 \\
Base de datos & PostgreSQL 16 (esquema \texttt{db/schema.sql}); SQLite para desarrollo y pruebas & ADR-003 \\
Hospedaje del ambiente de pruebas & \pendiente{} & ADR-004 \\
Integración continua & GitHub Actions (\texttt{.github/workflows/ci.yml}) & --- \\
\end{longtable}
}

\subsection{9.2 Instrucciones de instalación y
ejecución}\label{instrucciones-de-instalaciuxf3n-y-ejecuciuxf3n}

Los mismos pasos están en el \texttt{README.md} del repositorio. Requisitos: Python 3.12 o superior y Git.

\begin{enumerate}
\tightlist
\item Clonar el repositorio y entrar a \texttt{backend/}.
\item Crear un entorno virtual e instalar dependencias:\\
\texttt{python -m venv .venv} y \texttt{pip install -r requirements.txt}.
\item Copiar \texttt{.env.example} como \texttt{.env}. Sin cambios usa SQLite local; para PostgreSQL,
crear la base con \texttt{psql -d taller -f db/schema.sql} y poner su \texttt{DATABASE\_URL}. El
\texttt{.env} nunca se sube al repositorio.
\item Cargar los datos ficticios: \texttt{python -m app.semilla -{}-demo}. Crea las 4 bahías, un usuario por
rol (correos \texttt{@taller.test}), un cliente particular y una flotilla con crédito. La contraseña se
toma de \texttt{SEED\_PASSWORD} o se genera e imprime una sola vez.
\item Ejecutar: \texttt{uvicorn app.main:app -{}-reload} y abrir \texttt{http://localhost:8000/docs}.
\item Pruebas: \texttt{python -m pytest}.
\end{enumerate}

\subsection{9.3 Integración continua y flujo de
trabajo}\label{integraciuxf3n-continua-y-flujo-de-trabajo}

El pipeline de GitHub Actions se ejecuta en cada push a \texttt{main} y en cada Pull Request, de modo que
toda rama de historia o tarea se prueba antes de integrarse. Tiene tres trabajos:

\begin{enumerate}
\tightlist
\item \textbf{Pruebas del backend:} instala dependencias y ejecuta las 35 pruebas de pytest.
\item \textbf{Esquema en PostgreSQL:} levanta PostgreSQL 16 y carga \texttt{db/schema.sql} con
\texttt{ON\_ERROR\_STOP}, para detectar errores del esquema real.
\item \textbf{Documento de Ingeniería:} compila este documento LaTeX y publica el PDF como artefacto.
\end{enumerate}

\textbf{Flujo de trabajo:} GitHub Flow. Cada historia o tarea se trabaja en su propia rama
\texttt{tipo/número-descripción} (por ejemplo, \texttt{fix/9-version-stakeholders}), creada desde
\texttt{main}; los commits llevan el formato «Rol: mensaje» y la rama se integra con un Pull Request que
cierra su issue.
Por acuerdo del equipo, la rama main se integra mediante Pull Request con al menos una revisión y la CI aprobada
(CONTRIBUTING.md).

\figurapendiente{9.1}{Ejecución exitosa del pipeline de integración continua}{captura de la pestaña
Actions de GitHub con los tres trabajos (Pruebas del backend, Esquema en PostgreSQL, Documento de
Ingeniería) en verde para el último commit de main.}{.github/workflows/ci.yml}

\subsection{9.4 Evidencia del flujo
núcleo}\label{evidencia-del-flujo-nuxfacleo}

\pendiente{}.
Mientras tanto, el flujo completo está cubierto por las pruebas automatizadas CPA-01, CPA-13, CPA-22 y
CPA-35.

\figurapendiente{9.2}{Apertura de orden con checklist y fotografías}{pantalla de recepción en el
ambiente de pruebas con datos ficticios: kilometraje, combustible, checklist completo, las 4 fotografías
y el folio generado.}{docs/09-implementacion/evidencia/01-recepcion.png}

\figurapendiente{9.3}{Diagnóstico registrado}{pantalla de la tableta del jefe de mecánicos con la
bahía y el responsable asignados, el diagnóstico guardado y la orden en estado
Diagnosticada.}{docs/09-implementacion/evidencia/02-diagnostico.png}

\figurapendiente{9.4}{Cotización con autorización registrada}{portal del cliente en un celular: la
cotización desglosada, la selección de opcionales, la captura del código de 6 dígitos y la confirmación
con fecha y hora.}{docs/09-implementacion/evidencia/03-autorizacion.png}

\figurapendiente{9.5}{Lista de órdenes abiertas}{pantalla con las órdenes abiertas, su estado, bahía,
responsable y días de estancia.}{docs/09-implementacion/evidencia/04-ordenes-abiertas.png}

\section{10. Reporte final de pruebas e
incidencias}\label{reporte-final-de-pruebas-e-incidencias}

\subsection{10.1 Resultado de los casos de
prueba}\label{resultado-de-los-casos-de-prueba}

Corte parcial (29/10/2026). El reporte final se completa en el Sprint 6 con la
ejecución de aceptación con el cliente y el QA cruzado.

{\def\LTcaptype{none} % do not increment counter
\begin{longtable}[]{@{}
  >{\raggedright\arraybackslash}p{(\linewidth - 2\tabcolsep) * \real{0.5984}}
  >{\raggedright\arraybackslash}p{(\linewidth - 2\tabcolsep) * \real{0.4016}}@{}}
\toprule\noalign{}
\begin{minipage}[b]{\linewidth}\raggedright
\textbf{Indicador}
\end{minipage} & \begin{minipage}[b]{\linewidth}\raggedright
\textbf{Valor}
\end{minipage} \\
\midrule\noalign{}
\endhead
\bottomrule\noalign{}
\endlastfoot
Casos de prueba diseñados & 35 (CPA-01 a CPA-35, sección 5.2) \\
Casos ejecutados & 35, automatizados \\
Pasan & 35 \\
Fallan & 0. CPA-11 falló en su primera ejecución (el portal respondía 404 ante una cotización
vencida, en lugar del aviso de RN-11); se corrigió antes de integrar a main. \\
Bloqueados & 0 \\
Pruebas unitarias automatizadas (número y resultado en CI) & 35 pruebas de pytest; se ejecutan en
GitHub Actions en cada push y Pull Request \\
\end{longtable}
}

\subsection{10.2 Incidencias}\label{incidencias}

Incluye las incidencias internas, del QA cruzado y del cliente. El tiempo
de atención se mide desde que se registró hasta que se cerró.

{\def\LTcaptype{none} % do not increment counter
\begin{longtable}[]{@{}
  >{\raggedright\arraybackslash}p{(\linewidth - 10\tabcolsep) * \real{0.1497}}
  >{\raggedright\arraybackslash}p{(\linewidth - 10\tabcolsep) * \real{0.1495}}
  >{\raggedright\arraybackslash}p{(\linewidth - 10\tabcolsep) * \real{0.1496}}
  >{\raggedright\arraybackslash}p{(\linewidth - 10\tabcolsep) * \real{0.2094}}
  >{\raggedright\arraybackslash}p{(\linewidth - 10\tabcolsep) * \real{0.1497}}
  >{\raggedright\arraybackslash}p{(\linewidth - 10\tabcolsep) * \real{0.1922}}@{}}
\toprule\noalign{}
\begin{minipage}[b]{\linewidth}\raggedright
\textbf{Severidad}
\end{minipage} & \begin{minipage}[b]{\linewidth}\raggedright
\textbf{Registradas}
\end{minipage} & \begin{minipage}[b]{\linewidth}\raggedright
\textbf{Cerradas}
\end{minipage} & \begin{minipage}[b]{\linewidth}\raggedright
\textbf{Rechazadas (con justificación)}
\end{minipage} & \begin{minipage}[b]{\linewidth}\raggedright
\textbf{Abiertas}
\end{minipage} & \begin{minipage}[b]{\linewidth}\raggedright
\textbf{Tiempo promedio de atención}
\end{minipage} \\
\midrule\noalign{}
\endhead
\bottomrule\noalign{}
\endlastfoot
S1 & & & & & \\
S2 & & & & & \\
S3 & & & & & \\
S4 & & & & & \\
\end{longtable}
}

\subsection{10.3 Defectos conocidos que pasan a la Fase
2}\label{defectos-conocidos-que-pasan-a-la-fase-2}

{\def\LTcaptype{none} % do not increment counter
\begin{longtable}[]{@{}
  >{\raggedright\arraybackslash}p{(\linewidth - 8\tabcolsep) * \real{0.1174}}
  >{\raggedright\arraybackslash}p{(\linewidth - 8\tabcolsep) * \real{0.3483}}
  >{\raggedright\arraybackslash}p{(\linewidth - 8\tabcolsep) * \real{0.1282}}
  >{\raggedright\arraybackslash}p{(\linewidth - 8\tabcolsep) * \real{0.1282}}
  >{\raggedright\arraybackslash}p{(\linewidth - 8\tabcolsep) * \real{0.2779}}@{}}
\toprule\noalign{}
\begin{minipage}[b]{\linewidth}\raggedright
\textbf{Issue}
\end{minipage} & \begin{minipage}[b]{\linewidth}\raggedright
\textbf{Descripción}
\end{minipage} & \begin{minipage}[b]{\linewidth}\raggedright
\textbf{Severidad}
\end{minipage} & \begin{minipage}[b]{\linewidth}\raggedright
\textbf{Prioridad}
\end{minipage} & \begin{minipage}[b]{\linewidth}\raggedright
\textbf{Motivo por el que no se corrigió}
\end{minipage} \\
\midrule\noalign{}
\endhead
\bottomrule\noalign{}
\endlastfoot
& & & & \\
& & & & \\
\end{longtable}
}

\subsection{10.4 QA cruzado realizado al otro
equipo}\label{qa-cruzado-realizado-al-otro-equipo}

No aplica todavía: el QA cruzado al otro equipo lo realiza el QA Lead en
el Sprint 6 (casos ejecutados, defectos reportados y correcciones
verificadas en la ronda 2).

\section{11. Paquete de transferencia a la Fase
2}\label{paquete-de-transferencia-a-la-fase-2}

Lo necesario para que un equipo, incluso uno distinto, continúe el
proyecto en Ingeniería de Software II. Se completa en el Sprint 6.

{\def\LTcaptype{none} % do not increment counter
\begin{longtable}[]{@{}
  >{\raggedright\arraybackslash}p{(\linewidth - 2\tabcolsep) * \real{0.3204}}
  >{\raggedright\arraybackslash}p{(\linewidth - 2\tabcolsep) * \real{0.6796}}@{}}
\toprule\noalign{}
\begin{minipage}[b]{\linewidth}\raggedright
\textbf{Apartado}
\end{minipage} & \begin{minipage}[b]{\linewidth}\raggedright
\textbf{Contenido}
\end{minipage} \\
\midrule\noalign{}
\endhead
\bottomrule\noalign{}
\endlastfoot
Backlog restante priorizado & Historias con release Fase 2 de la
sección 4.5: HU-05, HU-11, HU-19, HU-24 y HU-29. \\
Deuda técnica conocida & Se registra al cierre del Sprint 6. \\
Riesgos abiertos & Registro de riesgos de la sección 2.4. \\
Decisiones pendientes con el cliente & Hospedaje del ambiente de
pruebas (ADR-004) y validación de RN-12 y RN-13 (sección 4.11). \\
Accesos y configuración & Repositorio en GitHub y GitHub Actions; las
credenciales no se incluyen en el documento ni en el repositorio. \\
Recomendaciones & Se registran al cierre del Sprint 6. \\
\end{longtable}
}

\section{12. Retrospectiva del proyecto y contribución
individual}\label{retrospectiva-del-proyecto-y-contribuciuxf3n-individual}

\subsection{12.1 Retrospectiva del
proyecto}\label{retrospectiva-del-proyecto}

{\def\LTcaptype{none} % do not increment counter
\begin{longtable}[]{@{}
  >{\raggedright\arraybackslash}p{(\linewidth - 4\tabcolsep) * \real{0.3333}}
  >{\raggedright\arraybackslash}p{(\linewidth - 4\tabcolsep) * \real{0.3333}}
  >{\raggedright\arraybackslash}p{(\linewidth - 4\tabcolsep) * \real{0.3333}}@{}}
\toprule\noalign{}
\begin{minipage}[b]{\linewidth}\raggedright
\textbf{¿Qué funcionó bien?}
\end{minipage} & \begin{minipage}[b]{\linewidth}\raggedright
\textbf{¿Qué no funcionó?}
\end{minipage} & \begin{minipage}[b]{\linewidth}\raggedright
\textbf{¿Qué haríamos distinto en la Fase 2?}
\end{minipage} \\
\midrule\noalign{}
\endhead
\bottomrule\noalign{}
\endlastfoot
\pendiente{} & \pendiente{} & \pendiente{} \\
\end{longtable}
}

\subsection{12.2 Contribución
individual}\label{contribuciuxf3n-individual}


{\def\LTcaptype{none} % do not increment counter
\begin{longtable}[]{@{}
  >{\raggedright\arraybackslash}p{(\linewidth - 6\tabcolsep) * \real{0.1923}}
  >{\raggedright\arraybackslash}p{(\linewidth - 6\tabcolsep) * \real{0.3205}}
  >{\raggedright\arraybackslash}p{(\linewidth - 6\tabcolsep) * \real{0.2735}}
  >{\raggedright\arraybackslash}p{(\linewidth - 6\tabcolsep) * \real{0.2137}}@{}}
\toprule\noalign{}
\begin{minipage}[b]{\linewidth}\raggedright
\textbf{Integrante}
\end{minipage} & \begin{minipage}[b]{\linewidth}\raggedright
\textbf{Principales aportes}
\end{minipage} & \begin{minipage}[b]{\linewidth}\raggedright
\textbf{Evidencia (issues, PR, documentos)}
\end{minipage} & \begin{minipage}[b]{\linewidth}\raggedright
\textbf{Revisiones de código o documentos hechas a otros}
\end{minipage} \\
\midrule\noalign{}
\endhead
\bottomrule\noalign{}
\endlastfoot
Ángel Yael Avendaño Guzmán (Product Owner) & Enlace con el cliente,
levantamiento de requerimientos en entrevistas, priorización del Product
Backlog y aceptación de entregables. & Rama \texttt{rol/product-owner} &
--- \\
Abdías Daniel Chávez Cruz (Scrum Master / Tech Lead) & Facilitación de
Scrum, arquitectura del backend en FastAPI, diseño de la base de datos
relacional y gestión técnica de Pull Requests. & Rama
\texttt{rol/scrum-master}; \texttt{backend/}; \texttt{db/schema.sql} &
--- \\
Dulce María Estudillo Rosado (QA Lead / Developer) & Plan de pruebas,
casos de aceptación, validación de calidad y QA cruzado al otro
equipo. & Rama \texttt{rol/qa-lead}; sección 5 & --- \\
Isis Alhelí González Maldonado (Developer) & Especificación de requisitos
(SRS), historias de usuario y modelado de procesos del negocio. & Rama
\texttt{rol/developers}; sección 4 & --- \\
Isis Yamilet Meléndez Hernández (Developer) & Integración y compilación
del Documento de Ingeniería, flujos de CI en GitHub Actions y apoyo en
desarrollo. & Rama \texttt{rol/developers}; \texttt{docs/documento/};
\texttt{.github/workflows/} & --- \\
\end{longtable}
}

\section{13. Referencias}\label{referencias}

Schwaber, K., \& Sutherland, J. (2020). \emph{La Guía de Scrum}.
\href{https://scrumguides.org/}{\ul{https://scrumguides.org/}}

ISO/IEC/IEEE. (2018). \emph{ISO/IEC/IEEE 29148:2018 Systems and software
engineering --- Life cycle processes --- Requirements engineering}. ISO.

ISO/IEC. (2023). \emph{ISO/IEC 25010:2023 Systems and software
engineering --- SQuaRE --- Product quality model}. ISO.

ISO/IEC/IEEE. (2021). \emph{ISO/IEC/IEEE 29119-3:2021 Software and
systems engineering --- Software testing --- Part 3: Test
documentation}. ISO.

Object Management Group. (2017). \emph{OMG Unified Modeling Language,
Version 2.5.1}.
\href{https://www.omg.org/spec/UML/2.5.1}{\ul{https://www.omg.org/spec/UML/2.5.1}}

Object Management Group. (2021). \emph{Decision Model and Notation (DMN),
Version 1.3}.
\href{https://www.omg.org/spec/DMN/1.3}{\ul{https://www.omg.org/spec/DMN/1.3}}

Camunda. (s.f.). \emph{DMN reference} [Documentación].
\href{https://docs.camunda.io/docs/components/modeler/dmn/}{\ul{https://docs.camunda.io/docs/components/modeler/dmn/}}
\section{Anexos}\label{anexos}

Los anexos se numeran con letras. Los anexos B a E (minutas de las Sprint
Reviews, detalle de historias y casos de prueba, manual de uso y aviso de
privacidad) se agregan en el corte que los requiere.

\subsection{Anexo A. Guion y minuta de la entrevista con el cliente}\label{anexo-a}

\subsubsection{A.1 Guion de la entrevista inicial}\label{anexo-a1}

Objetivo: entender cómo opera hoy el taller, qué problemas le cuestan
dinero o clientes y qué reglas sigue aunque no estén escritas. Las
preguntas se agrupan por etapa del servicio; en cada una se pide un
ejemplo reciente y se pregunta por las excepciones (quién puede romper la
regla y si queda registro).

\paragraph{Negocio y operación general}
\begin{enumerate}
\tightlist
\item ¿Qué tipos de vehículos atienden y cuáles no reciben?
\item ¿Cuántas bahías de trabajo tienen y cuántos vehículos caben en
espera? ¿Qué pasa cuando todo está ocupado?
\item ¿Quiénes trabajan en el taller y qué hace cada uno durante el día?
\item ¿Qué problemas les cuestan hoy más dinero o más clientes?
\end{enumerate}

\paragraph{Citas}
\begin{enumerate}
\tightlist
\setcounter{enumi}{4}
\item ¿Cómo agenda hoy un cliente su cita y quién la registra?
\item Si un cliente pidiera la cita en línea, ¿debe apartar lugar de
inmediato o alguien la confirma?
\end{enumerate}

\paragraph{Recepción del vehículo}
\begin{enumerate}
\tightlist
\setcounter{enumi}{6}
\item ¿Qué datos del vehículo y del cliente anotan al recibirlo?
\item ¿Han tenido reclamos por golpes o faltantes? ¿Cómo los resolvieron
y qué evidencia tenían?
\item ¿Qué hacen cuando no funciona el indicador de combustible?
\end{enumerate}

\paragraph{Diagnóstico, bahías y mecánicos}
\begin{enumerate}
\tightlist
\setcounter{enumi}{9}
\item ¿Quién decide a qué bahía y a qué mecánico va cada vehículo?
\item ¿Pueden trabajar varios mecánicos en la misma orden? ¿Quién
responde por el trabajo?
\item Si una refacción tarda varios días en llegar, ¿qué pasa con el
vehículo y con la bahía?
\end{enumerate}

\paragraph{Cotización y autorización}
\begin{enumerate}
\tightlist
\setcounter{enumi}{12}
\item ¿Cómo autoriza hoy el cliente una cotización y qué queda como
prueba?
\item ¿Han tenido disputas por cobros que el cliente dice no haber
autorizado? ¿Cómo se resolvieron?
\item ¿El cliente puede aceptar unos conceptos y rechazar otros?
\item ¿Qué hacen si durante la reparación aparece un trabajo que sube el
costo?
\item ¿Cuánto tiempo es válida una cotización?
\end{enumerate}

\paragraph{Cobro, entrega y seguimiento}
\begin{enumerate}
\tightlist
\setcounter{enumi}{17}
\item ¿Se entrega alguna vez un vehículo sin pagar completo? ¿Quién lo
autoriza y para qué clientes?
\item ¿Cuántas llamadas reciben para preguntar por el avance y quién las
atiende?
\item ¿Durante cuánto tiempo debería el cliente poder consultar su
vehículo en línea?
\item ¿Qué garantía dan sobre los trabajos y cómo la consultan hoy?
\end{enumerate}

\paragraph{Infraestructura}
\begin{enumerate}
\tightlist
\setcounter{enumi}{21}
\item ¿Qué equipos tienen (computadoras, tabletas, internet) y qué tan
seguido se va la luz o el internet?
\end{enumerate}

\subsubsection{A.2 Minuta de la entrevista del 19/10/2026}\label{anexo-a2}

{\small
\begin{longtable}[]{@{}
  >{\raggedright\arraybackslash}p{0.25\linewidth}
  >{\raggedright\arraybackslash}p{0.7\linewidth}@{}}
\toprule
\textbf{Dato} & \textbf{Contenido} \\
\midrule
\endhead
\bottomrule
\endlastfoot
Fecha & 19/10/2026 \\
Hora y modalidad & \pendiente{} \\
Asistentes por el cliente & \pendiente{} \\
Asistentes por el equipo & \pendiente{} \\
Elaboró & Product Owner (Ángel Yael Avendaño Guzmán) \\
Objetivo & Levantar los requisitos iniciales del sistema a partir del
anteproyecto y del guion A.1. \\
\end{longtable}
}

\textbf{Problemas expresados por el cliente.} Disputas por cobros que el
cliente dice no haber autorizado («dijo que sí, pero cobra otra cosa»),
porque las autorizaciones son verbales o por mensaje; reclamos por daños
sin checklist ni fotografías del estado de llegada; muchas llamadas para
preguntar por el avance, que interrumpen a la recepcionista; bahías y
carga de trabajo asignadas de palabra, con vehículos detenidos sin
justificación; e historial de reparaciones y garantías archivado en
carpetas (sección 3.1).

\textbf{Hallazgos y acuerdos.} Cada hallazgo tiene un identificador que
se cita como fuente en las reglas de negocio de la sección 4.7.

{\small
\begin{longtable}[]{@{}
  >{\raggedright\arraybackslash}p{0.07\linewidth}
  >{\raggedright\arraybackslash}p{0.62\linewidth}
  >{\raggedright\arraybackslash}p{0.25\linewidth}@{}}
\toprule
\textbf{ID} & \textbf{Hallazgo o acuerdo} & \textbf{Se refleja en} \\
\midrule
\endhead
\bottomrule
\endlastfoot
E-01 & Toda autorización de cotización, remota o presencial, debe tener
una doble verificación con un código de 6 dígitos. & RN-02 \\
E-02 & El cliente puede aprobar o rechazar por separado los conceptos
opcionales de una cotización; solo se ejecuta lo autorizado. & RN-03 \\
E-03 & El enlace de seguimiento del cliente funciona mientras el vehículo
está en el taller y vence al entregarlo y saldarlo. & RN-01 \\
E-04 & El taller tiene 4 bahías (un vehículo por bahía) y un
estacionamiento de espera para 5 vehículos. & RN-05 \\
E-05 & Si una refacción tarda varios días, el vehículo pasa al
estacionamiento para liberar la bahía. & RN-06 \\
E-06 & Solo se atienden vehículos ligeros: automóviles, motocicletas y
camionetas de hasta 3 toneladas. & RN-07 \\
E-07 & Las citas pedidas en línea no apartan lugar hasta que la
recepcionista las confirma. & RN-08 \\
E-08 & Varios mecánicos pueden colaborar en una orden, pero siempre hay
un responsable. & RN-09 \\
E-09 & No se entrega un vehículo con saldo pendiente; solo el gerente
puede autorizar la entrega a crédito a flotillas empresariales
registradas. & RN-10 \\
E-10 & Hay cortes de luz e internet; el personal necesita seguir
capturando sin conexión. & RNF-08, RNF-09 \\
E-11 & El jefe de mecánicos trabajará con 1 o 2 tabletas en un ambiente
con ruido, polvo y manos sucias. & 4.3.1, RNF-02, RNF-03 \\
\end{longtable}
}

\textbf{Pendientes para la siguiente sesión con el cliente.} Confirmar el
mínimo de 4 fotografías al recibir el vehículo (RN-12) y la fuente de la
garantía de 30 días sobre la mano de obra (RN-13). Responsable: Product
Owner.

\subsection{Anexo F. Atención de observaciones de la Review 1}\label{anexo-f}

Observaciones recibidas el 10/11/2026 sobre la versión v0.1. Cada
observación tiene su issue con las etiquetas \texttt{tipo:defecto-doc},
severidad y \texttt{origen:cliente}, y se corrige en su propia rama con
un Pull Request que cierra el issue.

Cada corrección aumenta en uno el tercer dígito de la versión del
documento: la v0.2 es la entrega de la Review 1 sin correcciones, la
v0.2.1 incluye la primera corrección, la v0.2.2 la segunda, y así
sucesivamente. La versión se registra en la portada, en el control de
versiones, en esta tabla y con una etiqueta de Git del mismo nombre.

{\small
\begin{longtable}[]{@{}
  >{\raggedright\arraybackslash}p{0.06\linewidth}
  >{\raggedright\arraybackslash}p{0.09\linewidth}
  >{\raggedright\arraybackslash}p{0.07\linewidth}
  >{\raggedright\arraybackslash}p{0.06\linewidth}
  >{\raggedright\arraybackslash}p{0.61\linewidth}@{}}
\toprule
\textbf{ID} & \textbf{Versión} & \textbf{Issue} & \textbf{PR} & \textbf{Cómo se resolvió} \\
\midrule
\endhead
\bottomrule
\endlastfoot
B11 & v0.2.1 & \#9 & \#15 & Se retiraron de 3.4 y de la Figura 3.1 al equipo de
desarrollo y al docente evaluador, que no son stakeholders del producto.
El documento se versiona según la tabla 0.2 del machote: la v0.1
(Review 0) contiene solo portada, 1, 2.1--2.4 y 3, y la v0.2 (Review 1)
agrega 4.1--4.7; cada entrega incluye únicamente las secciones de su
corte. \\
B5 & v0.2.2 & \#10 & \#16 & Se cambió la estrategia de ramas a GitHub Flow: una
rama por historia o tarea con la convención
\texttt{tipo/número-descripción}. Se actualizaron CONTRIBUTING.md, el
README y la sección 2.1; las ramas por integrante
(\texttt{rama/<NOMBRE>}) ya no reciben trabajo nuevo. Las correcciones
de esta revisión ya usan la convención
(\texttt{fix/9-version-stakeholders}, \texttt{docs/10-ramas-por-historia},
\texttt{docs/20-anexo-entrevista}). \\
B3 & v0.2.3 & \#20 & \#21 & Se agregó el Anexo A con el guion de la entrevista
(22 preguntas) y la minuta del 19/10/2026, con los hallazgos E-01 a E-11.
Cada regla de 4.7 que proviene de la entrevista cita en su fuente el
hallazgo de la minuta. Los anexos ya no contienen texto del machote. La
sección 4.11 no forma parte de este corte y se actualiza en la Review 2. \\
B4 & v0.2.4 & \#28 & \#61 & Las 30 historias se registraron como issues
(\#31 a \#60) y 4.5 muestra el número de cada una. En 2.3 se
registraron los puntos del Sprint 1 y la ruta a la minuta de la Review 1
(\ruta{docs/minutas/2026-11-10-review-1.md}). En 12.2 se anotaron los
commits, los Pull Requests y las revisiones de cada integrante, con datos
de GitHub. Siguen pendientes las fechas de los sprints, la minuta de la
Review 0 y las retrospectivas. \\
B2 & v0.2.5 & \#29 & \#62 & Primera parte: se elaboró la propuesta de
alcance (\ruta{docs/04-requisitos/propuesta-release-v0.1.0.md}), que
deja en el Release v0.1.0 las 11 historias del flujo núcleo (40 puntos)
y pasa 14 a la Fase 2. Falta negociarla con el cliente y actualizar 3.3,
3.6, 4.5 y el story map; el issue sigue abierto. \\
B6 & v0.2.6 & \#30 & \#63 & Primera parte: se agregó la sección 2.6 con
la declaración de uso de inteligencia artificial. Falta que cada autor
declare el contenido original de las secciones 1 a 4 y realizar la
sesión interna con minuta; el issue sigue abierto. \\
\end{longtable}
}


\end{document}
